% 本文件由 LaTeX 排版 Agent 根据有效 translation.json 自动生成。
% author=Ambrose; work_id=3402; title=论圣神

\chapter{\OriginalTerm{论圣神}{On the Holy Spirit}}

% structure: p0001 source=h1 target=omitted reason=page-title-already-rendered-by-parent

卷一\newline{}卷二\newline{}卷三\par

\SourceRecord{Translated by H. de Romestin, E. de Romestin and H.T.F. Duckworth.；Nicene and Post-Nicene Fathers, Second Series；Vol. 10.；Edited by Philip Schaff and Henry Wace.；Buffalo, NY: Christian Literature Publishing Co.,；1896.；<http://www.newadvent.org/fathers/3402.htm>.}

% structure: p0001 source=h1 target=section reason=page-title

\section{\WorkTitle{论圣神（卷一）}}

% structure: p0002 source=h2 target=subsection reason=relative-heading-rank; base=h2

\subsection{引言}

关于\OriginalTerm{圣神}{Holy Spirit}的三卷著作，正如\Person{圣盎博罗削}自己所言，是继论述信德诸书之后的续作；这两部论著有时被合引为一，冠以\WorkTitle{论圣三}之名。然而，我们从\Person{格拉齐安}皇帝致\Person{圣盎博罗削}的书信以及他的回函可知，每部论著各自独立，而\WorkTitle{论圣神}则写于数年之后，即公元381年。\par

在卷一，\Person{圣盎博罗削}首先以寓意解经法阐释\Person{基德红}与羊毛的历史，在羊毛变干而禾场变湿的奇迹中，他看\OriginalTerm{圣神}{Holy Spirit}离开\Place{犹太人}而倾注于\Place{外邦人}的预象。接着切入更直接的主题，证明圣神超越一切受造界，且是真天主，特别提出一项论证：干犯圣神的罪，无论今世或来世，永不得赦免。他指明圣经如何称圣神为天主之神；圣神藉着先知和宗徒发言；圣神圣化世人，且由圣经所述之神秘圣油所预表。其后，\Person{圣盎博罗削}论述圣神与\OriginalTerm{天主圣三}{Trinity}另两位的同一性，并表明圣神的被派遣丝毫无损于此同一性，反而在天主三位之间存在着和平、爱德及其它德性的圆满合一。\par

卷二开篇，以处理卷一\Person{基德红}历史同样的方式，来处理\Person{三松}的历史。只要圣神与\Person{三松}同在，他便事事成功；一旦圣神离开，他便陷入不幸。书中揭示，圣神的权能与圣父及圣子的权能相同，且在计划与运作、赋予人生命方面，三位和谐一致。圣神是造物主，故此应受钦崇；祂与圣父及圣子共同建立教会；卷末则证明三位在运作上的合一。\par

卷三继续同一论题，指出先知与宗徒的使命，甚至圣子本身的使命，均应归诸圣神；然而圣子丝毫不受制于圣神，因为圣神同样由圣父及圣子接受派遣。接着论及圣神的天主性，并加以证明，此时亦趁便指明并非有三位天主或三位主，因为天主三位在圣德和性体上是合一的。全书则以撮述若干主要论证作结。\par

鲜可置疑的是，这正是遭到\Person{圣热罗尼莫}猛烈抨击的那部著作，而\Person{圣盎博罗削}便是那位作者。抨击全文可查阅\Person{鲁菲努斯}的\WorkTitle{护教书}（见本系列第三卷，第470页）。\Person{圣盎博罗削}被比作一只披上别鸟羽毛的寒鸦，且被指控\Quote{以拉丁文写出坏东西，却是取自希腊文的好东西}。\Person{圣热罗尼莫}甚至不厌其烦地将\Person{圣迪狄摩}的一本论圣神的著作译出（上述摘录即取自该译本序言），以期让不懂希腊文的人——\Person{圣热罗尼莫}如此希望——认出其剽窃之处。\par

\Person{鲁菲努斯}竭力为\Person{圣盎博罗削}辩护，并指出其对手的诸多矛盾之处，说道：\Quote{\Person{圣盎博罗削}撰写论圣神之书，不仅用言辞，更以自己的鲜血写成，因为他向迫害者献出生命的血，且在自己体内流淌；然而天主保全了他的生命，以从事未来的劳苦事工。}\par

事实是，\Person{圣盎博罗削}精通希腊文，既已着手撰写论圣神的著作，便研读了前人之作，并采用了\Person{圣巴西略}、\Person{圣迪狄摩}等人所论及的论点。伟大\Person{圣奥斯定}对此论著的意见，可对抗\Person{圣热罗尼莫}的意见。他说：\Quote{\Person{圣盎博罗削}在论述圣神这一深邃主题，并阐明圣神与圣父及圣子同等时，却采用了浅显的论述风格；因为他的主题无需华丽的辞藻，而是需要推动读者心灵的证据。}\par

% structure: p0010 source=h2 target=subsection reason=relative-heading-rank; base=h2

\subsection{正文}

\Emph{致格拉齐安皇帝。}\par

\Person{基德红}的被选，是吾主\OriginalTerm{降生成人}{Incarnation}的预象；山羊羔的祭献，预表基督身体为罪所作的补赎；牛犊的祭献，则预表废除世俗礼仪；而那三百名兵士，预表未来藉十字架而完成的救赎。\Person{基德红}寻求种种征兆，本身也是一个奥迹：羊毛的变干与变湿，预表\Place{犹太人}的失落与\Place{外邦人}的蒙召；用盆子接盛的水，预表为宗徒们洗脚。\Person{圣盎博罗削}祈求洗涤自己的污秽，并称颂\Person{基督}的慈爱。由天主之子所遣发的同一水泉，引发奇妙的皈依；然而，这水泉不能由任何其他遣发，因为它是圣神的倾注，而圣神不受任何外在权能支配。\par

1. 正如我们所读到的，当\Person{耶鲁巴耳}在橡树下打麦子\BibleRef{民 6:11}时，他领受了天主的讯息，为能将天主的子民从外族的权势下引入自由。他被选以承受\Emph{\OriginalTerm{恩宠}{grace}}，这并不足为奇，因为即使在那个时候，他已被安置在圣十字架与可敬之\Emph{智慧}的荫庇下，处于未来降生成人的预定奥迹中，正从藏匿之处打出可见的可结果实的谷粒，且\Emph{[奥妙地]}将圣徒中的选民从虚空的秕糠之渣滓里分开来。因为这些选民，如同受过真理之杖的训练，舍弃了旧人及其行为的冗余，被聚集到教会内，犹如进入榨酒池。因为教会是永恒泉源的榨酒池，因为从她那里涌流出天上葡萄树的汁液。\par

2. \Person{基德红}被这信息所感动，当他听说虽然成千上万的人失败了，天主却要通过一个人将祂的子民从敌人手中拯救出来，\BibleRef{民 6:14}便献上了一只小山羊，并按照天使的话，将肉和无酵饼放在磐石上，又把汤汁浇在上面。天使用手中杖的尖端一碰，火就从磐石中发出，于是他所献的祭品便被焚尽了。\BibleRef{民 6:19}由此看来，那磐石显然是\OriginalTerm{基督圣体}{the Body of Christ}的预像，因为经上写道：\BibleQuote{他们饮了那伴随他们的磐石，那磐石就是基督。}\BibleRef{格前 10:4}这当然不是指祂的\OriginalTerm{天主性}{His Godhead}，而是指祂的肉身，那肉身以祂\OriginalTerm{圣血}{His Blood}的永恒之流，灌溉了渴慕者的心灵。\par

3. 即使在那个时候，这奥秘就已宣示：主耶稣在祂的肉身上，被钉十字架时，将除免全世界的罪恶，不仅是肉身的恶行，也包括灵魂的私欲。因为小山羊的肉象征行为的罪，汤汁象征欲望的诱惑，正如经上写道：\BibleQuote{百姓起了贪欲，说：\InnerQuote{谁给我们肉吃？}}\BibleRef{户 11:4}天使伸出他的杖，触摸磐石，有火从中发出，\BibleRef{民 6:21}这显示主的\OriginalTerm{肉身}{Flesh}，充满\OriginalTerm{圣神}{the Divine Spirit}，将要烧尽人类脆弱的一切罪过。因此主也说：\BibleQuote{我来是把火投在地上。}\BibleRef{路 12:49}\par

4. 于是，那受教并预知未来之事的人，遵守了天上的奥秘，因此照告诫，宰杀了父亲原来预备献给偶像的公牛犊，而自己却将一头七岁大的牛犊献给天主。\BibleRef{民 6:26}他这样做，极其清楚地表明：在主来临之后，所有外邦人的祭献都要废除，唯有主苦难的\OriginalTerm{祭献}{sacrifice}，才应为人民的救赎而奉献。因为那头牛犊是基督的预像，正如\Person{依撒意亚}所说，在祂内蕴藏着\OriginalTerm{圣神七恩}{the seven gifts of the Spirit}的圆满。\BibleRef{依 11:2}\Person{亚巴郎}在看见了主的日子时，也曾献上这牛犊，并满心欢喜。\BibleRef{若 8:56}那一位就是曾在不同预像下被祭献的：有时为小山羊的预像，因为祂是赎罪祭；有时为绵羊的预像，因为祂是不抵抗的祭品；有时为牛犊的预像，因为祂是无玷的祭品。\par

5. 圣\Person{基德红}预先看到了这奥秘。接着他挑选了三百人作战，以此显示世界得以摆脱更凶恶敌人的侵袭，不是靠人数的众多，而是靠十字架的奥秘。尽管他勇敢而忠诚，但他仍向上主请求未来胜利的更充分证据，说：\BibleQuote{上主啊，如果你真要用我的手拯救以色列，正如你所说的，看，我要把一块羊毛毯放在打谷场上，若是羊毛上有露水，而遍地都是干的，我就知道你必照你的许诺，藉我的手拯救百姓。事实果然如此。}\BibleRef{民 6:36}之后，他又请求让露水降在遍地，而羊毛却是干的。\par

6. 也许有人会问，他岂不是显得缺乏信德吗？他在得到许多标记的指示后，还要求更多。但是，那以奥秘方式说话的人，怎能显得像是疑惑或缺乏信德呢？他并非疑惑，而是谨慎，以免我们疑惑。因为他的祈祷既有效果，他怎能疑惑呢？他若不明白天主的信息，又怎能毫无畏惧地开始作战呢？羊毛上的露水象征犹太人的\OriginalTerm{信德}{faith}，因为天主的话如同露水降下。\par

7. 因此，当全世界因外邦迷信而枯干时，那来自天上眷顾的露水便降在羊毛上。但在\Place{以色列}家迷失的羊\BibleRef{玛 15:24}（我认为这正是犹太羊毛预像所预示的），我说，那些羊\BibleQuote{拒绝了活水的泉源}\BibleRef{耶 2:13}之后，滋润信德的露水在犹太人心中枯竭了，神圣泉源便将它的水流转向外邦人的心田。由此，如今全世界都被信德的露水所滋润，而犹太人却失去了他们的先知和顾问。\par

8. 他们遭受无信的干旱，并不奇怪，因为上主已剥夺了预言之雨的滋润，说：\BibleQuote{我将命令我的云不再向那葡萄园降雨。}\BibleRef{依 5:6}因为有一种施救\OriginalTerm{恩宠}{grace}的滋润之雨，正如\Person{达味}也说：\BibleQuote{他降临如落在羊毛上的雨，又如滴在地上的甘露。}天主的\WorkTitle{圣经}应许我们，在救主降临时，这雨将遍及全地，以\OriginalTerm{圣神}{the Divine Spirit}的露水灌溉世界。如今，主已经来到，雨水也已降下；主带着天上的甘露一同来临，因此我们这些先前干渴的人，现在得以畅饮，以内在的畅饮汲取那\OriginalTerm{圣神}{that Divine Spirit}。\par

9. 圣\Person{基德红}预见了这一点，即外邦民族也将通过接受信德而畅饮，因此他更仔细地询问，因为圣人们的谨慎是必要的。甚至\Person{农}的儿子\Person{若苏厄}，当他看见天军统帅时，也询问：\BibleQuote{你是来帮助我们，还是来帮助我们的敌人？}\BibleRef{苏 5:13}以免他或许被仇敌的诡计所欺骗。\par

10. 他把羊毛放在打谷场上，而不是田野或草地上，这并非没有理由，因为那里是麦子的收成所在：\BibleQuote{庄稼多，工人少。}\BibleRef{路 10:2}因为通过对主的信德，即将收获一个德行丰硕的收成。\par

11. 再者，他使犹太人的羊毛变干，将羊毛上的露水挤入盆中，使盆盛满了水，却不用那露水洗自己的脚，这并非没有缘由。如此伟大奥迹的特权，应留给另一位。我们期盼的那一位，惟有祂能洗净万民的污秽。\Person{基德红}不够伟大，不足以将这奥迹归于自己，但\BibleQuote{人子来不是受服事，而是服事人。}\BibleRef{玛 20:28} 因此，我们要认清这些奥迹在谁身上得以完成。不在\Person{基德红}身上，因为那时尚在开始阶段。所以外邦人起初被越过，因为旱像仍在外邦人身上，而以色列便胜过他们，因为那时露水存留在羊毛上。\par

12. 现在我们来看天主的福音。我发现主脱下祂的衣裳，以手巾束腰，向盆里倒水，洗门徒们的脚。\BibleRef{若 13:4} 那水正是天上的露水，这早已预告，即主耶稣基督要在那属天的露水中洗祂门徒的脚。如今，让我们心灵的脚也伸展出来。主耶稣也愿意洗我们的脚，因为祂不仅对\Person{伯多禄}，而是对每一位信友说：\BibleQuote{我若不洗你，你就与我无分。}\BibleRef{若 13:8}\par

13. 主耶稣，请来吧，脱下祢为我而穿的衣裳；请祢脱下，好以祢的仁慈覆盖我们。请为我们以手巾束腰，好以祢不朽的恩赐束起我们。请将水倒入盆中，不仅洗我们的脚，也洗我们的头；不仅洗身体的，也洗灵魂的足迹。我愿意脱去我们软弱的一切污秽，使我也能说：\BibleQuote{夜间我脱了长衣，怎能再穿上？我洗了脚，怎能再玷污？}\BibleRef{歌 5:3}\par

14. 何等的高超啊！作为仆人，祢洗门徒们的脚；作为天主，祢从天上降下露水。祢不仅洗脚，还邀请我们与祢同席，并以祢的尊威为我们立下榜样，劝勉我们说：\BibleQuote{你们称我 \InnerQuote{师傅}、\InnerQuote{主}，说得正对，我原来就是。若我为主、为师傅的，给你们洗脚，你们也该彼此洗脚。}\BibleRef{若 13:13}\par

15. 那么，我也愿意洗我弟兄们的脚，我愿意履行我主的诫命，我不以自己为耻，也不轻看祂先前所做的。谦卑的奥迹何其美好，因为当我洗去他人的污秽时，也洗去了我自己的。但并非所有人都能领会这奥迹。\Person{亚巴郎}确实愿意洗脚，\BibleRef{创 18:4} 但那是出于好客之情。\Person{基德红}也曾愿意给显现给他的主的天使洗脚，但他的意愿只限于一位；他愿意，像是服侍，而非赐予与自己共融。这是一个无人知晓的伟大奥迹。最后，主对\Person{伯多禄}说：\BibleQuote{我所做的，你现在还不明白，但以后就会明白。}\BibleRef{若 13:7} 我说，这是一个神圣的奥迹，就连那些洗脚的人也将探询。因此，我们得以堪当与基督有份，并非凭藉那属天奥迹的普通水。\par

16. 也有一种水，我们倾注到自己灵魂的盆中，这水来自羊毛，来自\BibleBook{}{民长纪}，也来自\BibleBook{}{圣咏集}。这是自天而来的信息之水。主耶稣啊，就让这水进入我的灵魂，进入我的肉身，藉着这雨的滋润，让我们心灵的峡谷和内心的田野得以青绿。愿祢的甘露降临我身，倾注恩宠与不朽。洗净我心灵的脚步，使我不再犯罪。洗净我灵魂的脚跟，使我能抹去诅咒，使我在灵魂的脚上感受不到蛇的咬伤\BibleRef{创 3:15}，而如祢曾吩咐跟随祢的人那样，以未受伤的脚践踏蛇和蝎子\BibleRef{路 10:19}。祢已救赎了世界，也请救赎一个罪人的灵魂。\par

17. 这正是祢慈爱的非凡卓越之处，祢藉此一个接一个地救赎了整个世界。\Person{厄里亚}曾受差遣往一位寡妇那里去；\Person{厄里叟}洁净了一人；主耶稣啊，祢在今天却洁净了成千上万的人。在\Place{罗马}城有多少，在\Place{亚历山大里亚}有多少，在\Place{安提约基}有多少，在\Place{君士坦丁堡}又有多少！因为君士坦丁堡也领受了天主的话，并领受了祢审判的明证。当她将亚略派的毒害怀在心中时，便常受邻近战争的惊扰，四周敌对的武器回响不绝。但当她拒绝那些与信仰相异的人后，她竟以恳求者的身份接纳了那位她一向畏惧的敌人——君王的审判者，在他死后将他埋葬，并使他入墓。那么，祢在君士坦丁堡洁净了多少人，最后，在今日的全世界，祢又洁净了多少！\par

18. \Person{达玛稣}未洁净，\Person{伯多禄}未洁净，\Person{盎博罗削}未洁净，\Person{额我略}未洁净；因为我们执行职务，但圣事是祢的。赐予属神之物的能力不在于人，而是，主啊，是祢和圣父的恩赐，正如祢藉先知们所说的：\BibleQuote{我要将我的神倾注在一切有血肉的人身上：你们的儿子和女儿要说预言。}\BibleRef{岳 2:28} 这就是那预像中的天上之露，这就是那恩泽的雨，正如我们读到：\Quote{恩泽的雨，为祢的产业而分施。} 因为圣神不受任何外在权势或法律的支配，而是自己自由的宰制者，按照自己意志的决断分施一切，如我们读到的，随自己的心意，各别地分施给人。\BibleRef{格前 12:11}\par

% structure: p0031 source=h2 target=subsection reason=relative-heading-rank; base=h2

\subsection{第一章。}

圣\Person{盎博罗削}以赞颂皇帝开始他的论证，既因皇帝的信德，也因将大殿归还教会；然后他力言，他的对手若承认圣神不是仆人，便不能否认祂高于万有，又补充说，当同一个圣神说\Quote{万物都侍奉祢}时，已清楚表明祂与受造物有区别；他也用其他证据证实了这一点。\par

19. 所以，圣神并不在万有之中，而是在万有之上。因为（最仁慈的皇帝，你既已充分领受了关于天主圣子的教导，甚至能亲自教导他人），我就不再耽搁你了，因为你渴望并要求更精确地[关于祂]得知一些事，尤其最近你对这类论据如此满意，以至于在无人敦促的情况下，你便下令将\OriginalTerm{大殿}{Basilica}归还给教会。\par

20. 因此，我们获得了你信德的恩宠和我们自己信德的赏报；因为我们只能将此归于圣神的恩宠：当所有人都没有意识到时，你突然归还了\OriginalTerm{大殿}{Basilica}。我要说，这是恩赐，这是圣神的工程，祂那时确实被我们宣讲，却在你内工作。\par

21. 我并不懊悔先前的损失，因为那次扣押\OriginalTerm{大殿}{Basilica}的结果获得了一种利息般的收获。因为你扣押\OriginalTerm{大殿}{Basilica}，是为证明你的信德。所以你的虔诚实现了它的意图，它扣押是为了证明，又如此证明以致归还。我没有失去果实，并且我得到你的判断，而向所有人显明：尽管行动有所不同，在你内却没有意见的分歧。我要说，向所有人显明：扣押不是出于你自己，归还却是出于你自己。\par

22. 现在让我们用证据来确立我们刚才所说的。讨论的第一点是：万有都服事。万有都服事是明显的，因为经上记载：\BibleQuote{万有都服事你。}这是圣神藉先知说的。祂没有说我们服事，而是说\BibleQuote{服事你，}好让你相信祂自己被排除在服事之外。所以，既然万有都服事，而圣神不服事，圣神当然不包含在万有之中。\par

23. 因为如果我们说圣神包含在万有之中，那么当我们读到圣神通晓天主的深奥事理，\BibleRef{格前 2:10} 时，无疑就是否认天主圣父在万有之上。因为圣神既出于天主，又是祂口中的圣神，既然天主——圣神属于祂——在万有之上，拥有确然的圆满完美和完美能力，我们怎能说圣神包含在万有之中呢？\par

25. 但为避免反对者认为宗徒犯了错误，就让他们学习他所追随、作为其信仰权威的那一位。主在福音中说：\BibleQuote{当\OriginalTerm{护慰者}{Paraclete}来临时，就是我从父那里派遣给你们的，那发自父的真理之神，祂要为我作证。}\BibleRef{若 15:26} 所以圣神既发自父，又为子作证。因为这既忠信又真实的见证者为父作证，对于表达天主的威严，没有什么见证比这更充分；对于天主权能的合一，没有什么比这更显明，因为圣神与子有同样的知识，子是父奥秘的见证者和不可分离的分享者。\par

26. 因此，祂把受造物的团体和数目排除在天主的知识之外，但因不排除圣神，祂显示圣神不属于受造物的团体。所以在福音中所读的那段话：\BibleQuote{从来没有人见过天主，只有在父怀里的\OriginalTerm{独生子}{Only-begotten Son}，祂已彰显了祂。}也与排除圣神有关。因为那连天主的深奥事理都洞悉的，怎会没有见过天主呢？那知道属天主的事的，怎会没有见过天主呢？那出于天主的，怎会没有见过天主呢？所以，既然确定从来没有人见过天主，而圣神却见过祂，显然圣神被排除了。因此，那被排除在万有之外的，是在万有之上。\par

% structure: p0040 source=h2 target=subsection reason=relative-heading-rank; base=h2

\subsection{第2章}

\BibleQuote{万有都是藉着祂而造成的}这句话，并不能证明圣神包含在万有之中，因为圣神并非受造。否则，用别的经文也能证明，子，甚至圣父自己，都必须被列在万有之中，这将是类似的不敬。\par

27. 仁慈的皇帝，这似乎是对我们正确感受的完整叙述，但对不虔敬者似乎并非如此。请看他们所追求的是什么。因为异端者惯常说，圣神应算在万有之中，因为论到天主圣子经上写着：\BibleQuote{万有都是藉着祂而造成的。}\BibleRef{若 1:3}\par

28. 一个不持守真理、陷入颠倒陈述次序的论据，会是多么混乱！因为这样的论据，若要证明圣神在万有之中，除非他们证明圣神是受造的，才有效力。因为圣经说，一切受造的都是藉着子而造成的；但我们既然未受教导说圣神是受造的，就绝对不能证明祂在万有之中，祂既非如万有那样受造，也非被创造。对我而言，这个证据有助于确立每一点：第一，它证明祂在万有之上，因为祂不是受造的；第二，因为祂在万有之上，可见祂不是受造的，也不应算在那些受造物之中。\par

29. 但是，若有人因为圣史说万有都是藉着\OriginalTerm{圣言}{Word}而造成的，而没有将圣神排除在外（虽然天主的圣神在\Person{若望}内说：\BibleQuote{万有都是藉着祂而造成的}，并没有说我们这些受造物都是万有；而主自己明确指示，天主的圣神在圣史们内说话，说：\BibleQuote{因为说话的不是你们，而是你们父的圣神在你们内说话}），\BibleRef{玛 10:20} 然而，若有人如我所说，在此没有排除圣神，而将祂列入万有之中，那么他相应地也就不会在宗徒说：\BibleQuote{可是为我们，只有一个天主，就是父，万有都出于祂，我们也归于祂。} \BibleRef{格前 8:6} 那段话中排除天主的子。但为了让他知道子不在万有之中，让他读随后的经文，因为当他接着说：\BibleQuote{并且只有一个主，就是耶稣基督，万有都藉着祂而有}，\BibleRef{格前 8:6} 他当然从万有中排除了天主的子，而这人也排除了圣父。\par

30. 然而，贬损圣父、圣子或圣神的尊威，是同等的不敬。因为凡不信圣子者，也不信圣父；凡不信圣神者，亦不信天主子；没有真理的规范，信德亦无法存立。因为凡开始否认圣父、圣子与圣神的能力一体者，必然无法在不分割之处证明其分裂的信德。因此，既然完全的虔敬在于正确地相信，那么完全的不虔也就在于错误地相信。\par

31. 因此，那些认为应将圣神列入万有之中的人，因为他们读到万有都是藉着圣子而造成的，也就必然认为圣子应被列入万有之中，因为他们读到：\BibleQuote{万有都出于天主。}\BibleRef{格后 5:18} 但结果，他们若不把圣子从一切受造物中区分出来，也就不会把圣父从万有中区分出来，因为正如万有都出于圣父，同样，万有也都是藉着圣子。而宗徒因在圣神内的先见之明，正是用了这句话，以免那些不虔敬的人，在听到圣子说过：\BibleQuote{我父将赐给我的，比一切都大，}\BibleRef{若 10:29} 以后，似乎把圣子也列入了万有之中。\par

% structure: p0047 source=h2 target=subsection reason=relative-heading-rank; base=h2

\subsection{第三章}

宗徒所说的万有都出于圣父、藉着圣子，并不使圣神与他们的团体分离，因为归于一位的，也归属于每一位。因此，那些因基督之名受洗的人，若相信三位，就被视为因圣父和圣神之名受洗；否则，这圣洗便无效。这同样适用于因圣神之名受洗。如果因一段经文就将圣神与圣父及圣子分离，那么从其他经文必然推出父从属于子。圣子受天使崇拜，而非受圣神崇拜，因为圣神是圣子的见证者，而非仆役。圣经说到圣子在万有之先，这是就受造物而言。圣神的伟大尊威由不得赦免亵渎圣神的罪所证明。这罪如何不得赦免，以及圣神如何是唯一的。\par

32. 但也许有人会说，作者之所以说万有都出于圣父，万有都藉着圣子，\BibleRef{格前 8:6} 却未提及圣神，是有其理由的，并由此得到论证的立足点。但如果他固执于这种乖谬的解释，他将在多少经文中发现，圣经虽未提及圣父或圣子，却断言了圣神的德能，只是留待领悟而已？\par

40. 那么，若圣神的恩宠被断言，难道天主圣父或独生子的恩宠就被否认了吗？绝非如此；因为正如父在子内，子在父内，同样，\BibleQuote{天主的爱，藉着赐给我们的圣神，已倾注在我们心中了。}\BibleRef{罗 5:5} 正如在基督内蒙福的人，是因圣父、圣子及圣神之名而蒙福，因为名号是一，德能是一；同样，论及任何属神的工程，无论是圣父的，或圣子的，或圣神的，它不仅归于圣神，也归于圣父和圣子，不仅归于圣父，也归于圣子和圣神。\par

41. 同样，甘达刻女王的厄提约丕雅太监，在基督内受洗时，获得了圆满的圣事。而那些自称不知道有什么圣神的人，虽然声称自己已受若翰的洗，后来仍受洗了，因为若翰是为了罪之赦，因即将来临的耶稣之名施洗，而非因自己的名。所以他们不认识圣神，因为按若翰施洗的方式，他们并没有因基督之名受洗。原来若翰虽未以圣神施洗，却已宣讲了基督和圣神。后来，当他被问及自己是否就是基督时，他回答说：\BibleQuote{我以水洗你们，但有一位比我更强的要来，我就是解他的鞋带也不配，他要用圣神和火洗你们。}因此，他们既未因基督之名受洗，也未怀有对圣神的信德，便不能领受这圣洗圣事。\par

42. 于是他们因主耶稣基督之名受洗，针对他们，圣洗并未重复举行，而是另行施予，因为圣洗只有一个。但在没有完整的圣洗圣事之处，便不被视为有开始，也谈不上任何形式的洗礼。唯有宣认圣父、圣子及圣神，圣洗才是完整的。若否认其中一位，你就推翻了整体。正如你在言语上只提及一位，或圣父，或圣子，或圣神，但在你的信仰中并不否认圣父、圣子或圣神，信德的奥迹便是完整的；同样，你虽然指名圣父、圣子和圣神，却削减其中一位——或父、或子、或圣神——的德能，整个奥迹便成了空虚。最后，那些曾说过：\BibleQuote{我们连有没有圣神都没有听到过，后来便因主耶稣基督之名受了洗。}这是恩宠额外的丰盛，因为现在通过保禄的宣讲，他们认识了圣神。\par

43. 这也不该被视为与此矛盾：即虽然有时未提及圣神，圣神仍被相信，而言语未提及的，在信仰中获得了表达。因为当说\Quote{因我们的主耶稣基督之名}时，藉着这名号的一体性，奥迹得以圆满，而圣神并不与基督的圣洗分离，因为若翰是施洗以使人悔改，基督则是在圣神内施洗。\par

44. 现在让我们思考：正如我们读到，以基督之名举行的圣洗圣事是完整的，那么，当单独提及圣神时，奥迹的完整性是否会有所欠缺？让我们继续论证：提及一位，即已指明天主圣三。如果你提及基督，你便同时暗示了以圣神傅了圣子的圣父、受傅的圣子本身，以及圣子受傅时所凭的圣神。因为经上记载：\BibleQuote{这位纳匝肋的耶稣，天主以圣神傅了祂。}\BibleRef{宗 10:38} 如果你提及圣父，你同样意指祂的圣子和祂口中的圣神，只要你心中明白这一点。如果你提及圣神，你也同样提及圣父——圣神由祂所发——以及圣子，因为圣神也是圣子的神。\par

45. 因此，为使权威也能与理性结合，圣经指出我们也可以正当地在圣神内受洗，正如主所说：\BibleQuote{至于你们，要受圣神的洗。}\BibleRef{宗 1:5} 在另一处，宗徒又说：\BibleQuote{因为我们众人都在这身体内受洗，归于一位圣神。}\BibleRef{格前 12:13} 工程是同一的，因为奥迹是同一的；洗礼是同一的，因为为世界的死亡是同一的；因此，工作是一体的，宣示也是一体的，二者不可分离。\par

46. 但若在此处，圣神被分离于父与子的运作，因为经上说：\BibleQuote{万物都出于天主，万物都藉着圣子}\BibleRef{格前 8:6}，那么，同样地，当宗徒论及基督说：\BibleQuote{祂是在万有之上，世世代代应受赞美的天主}\BibleRef{罗 9:5} 时，他便不仅将基督置于所有受造物之上，而且（这是亵渎的话）也置于圣父之上。但绝对不可，因为圣父不属于万有，也不属于祂自己受造物之群。整个受造界在下，而万有之上是圣父、圣子和圣神的天主性。前者服侍，后者统治；前者从属，后者为王；前者是工程，后者是工程的创造者；前者无一例外地朝拜，后者则无一例外地受朝拜。\par

47. 最后，关于圣子，经上记载：\BibleQuote{天主的众天使都要朝拜祂。}\BibleRef{希 1:6} 你找不到\Quote{让圣神朝拜}这样的话。再往后：\BibleQuote{天主曾向哪一位天使说过：\InnerQuote{你坐在我右边，等我使你的仇敌作你的脚凳}？众天使岂不都是服役的神，被派遣服侍吗？}\BibleRef{希 1:14} 当他提到\Emph{众}时，包括圣神吗？当然不，因为天使和其他大能者注定要服侍并服从天主子。\par

48. 但事实上，圣神不是仆役，而是圣子的见证者，正如圣子亲自论及祂说：\BibleQuote{祂要为我作证。}\BibleRef{若 15:26} 因此，圣神是圣子的见证者。见证者知晓一切，正如圣父是见证者。因为你在后文中读到：\BibleQuote{我们所得的救恩，是由天主以征兆、奇事、各种异能，以及圣神的分配，与祂一同作证而证实了的。}\BibleRef{希 2:3} 那位随己意分配者，无疑在万有之上，而非在万有之中，因为分配是行动者的恩赐，而非工程本身固有的部分。\par

49. 如果圣子在万有之上——我们的救恩藉着祂得以开始，以便被宣讲——那么，圣父，那位以征兆和奇事为我们的救恩作证并予以确证者，也无疑被排除在万有之外。同样，圣神，那位以祂各种恩赐为我们的救恩作证者，不应被列入受造物之群，而应与圣父和圣子并列；祂在分配时，并不因分割自己而被分割，因为祂是不可分割的，当祂赐予万众时毫无损失，正如圣子在圣父接受国度时毫无损失\BibleRef{格前 15:24}，圣父将属祂的交给圣子时也毫无亏损。因此，我们从圣子的证据得知，在属神恩宠的分配中毫无损失；因为那位随意嘘气的\BibleRef{若 3:8} 处处都无损失。关于这一权能，我们将在后文更充分地论述。\par

50. 同时，既然我们的意图是按适当次序证明圣神不应被列在万有之中，就让我们以宗徒——他们质疑其言论——作为这一立场的权威。因为何为\Quote{万有}，无论是可见的还是不可见的，他亲自声明说：\BibleQuote{因为在天上和地上，万有都是在祂内受造的。}\BibleRef{哥 1:16} 你看到，\Quote{万有}指的是天上的事物和地上的事物，因为天上也有不可见的受造物。\par

51. 但为了使无人不知，他补充了所指的对象：\BibleQuote{无论是上座者、宰制者、率领者、掌权者，万有都是藉着祂，并且为了祂而受造的；祂在万有之先就有，万有都赖祂而存在。}\BibleRef{哥 1:16} 那么，他在这里把圣神包括在受造物之中了吗？或者说，当他说天主子在万有之先时，难道要认为他说祂在圣父之先吗？绝非如此；因为正如他在这里说万有是藉着圣子受造，天上万有都赖祂而存在，同样，无法怀疑天上万有都在圣神内拥有其力量，因为我们读到：\BibleQuote{因上主的话诸天造成，因祂口中的圣神万象生成。}\BibleRef{咏 33:6} 那么，那位是天上地上万物力量之源者，就是在万有之上。因此，那位在万有之上的，当然不侍奉；不侍奉的是自由的；自由的拥有主权的特权。\par

52. 如果我一开始就提出这一点，它会被否认。但正如同他们否认较小的，以免人相信较大的，让我们也先提出较小的事，好使他们要么在较小的事上显露背信，要么，如果他们承认较小的事，我们就可以从较小的推论较大的。\par

53. 我认为，最仁慈的皇帝，那些胆敢将圣神归入万物的人已彻底被驳倒了。但为使他们知道他们不仅被宗徒们的证言、也被我们的主的证言所逼压；他们怎么还敢将圣神算作万物之一呢？因为主亲口说过：\BibleQuote{凡出言干犯人子的，可得赦免；但出言干犯圣神的，在今世及来世，都不得赦免。}\BibleRef{玛 12:32}那么，谁还敢把圣神列于受造物之中呢？或者谁竟如此盲目，认为干犯任何受造物都不会不得赦免呢？因为如果犹太人因敬拜天上的万象而被夺去了天主的保护，而敬拜并宣认圣神的人却为天主所接纳，那不宣认祂的人则被定为无可赦免的亵圣之罪：由此必然推出，圣神不能列于万物之中，而是超越万有，冒犯祂的惩罚是永罚。\par

54. 但请仔细留意主为什么说：\BibleQuote{凡出言干犯人子的，可得赦免；唯独出言干犯圣神的，在今世及来世，都不得赦免。}\BibleRef{玛 12:32}干犯圣子与干犯圣神有所不同吗？因为他们的尊威同属一体，共同无别，所以冒犯也是一样的。但若有人因可见的人性躯体而被误导，对基督的身体（对我们来说，它不应被视为低贱的，因为它是贞洁的宫殿，童贞女的果实）有所轻视，他虽犯了罪，但并非被拒于怜悯之外，藉着信德仍可获得宽恕。然而，若有人否认圣神的尊威、威能和永远的能力，并认为魔鬼不是因天主的神而是因\OriginalTerm{贝尔则步}{Beelzebub}被逐出，那么，在充满亵渎之处，便无法获得宽恕；因为凡否认圣神的人，也就否认了圣父和圣子，因为同一位圣神就是天主的神，也是基督的神。\par

% structure: p0065 source=h2 target=subsection reason=relative-heading-rank; base=h2

\subsection{第四章}

圣神就是那位曾在先知和宗徒们内说话的那一位，也是天主的神及基督的神；并且圣经称祂为护慰者、生命之神和真理之神。\par

55. 没有人会怀疑圣神是唯一的，尽管许多人曾怀疑天主是否唯一。因为许多异端者曾声称旧约的天主是一位，新约的天主是另一位。但正如圣父是唯一的，祂如我们读到的，在古时曾藉着先知向祖先发言，而在这末期则藉着祂的圣子向我们发言；\BibleRef{希 1:1} 圣子也是唯一的，按照旧约的记载，祂曾被亚当触怒，\BibleRef{创 3:17} 曾显现给亚巴郎，\BibleRef{创 18:22} 曾受雅各伯敬拜；\BibleRef{创 28:17} 同样，圣神也是唯一的，祂曾在先知们内工作，\BibleRef{伯后 1:21} 曾嘘气于宗徒们，\BibleRef{若 20:22} 并在圣洗圣事中与圣父和圣子相结合。\BibleRef{玛 28:19} 达味论到祂说：\Quote{不要从我身上收回祢的圣神。}又在另一处论到祂说：\Quote{我往何处去躲避祢的神？}\par

56. 为使你知道天主的神与圣神是同一的，正如我们在宗徒书信中读到的：\BibleQuote{除非受天主的神感动，没有一个人会说\InnerQuote{诅咒耶稣}；除非受圣神感动，也没有一个人能说\InnerQuote{主耶稣}。}\BibleRef{格前 12:3}宗徒称祂为天主的神。祂也称祂为基督的神，如你读到的：\BibleQuote{你们已不属于肉性，而是属于圣神，只要天主的神住在你们内。}\BibleRef{罗 8:9}随后又说：\BibleQuote{如果那使耶稣从死人中复活者的神住在你们内。}\BibleRef{罗 8:11}因此，同一位就是天主的神，也就是基督的神。\par

57. 同样，祂也是生命之神，正如宗徒所说：\BibleQuote{因为在基督耶稣内生命之神的法律，已使我脱离罪与死亡的法律。}\BibleRef{罗 8:2}\par

58. 这位宗徒称为生命之神的，主在福音中则称为护慰者，也即真理之神，如你所见的：\BibleQuote{我也要求父，祂必赐给你们另一位\OriginalTerm{护慰者}{Paraclete}，使祂永远与你们同在：祂是世界所不能领受的真理之神，因为世界看不见祂，也不认识祂。}\BibleRef{若 14:16}如此，你们有了护慰者圣神，也称为真理之神，无形之神。那么，有些人怎么认为圣子在其天主性上是可见的呢，而世界连圣神也看不见呢？\par

59. 现在领受主的话吧：那称为真理之神的圣神，就是主所论的圣神，因为在这卷书的末了，你读到：\BibleQuote{领受圣神吧！}\BibleRef{若 20:22}伯多禄也教导我们，那位是圣神的神就是主的神，他说：\BibleQuote{阿纳尼雅！为什么撒殚充满了你的心，使你欺骗圣神呢？}\BibleRef{宗 5:3}稍后，他又对阿纳尼雅的妻子说：\BibleQuote{你们为什么共谋试探主的神呢？}\BibleRef{宗 5:9}当他说\Quote{你们}时，他表明他是在谈论他之前对阿纳尼雅所谈论的那同一位神。因此，主的神就是圣神。\par

60. 主也清楚地表明，圣父的神就是圣神，按照玛窦福音，祂说，在受迫害时，你们不要思虑该说什么：\BibleQuote{因为说话的不是你们，而是你们圣父的神在你们内说话。}\BibleRef{玛 10:20}按照圣路加福音，祂又说：\BibleQuote{你们不要思虑怎样申辩，或该说什么话，因为在那时刻，天主圣神必要教给你们应说的话。}\BibleRef{路 12:11}因此，虽然有许多神被称为神，如经上说：\Quote{祂使祂的天使成为神，}但天主的神却是唯一的。\par

61. 宗徒和先知们都领受了那同一圣神，正如\OriginalTerm{蒙选的器皿}{vessel of election}，\OriginalTerm{外邦人的导师}{Doctor of the Gentiles}所言：\Quote{因为我们众人都饮于一圣神；}\BibleRef{格前 12:13} 祂仿佛是那不可分割的，却倾注于灵魂，流入感官，好能熄灭这世俗的焦渴。\par

% structure: p0074 source=h2 target=subsection reason=relative-heading-rank; base=h2

\subsection{第五章}

圣神既然圣化受造物，那么祂既非受造物，也不受变化影响。祂常是美善的，因为祂由圣父和圣子所赐予；也不应把祂列入那些所谓会消逝的事物中。必须承认祂是美善的根源。祂是天主口唇的圣神，邪恶的改正者，且自身是美善的。最后，因为圣经称祂为美善，并在圣洗中与圣父及圣子相结合，祂就绝不可能被否认是美善的。然而，经上不说祂在美善上进步，却说是被成全，这使祂有别于一切受造物。\par

62. 因此，圣神不属于有形之物的实体，因为祂把无形可见的恩宠赋予有形之物；祂也不属于无形受造物的实体，因为它们领受祂的圣化，并藉着祂而超越宇宙的其他工程。无论是天使、宰制者，或是掌权者，每一受造物都期待圣神的恩宠。因为我们藉着圣神成为子女，正如经上所载：\Quote{天主派遣了祂圣子的圣神，进入我们心中，呼喊：\InnerQuote{阿爸，父啊！}这样，你们不再是奴仆，而是儿子了；}\BibleRef{迦 4:6} 同样，一切受造物也等待天主子女的显扬；他们确实是因圣神的恩宠而成为天主的子女。因此，一切受造物本身也将因圣神恩宠的启示而改变，\Quote{并脱离败坏的控制，得享天主子女的光荣自由。}\par

63. 因此，一切受造物都易于变化，不仅那些因某种罪过或外在元素的情况而改变了的，而且那些虽因本性缺陷可能败坏，但因谨慎的纪律尚未如此的也一样；因为，正如我们在先前论述中所指出的，天使的本性显然能够变化。确实适宜断定：一个的性质怎样，其他的性质也怎样。那么，其余的本性都能变化，但纪律更为优越。\par

64. 因此，一切受造物都能变化，但圣神是美善的，不能变化，也不因任何过错而改变；祂消除所有人的过错，并宽恕他们的罪。那么，祂怎能变化呢？祂通过圣化在别人身上促成向恩宠的改变，自己却不改变。\par

65. 祂既常是美善的，怎能变化？因为圣神，那一切美善之物藉以分施给我们的，从不邪恶。因此，两位圣史在同一处，用词虽异，却陈述了同一件事，正如你在\WorkTitle{玛窦福音}中读到：\Quote{你们纵然不善，尚且知道把好的东西给你们的儿女，何况你们在天之父，岂不更将好的赐与求祂的人？}\BibleRef{玛 7:11} 但依照\WorkTitle{路加福音}，你会发现这样写：\Quote{何况在天之父，岂不更将圣神赐与求祂的人吗？}\BibleRef{路 11:13} 那么，我们注意到，按主的判断，藉圣史的见证，圣神是美善的，因为一位用\Quote{好的}代替圣神，另一位用圣神代替\Quote{好的}。所以，如果圣神就是那美善的，祂怎能不善？\par

66. 我们也没有忽略，有些抄本同样按圣路加记载：\Quote{何况你们的天之父，岂不更将好礼物赐与求祂的人吗？} 这好礼物是圣神的恩宠，就是主耶稣被钉在十字架的刑木上，战胜死亡并夺去其权势后，凯旋归来，从天上倾注下来的，正如经上所载：\Quote{祂上升高天，俘虏了俘虏，将好礼物赐与人。} 说\Quote{礼物}说得好，因为圣子曾赐予了，经上说：\Quote{有一个婴孩为我们诞生了，有一位圣子赐给了我们；}\BibleRef{依 9:6} 同样，圣神的恩宠也赐予了。但我何必犹豫说，圣神也赐予了我们呢？既然经上写着：\Quote{天主的爱，藉着赐给我们的圣神，已倾注在我们心中了。}\BibleRef{罗 5:5} 既然被俘的心灵肯定不能领受祂，主耶稣便先俘虏了俘虏，使我们的情感获得释放，好让祂倾注属神恩宠的恩赐。\par

67. 祂说\Quote{俘虏了俘虏}说得好。因为基督的胜利是自由的胜利，为众人赢得了恩宠，未伤害任何人。因此，在释放众人时，没有人被俘。因为在主受难之际，唯有不义独行其道，它曾俘获所有它所掌握的，如今不义反转身来，自己被俘了，不再隶属于\OriginalTerm{贝里雅耳}{Belial}，而归属于基督；事奉基督便是自由。\Quote{因为谁若在主内蒙召为奴仆，就是主的自由人。}\BibleRef{格前 7:22}\par

68. 但回到主题。\Quote{人人都离弃了正道，一同败坏了；没有行善的人，实在没有一个。} 如果他们排除圣神，他们自己也承认祂不在\Quote{人人}之中；如果他们不排除祂，那么他们也承认祂在\Quote{人人}中一同败坏了。\par

69. 但让我们思考，祂自身是否具有美善，因为祂是美善的源头和根本。因为正如圣父与圣子具有美善，圣神也同样具有美善。宗徒也教导了这一点，他说：\Quote{圣神的效果是：平安、仁爱、喜乐、忍耐、良善。}\BibleRef{迦 5:22} 因为谁怀疑祂是美善的呢？祂的效果就是良善。因为\Quote{好树结好果子。}\BibleRef{玛 7:17}\par

70. 所以如果天主是善的，那么祂口中的那位圣神，就是洞察天主深奥事理的那位，怎可能不善呢？邪恶的感染怎能进入天主的深奥事理呢？由此可见，那些否认天主圣子是善的人是多么愚妄，因为他们无法否认基督的圣神是善的，而天主圣子论到祂说：\BibleQuote{因此我曾说过：他要领受我的。}\BibleRef{若 16:15}\par

71. 或者，难道圣神不善吗？祂使最坏的人成为善人，除去罪恶，毁灭邪恶，排除罪行，倾注美好的恩赐，使迫害者成为宗徒，使罪人成为司铎。经上说：\BibleQuote{从前你们原是黑暗，但现在你们在主内却是光明。}\BibleRef{弗 5:8}\par

72. 但我们何必推迟呢？如果他们要求陈述（既然他们不否认事实），就让他们听：圣神是善的，因为\Person{达味}说：\BibleQuote{愿你的善神引导我走上正直的途径。} 圣神除了充满美善，还是什么呢？祂虽然因本性而无法企及，却因着祂的良善能被我们领受；祂的能力充满万有，却只为义人所分享；祂本体单纯，德能丰盛，临在于每个人，将自己的恩赐分施给众人，而祂自己却完整地无所不在。\par

73. 天主圣子完全有理由说：\BibleQuote{你们要去使万民成为门徒，因父及子及圣神之名给他们授洗，}\BibleRef{玛 28:19} 祂并不耻于与圣神联合。那么，为什么有人厌恶那位在圣洗圣事中主都不曾耻于联合的圣神，在我们的敬礼中与父及子相提并论呢？\par

74. 所以圣神是善的，但祂的善不是获取来的，而是分施善。因为圣神不是从受造物领受，而是被领受；祂也不是被圣化，而是圣化人；因为受造物被圣化，而圣神圣化。在此事上，虽然用同一个词，但在本性上有区别。因为领受圣善的人与分赐圣善的天主都被称为圣，如经上所说：\BibleQuote{你们应是圣的，因为我是圣的。}\BibleRef{肋 19:2} 但是圣化与腐败不能分享同样的本性，因此圣神的恩宠与受造物不可能是同一实体。\par

75. 那么，既然整个不可见的受造界（有些人正确地相信其本体是有理性且无形体的），除了天主圣三之外，不是分施圣神的恩宠，而是获取，不是分享而是领受，所以整个受造界的群体就应与圣神分开。因此，让他们相信圣神不是受造物；或者，如果他们认为祂是受造物，又为何将祂与父相提并论？若他们视祂为受造物，为何将祂与天主圣子相连？但如果他们不认为祂应与父和子分离，他们就不视祂为受造物，因为圣化同一之处，本性也同一。\par

% structure: p0090 source=h2 target=subsection reason=relative-heading-rank; base=h2

\subsection{第六章}

虽然我们藉水和圣神受洗，但后者远胜于前者，因此不可将圣神与父及子分离。\par

76. 然而有许多人，因为我们藉水和圣神受洗，就认为水和圣神的功用没有区别，因此认为它们在本性上也没有差异。他们也没有注意到，我们被埋葬在水元素中，是为了能藉圣神重新复活。因为在水中是死亡的象征，在圣神中却是生命的保证，为使罪恶的身体藉着水而死，水仿佛将身体封在一种坟墓里，好使我们藉着圣神的能力，从罪恶的死亡中被更新，在天主内重生。\par

77. 因此，这三个见证是合一的，正如\Person{若望}所说：\BibleQuote{水、血和圣神。}\BibleRef{若一 5:8} 在奥迹上是合一的，不在本性上。水是埋葬的见证，血是死亡的见证，圣神是生命的见证。那么，如果水中有什么恩宠，并非来自水的本性，而是由于圣神的临在。\par

78. 我们是生活在水中呢，还是在圣神内？我们是在水中受印呢，还是在圣神内？因为我们在祂内生活，祂自己就是我们嗣业的抵押，正如宗徒写信给厄弗所人说的：\BibleQuote{在祂内你们相信以后，就受了恩许的圣神的印证；这圣神就是我们得嗣业的抵押。}\BibleRef{弗 1:13} 所以我们是在圣神内受印证，不是因本性，而是因天主，因为经上写道：\BibleQuote{那给我们傅油的是天主，祂也在我们身上盖了印，并赐下圣神在我们心里作抵押。}\par

79. 那么，我们是由天主在圣神内受印的。因为正如我们在基督内死，是为了再生，同样，我们受圣神印证，是为了能拥有祂的光辉、肖像和恩宠，这无疑就是我们的属神印记。因为虽然我们是在身体上明显地受印，实际上却是在心中受印，好使圣神在我们身上刻画出天上肖像的模样。\par

80. 那么，谁敢说圣神与父及子是分离的呢？因为藉着祂我们得以达到天主的肖像与模样，并且一如宗徒\Person{伯多禄}所说，我们得以分享天主性体。这其中当然不是肉性继承的嗣业，而是义子之恩宠的属神联系。为了叫我们知道这印记是在我们心上多于在身体上，先知说：\BibleQuote{上主，你向我们显示了你圣容的光辉，你在我心中充满了喜乐。}\par

% structure: p0097 source=h2 target=subsection reason=relative-heading-rank; base=h2

\subsection{第七章}

圣神不是受造物，因为祂是无限的，曾被浇灌在分散于各国的宗徒们身上，而且祂也圣化天使，并使我们与天使相等。同样，\Person{玛利亚}也充满了祂，主基督亦然，乃至一切或高或低的事物。而且一切祝福都源于祂的运作，如在\Place{贝特匝达}水动时所预表的。\par

81. 因此，一切受造物都被其自身本性的界限所限，而那些不可见的运作虽不能由场所和界限所限，却也被其本体属性所封闭；那么，谁敢称圣神为受造物呢？祂的能力不受限制、无所界限。因为祂常在万物中、处处存在，这无疑是天主性和主权的特质，因为经上说：\BibleQuote{大地和其中的万物，都属于主。}\par

81. 所以，当主派遣祂的仆人宗徒们时，为了让我们认识到受造物是一回事，而圣神的恩宠是另一回事，祂将他们派往不同的地方，因为所有人不能同时在各处。但祂将圣神赐给众人，将不可分割的恩宠的恩赐倾注在分散的宗徒们身上。因此，人虽不同，但工作的完成在众人内却是一样，因为圣神只有一个，论到祂说：\BibleQuote{当圣神降临于你们身上时，你们将充满圣神的德能，要在\Place{耶路撒冷}及全\Place{犹太}和\Place{撒玛黎雅}，并直到地极，为我作证人。} 见：\BibleRef{宗 1:8}\par

82. 所以，圣神是无限无界的、无穷的，祂将自己注入门徒们的心灵，他们遍布遥远地区的各分隔区域，直至全世界的遥远边界，没有任何事物能逃避或欺骗祂。因此，\Person{圣达味}说：\BibleQuote{我往何处，才能脱离你的神能？我逃往何处，才能逃避你的面容？} 对哪一位天使圣经说过这话？对哪一位宰制者？对哪一位大能者？对哪一位天使我们可曾发现其能力散播于众人之上？因为天使被派给少数人，而圣神却倾注在全体人民身上。那么，谁能怀疑那同时倾注于众人而看不见的是属神的；而那能被看见并由个人所拥有的，却是属于形体的？\par

83. 但是，正如圣化宗徒们的圣神不分享人性；同样，祂圣化天使、宰制者和大能者，也与受造物没有共通之处。但如果有人认为天使的圣德并非属灵的，而是属于其本性特质的另一种恩宠，那么他们实际上认定天使低于人类。因为他们自己也承认，不敢将天使与圣神相比，并且他们不能否认圣神倾注在人类身上；而圣神的圣化是属神的恩赐与恩惠，那么拥有更卓越圣化的人，必然被认定比天使更优越。但既然天使降到人间帮助人，就必须理解天使的本性之所以更高，是因为它领受了更多的圣神的恩宠，而赐予我们和他们的恩惠，都来自同一位赐予者。\par

84. 然而，那使人类较低的本性也能等同于天使领受的恩赐的恩宠，是何等伟大啊，正如主亲自应许说：\BibleQuote{你们要像天使一样，在天上。} 这并不是难事，因为那位在圣神内创造了那些天使的，也要用同样的恩宠使人类等同于天使。\par

85. 但是，对于哪一受造物，能说它充满万物呢，正如圣经论到圣神所记载的：\BibleQuote{我要将我的神倾注在一切有血肉的人身上。} 见：\BibleRef{岳 2:28} 这话不能用于天使。最后，\Person{加俾额尔}自己，当他奉命去见\Person{玛利亚}时，说：\BibleQuote{万福！充满恩宠者，} 见：\BibleRef{路 1:28} 明确宣告了那在她内的圣神的恩宠，因为圣神临于她，她将因天上的圣言而使她的胎充满恩宠。\par

86. 因为充满万物是主的作为，祂说：\BibleQuote{我充满天地。} 见：\BibleRef{耶 23:24} 那么，如果主充满天地，谁能断定圣神在主权和天主能力上没有份呢，因为祂充满了世界，以及全世界之外，充满了全世界的救赎主\Person{耶稣}？因为经上记载：\BibleQuote{耶稣充满圣神，从\Place{约旦河}回来。} 见：\BibleRef{路 4:1} 那么，除了那拥有同样圆满的，谁能充满那充满万物的呢？\par

87. 但为了避免他们反对说这是按肉身而言，尽管那位从肉身发出能力医治众人的，比众人更伟大；然而，正如主充满万物，同样我们读到有关圣神的话：\BibleQuote{主的神充满了世界。} 见：\BibleRef{智 1:7} 并且你会发现，论到所有与宗徒们聚集的人，这样说，\BibleQuote{他们充满圣神，勇敢地宣讲天主的圣言。} 见：\BibleRef{宗 4:31} 你看到圣神赐予圆满和勇敢，总领天使向\Person{玛利亚}宣告祂的运作，说：\BibleQuote{圣神要临于你。} 见：\BibleRef{路 1:35}\par

88. 你也在福音中读到，天使按时下到水池中，搅动池水，第一个下到水池里的，无论患什么病，必得痊愈。见：\BibleRef{若 5:4} 天使借着这个预象宣告了什么，岂不就是圣神的降临吗？这降临要在我们的时代实现，并要借着司铎的祈祷呼求，圣化诸水。因此，那位天使是圣神的先驱，因为借着圣神的恩宠，灵药将要施予我们灵魂和心智的疾病。所以，圣神拥有与天主父和基督相同的仆役。祂充满万物，拥有万物，如同天主父和子一样，在万物中运作一切。\par

89. 那么，还有什么比圣神的化工更神圣呢？因为天主亲自作证，圣神主宰他的祝福，说：\BibleQuote{我要将我的神倾注在你的后裔身上，将我的祝福降在你的子孙身上}，\BibleRef{依 44:3}。因为没有圣神的默感，任何祝福都不能圆满。因此，宗徒也认为没有比这更能为我们祝愿的了，正如他亲口说的：\BibleQuote{我们不断为你们祈祷，恳求你们充满对他旨意的认识，在一切智慧与属灵的明悟中，行事相称于主}，\BibleRef{哥 1:9}。他教导我们，这就是天主的旨意：我们应多在善行、善言和善意中行走，从而充满天主的旨意，他正是将他的圣神放在我们心中的那位。因此，如果领受圣神的人充满天主的旨意，圣父与圣子的旨意当然毫无区别。\par

% structure: p0109 source=h2 target=subsection reason=relative-heading-rank; base=h2

\subsection{第八章}

圣神只由天主赐予，然而并非全部赐给每一个人，因为除基督外，无人能完全领受他。爱德由圣神倾注，圣神以神秘的香液为预像，表明他与受造物毫无共同之处；并且，他既被说是由天主口中发出，就绝不能归于受造物，也不能归于可分之物，因为他是永恒的。\par

90. 同时注意，是天主赐予圣神。因为这不是人的工作，也不是人的恩赐；司铎所呼求的，是由天主赐予的，其中含有天主的恩赐和司铎的职务。如果宗徒\Person{保禄}断定自己不能凭自己的权威赐予圣神，并认为自己极不适合这职分，以至他愿我们被天主充满圣神，\BibleRef{弗 5:18}，谁还敢擅自将赐予这恩宠的权力归于自己呢？因此，宗徒在祈祷中发出这愿望，并未凭自己的任何权威声称有权；他渴望获得，而未敢擅自命令。\Person{伯多禄}也说，他不能强迫或束缚圣神。因为他这样说：\BibleQuote{既然如此，天主既然赐给他们同我们一样的恩宠，我是谁，怎能抗拒天主呢？} \BibleRef{宗 11:17}。\par

91. 但或许他们不为宗徒的榜样所动，那就让我们引用天主的圣言吧；因为经上记载：\BibleQuote{\Person{雅各伯}是我的仆人，我扶持他；\Person{以色列}是我所简选的，我的心灵喜悦他；我在他身上倾注了我的神。} \BibleRef{依 42:1} 主又借\Person{依撒意亚}说：\BibleQuote{上主的神临在我身上，因为他给我傅了油。} \BibleRef{依 61:1}\par

92. 那么，谁敢说圣神的\OriginalTerm{实体}{substance}是受造的呢？在他的光照下，我们在心中看见天主真理的美，以及受造物与\OriginalTerm{天主性}{Godhead}之间的差别，好使工程与其创造者区分开来。天主又对哪个受造物说过：\BibleQuote{我要倾注我的神}呢？\BibleRef{岳 2:28} 他没有说\Quote{神}，而是说\Quote{我的神}，因为我们不能领受圣神的圆满，而是按照我们的师傅按自己的旨意，从自己分施给我们的那么多。\BibleRef{斐 2:6} 因为正如天主子不以自己与天主同等为僭越，却空虚了自己，使我们能在心智中领受他；但他空虚自己，并非抛弃了自己的圆满，而是好叫那我所不能承受的圆满，按我所能的度量，将自己注入我内；同样，圣父也说，他将这神倾注在所有血肉之上；因为他并非整个地倾注，而是所倾注的，足以丰裕地分给众人。\par

93. 因此，圣神倾注在我们身上，但当主耶稣具有人的形体时，圣神却住在他身上，正如经上记载的：\BibleQuote{你看见圣神从天降下，住在他身上，他就是要以圣神施洗的那位。} \BibleRef{若 1:33} 在我们周围，是施予者以丰裕的供应所显示的慷慨；在他内，则永远存留着圣神的圆满。因此，他倾注了他认为足够我们用的，而所倾注的，既未分离，也未分割；但他拥有圆满的一体，借此可按我们力所能及的程度，照亮我们心灵的眼目。最后，我们所领受的，就如我们心智的进步所能获得的这么多，因为圣神恩宠的圆满是不可分割的，而是按我们本性的容量为我们所分享。\par

94. 因此，天主倾注圣神，天主的爱也借着圣神倾注；在这一点上，我们应认识到工程与恩宠的一体性。因为正如天主倾注了圣神，同样，\BibleQuote{天主的爱，借着所赐与我们的圣神，已倾注在我们心中了。} \BibleRef{罗 5:5} 好使我们明白，圣神并非受造物，而是天主之爱的\OriginalTerm{分施者}{dispenser}和丰裕的泉源。\par

95. 同样，为使你相信所倾注的，不能与受造物相通，而是天主性所独有的，圣子的名也被倾注，如你所读到的：\BibleQuote{你的名如同倾注的香液。} \BibleRef{歌 1:3} 这句话的效力无与伦比。正如香液密封在瓶中，只要局限于瓶窄小的空间，虽不能及于多人，却仍保持其芬芳；但当香液从那封闭的瓶中倾出时，便远播广扬；同样，基督的名字在他降临以色列子民之前，也如密封在瓶中的香液，局限在犹太人的心灵内。因为\Quote{天主在\Place{犹大}显扬，他的名在\Person{以色列}彰显；}也就是说，这名字被犹太人的窄瓶局限在他们狭隘的范围内。\par

96. 那时，这名确是伟大的，即使它仍停留在软弱和少数人的窄小范围内，但尚未将其伟大倾注于外邦人的心中，直至全世界的地极。然而，在他借着来临而照耀整个世界之后，他将他那神圣的名散布到一切受造物之中，并非由任何增加所充满（因为圆满不容增添），而是填补了空虚之处，好使他的名在全世界内成为奇妙的。因此，他名的倾注，意味着一种恩宠的丰盛洋溢和天上财富的富饶；因为凡是倾注的，都是从丰盛中溢出。\par

97. 因此，正如从天主之口生发的智慧不能被称为受造，从他心中所发出的圣言不能被造，在其中蕴藏着永恒尊威圆满的能力也不能被造；同样，从天主口中倾注的圣神也不能被视为受造，因为天主亲自显示了他们的合一性，以致他提到了倾注他的圣神。由此我们明白，天主圣父的恩宠与圣神的恩宠是同一的，且毫无分裂或减损地被分施到每个人的心中。因此，那被倾注的圣神，既不割裂，也不受任何形体的限制，也不被分裂。\par

98. 因为圣神怎可能经由某种分配而被分割呢？若望论及天主说：\BibleQuote{我们所以知道，他住在我们内，是借他赐给我们的圣神。}\BibleRef{若一 3:24} 但那始终内住的，自然不改变；因此，若不受任何改变，便是永恒的。所以圣神是永恒的，然而受造物却易于犯错，因而受制于变化。凡受制于变化的，不能是永恒的，因此圣神与受造物毫无共同之处，因为圣神是永恒的，但一切受造物都是暂时的。\par

99. 宗徒同样指出圣神是永恒的，因为：\BibleQuote{如果公山羊和牛犊的血，并母牛的灰烬，洒在不洁的人身上，可使人获得肉身的洁净，那么基督的血，岂不是更……他借着永恒的圣神，把自己无瑕可疵地奉献给天主吗？}\BibleRef{希 9:13} 因此，圣神是永恒的。\par

% structure: p0121 source=h2 target=subsection reason=relative-heading-rank; base=h2

\subsection{第九章}

圣神被恰当地称为基督的油膏和喜乐之油；为何基督自己不是油膏，既然他以圣神受傅。圣神被称作油膏并不奇怪，因为父与子也被称作神。他们之间并无混淆，因为唯有基督承受了死亡，此处便论及他的救赎十字架。\par

100. 许多人曾认为，圣神即是基督的油膏。称之为油膏是合宜的，因为他被称为喜乐之油，汇集诸多恩宠，散发出甘饴的芬芳。但是，全能者天主圣父以圣神傅了他，使他成为司祭长，他与其他人在法律之下以预表方式受傅者不同，他即是按法律在肉身内受傅，又在真理中充满来自圣父在超越法律之上的圣神的德能。\par

101. 这便是那喜乐之油，关于它先知说：\Quote{天主，你的天主以喜乐之油傅了你，胜过你的同伴。} 最后，伯多禄说耶稣是以圣神受傅的，正如你所读到的：\BibleQuote{你们都知道：在若望宣讲洗礼以后，从加里肋亚开始，在全犹太所传布的话：关于纳匝肋人耶稣，天主怎样以圣神傅了他。}\BibleRef{宗 10:37} 因此，圣神就是那喜乐之油。\par

102. 他说喜乐之油，是说得极好的，免得你以为圣神是一个受造物；因为这类油的本性绝不容与另一种湿气混合。同样，喜乐也不傅抹身体，而是照亮内心深处，正如先知所说：\Quote{你使我心中欢愉。} 所以，谁若想把油与更湿的东西混合，便是徒劳无功，因为油性较轻，当其他东西沉淀时，它便上浮而分离。那些可悲的小贩怎会想起以他们的诡计将喜乐之油与其他受造物相混合呢？因为实际上，形体的东西不能与非形体的相混合，受造物也不能与非受造者相混合。\par

102. 那用来傅基督的，称为喜乐之油，是极恰当的；因为为他所寻求的，并非寻常或普通的油，或是为包扎伤口，或是为解热；因为世界的救恩并不寻求减轻他的伤痛，他疲倦身体的永恒大能也不需要恢复精力。\par

103. 如果拥有喜乐之油的那位，使濒死者欢欣，从世界除去忧愁，摧毁了哀伤死亡的气息，这并不足为奇。所以宗徒说：\BibleQuote{因为我们就是献于天主的基督的芬芳；}\BibleRef{格后 2:15} 这确然表示他是在谈论属神的事。而当天主子亲自说：\BibleQuote{上主的神临于我身上，因为他给我傅了油，}\BibleRef{路 4:18} 他便指出了圣神的油膏。因此，圣神便是基督的油膏。\par

104. 或者，由于耶稣的名如同倾注的香液，如果他们认为在油膏之名下，所表达的是基督自己，而非基督的圣神，那么，当宗徒伯多禄说主耶稣以圣神受傅时，毋庸置疑，圣神也同样被称为油膏。\par

105. 但这又有何奇？因为父与子皆被称为神。这一点，当我们开始论述圣名的合一性时将更充分地论及。然而，此处既有最恰当的位置，为免我们似乎未下结论便略过，请他们阅读：父被称为神，如主在福音中所说：\BibleQuote{天主是神；}\BibleRef{若 4:24} 而基督也被称为神，因为耶肋米亚说过：\BibleQuote{我们面前的神，即主基督。}\BibleRef{哀 4:20}\par

106. 所以，父是灵，基督也是灵，因为凡不是受造的身体就是灵，但圣神并不与父及子混合，而是与父及子有区别。因为圣神没有死，他既未取肉身就不能死，永恒的\OriginalTerm{天主性}{Godhead}也不能死，但基督按肉身而言死了。\par

107. 确实，他死在由童贞女所取的，而非由圣父所有的，因为基督是以他被钉死的那一性死去的。但圣神不能被钉，他没有肉和骨，但圣子被钉，他取了肉和骨，好使我们在那十字架上肉身的诱惑得以死去。因为他取了那不属于他的，为隐藏那本是他的；他隐藏了那本是他的，好使自己在那不属于他的中受试探，而那不属于他的得以被救赎，为借着那不属于他的，呼召我们归于那本是他的。\par

108. 啊，那十字架的神圣奥秘，在其上软弱被悬，能力得自由，邪恶被钉，胜利的旗帜高扬。所以有一位圣人说：\Quote{因敬畏你，用钉刺透我的肉躯；}他不是说用铁钉，而是用敬畏和信德之钉。因为德行的束缚比惩罚的束缚更强。最后，他的信德束缚了\Person{伯多禄}，当他跟随主直到\OriginalTerm{大司祭}{high priest}的院子时，无人束缚他，惩罚并未松开那被信德束缚的人。再者，当他被犹太人束缚时，祈祷释放了他，惩罚不能拘禁他，因为他没有背离基督。\par

109. 所以，你也要将罪钉死，好使你向罪死；向罪死的人，是为天主而活；你当为那位不怜惜自己的儿子，反倒在他的身体上钉死我们的私欲偏情者而活。因为基督为我们死了，好使我们在他复活的身体内生活。所以，不是我们的生命，而是我们的罪债在他内死了，如经上说：\BibleQuote{他在自己的身上，亲自承担了我们的罪过，上了木架，为叫我们死于罪恶，而活于正义；你们是因他的创伤而获得了痊愈。}\BibleRef{伯前 2:24}\par

110. 因此，那十字架的木头，宛如我们救恩的船，我们的通道，而非惩罚，因为除了那永恒救恩的通道外，别无救恩。当我期待死亡时，我感受不到它；当我轻看惩罚时，我不受苦；当我不在乎恐惧时，我不知其为何物。\par

111. 那么，那位我们因其创伤而得痊愈的，不就是主基督吗？关于他，同一\Person{依撒意亚}曾预言他的创伤使我们得医治，\BibleRef{依 53:5} 关于他，宗徒\Person{保禄}在书信中写道：\BibleQuote{那不认识罪的，替我们成了罪。}\BibleRef{格后 5:21} 这在他身上确实是神圣的，即他的肉身从未犯罪，受造的身体在他内也未犯罪。如果\OriginalTerm{天主性}{Godhead}本身不犯罪，这有什么稀奇，因为它没有犯罪的诱因？但如果唯天主无玷，那么确实，所有受造物就其本性而言，如我们所说，都可能易犯罪。\par

% structure: p0136 source=h2 target=subsection reason=relative-heading-rank; base=h2

\subsection{第十章}

圣神赦罪，是他与父及子所共有的，而非与天使共有。\par

112. 那么，请告诉我，无论你是何人，凡否认圣神的\OriginalTerm{天主性}{Godhead}的。圣神不可能犯罪，他反而是赦罪者。天使能赦罪吗？\OriginalTerm{总领天使}{Archangel}能吗？当然不能，唯有父、子及圣神能。因为凡能赦免的，也必能避免。\par

113. 但也许有人会说，\OriginalTerm{色辣芬}{Seraph}对\Person{依撒意亚}说：\BibleQuote{看，这炭接触了你的口唇，你的邪恶已经消除，你的罪孽已获赦免。}\BibleRef{依 6:7} 他说\InnerQuote{必除去}和\InnerQuote{必洗净}，而非\InnerQuote{我要除去}，乃是那来自天主的祭台上的火，即圣神的恩宠。因为除了圣神的恩宠，我们还能虔诚地理解为在祭台上的是什么呢？当然不是林中的木材，也不是烟灰和木炭。还有什么比依照奥迹理解此事更能符合虔诚呢？这奥迹借着\Person{依撒意亚}的口得以显示，即众人将因基督的苦难得以洁净，他按肉身而言，犹如一块炭，烧尽了我们的罪过，正如你在\BibleBook{匝加利亚}{书}所读到的：\BibleQuote{这岂不是从火中抽出来的木柴吗？而\Person{耶叔亚}身穿污秽的衣服。}\BibleRef{匝 3:2}\par

114. 最后，为使我们知道这普遍救赎的奥迹由先知们极为清晰地启示出来，我们在这里还有：\Quote{看，它已除去你的罪；}不是基督除去了自己的罪，他没有犯罪，而是在基督的肉身内全人类都应从自己的罪中释放。\par

115. 但即使\OriginalTerm{色辣芬}{Seraph}除去了罪，那也是作为天主派遣为此奥迹服务的一位。因为\Person{依撒意亚}这样说：\BibleQuote{因为有一个色辣芬被派到我这里来。}\BibleRef{依 6:6}\par

% structure: p0142 source=h2 target=subsection reason=relative-heading-rank; base=h2

\subsection{第十一章}

圣神被派遣到所有人，却不从一处移到另一处，因为他不受时间或空间的限制。他由子出发，正如子由圣父出发，他永远居于他内；也当我们领受时，他来到我们中间。他也如同圣父自己那样来临，绝不能与圣父分离。\par

116. 圣神也诚然被说成受派遣，但\OriginalTerm{色辣芬}{Seraph}只被派到一个人那里，圣神则被派到所有人。色辣芬受派去服务，圣神则行奥迹之事。色辣芬执行所命之事，圣神则随心分配。色辣芬从一处移到另一处，因为他不充满万有，反而自身被圣神充满。色辣芬按自己的本性以某种方式降临，但我们不能这样设想圣神，圣子论到他说：\BibleQuote{当\OriginalTerm{护慰者}{Paraclete}，就是我从父那里给你们派遣的，那发于父的\OriginalTerm{真理之神}{Spirit of Truth}来到时。}\BibleRef{若 15:26}\par

117. 因为如果圣神由一处出发并移到另一处，那么圣父自己也会被发现在某处，圣子也一样。如果圣父或圣子所派遣的圣神由一处出发，按照那些不虔敬的解释，圣神离开一处并前进，似乎将圣父与圣子都像某种物质身体一样留下了。\par

118. 我这样说，是针对那些主张圣神借由运动而降下的人。但圣父，祂不仅超越一切有形体的本性，而且也超越一切无形的受造物，不受任何地方的局限；圣子，作为一切受造物的创造者，超越一切受造物，也不被其自身工程的处所或时间所囿限；同样，真理之神，既是天主的神，也不受任何有形界限的拘束，因为祂是无形的，借着祂不可言喻的\OriginalTerm{天主性}{Godhead}的丰盛，远远超越全部理性的受造界，且在万有之上具有随意呼吸和随意感化的权能。\BibleRef{若 3:8}\par

119. 所以，圣神不是像从某个地方被派遣，也不是像从一个地方所由发；当祂由圣子所发时，正如圣子自己说的：\Quote{我由父而来，来到了世界上}，这句话摧毁了一切可视为由一处到另一处的想象。同样，当我们念及天主在内或在外时，我们也绝对不会把天主禁锢在任何身体内，或使祂与任何身体分离，而是在深刻而难以言喻的估量中权衡这些事，领悟到\OriginalTerm{天主性}{divine nature}的隐秘。\par

120. 最后，智慧如此说她由至高者的口中发出\BibleRef{德 24:5}，并非处于父之外，而是与父同在；因为\Quote{圣言与天主同在}\BibleRef{若 1:1}；不单与天主同在，也在天主内；因为祂说：\Quote{我在父内，父在我内}。然而，当祂从父出来时，祂并非离开某个地方，也不是像物体与物体分离；当祂在父内时，也不是像物体被包裹在物体里。同样，圣神由父和子所发时，也既不与父分离，也不与子分离。因为那身为祂口中的圣神的，怎能与父分离呢？这确实是祂永恒性的明证，也表达了这一\OriginalTerm{天主性}{Godhead}的合一。\par

121. 所以，那身为祂口中圣神的一位，是永远存在且常居不变的；但当我们领受祂时，祂似乎降下来，为能住在我们内，使我们不致远离祂的恩宠。对我们就祂似乎是降下来了，并不是祂实际下降，而是我们的心灵上升归向祂。关于这点，我们本可更详细讲述，只是我们还记得在前一部论著中已阐明：父曾说：\Quote{让我们下去，混乱他们的语言}\BibleRef{创 11:7}，以及子曾说：\Quote{谁爱我，必遵守我的话，我父也必爱他，我们要到他那里去，并在他那里作我们的住所}\BibleRef{若 14:23}。\par

122. 因此，圣神降临就如父降临一样，因为父在哪里，子也在哪里，子在哪里，圣神也在哪里。所以，不应设想圣神是单独降临的。祂的降临并非从一处到另一处，而是从秩序的部署到救赎的保全，从赐生命的恩宠到圣化的恩宠，使我们从地上转移到天上，从困苦到光荣，从奴役到国度。\par

123. 所以，圣神降临就如父降临。因为子曾说：\Quote{我和父要降临，并在他那里作我们的住所}\BibleRef{若 14:23}。父是以有形身体的方式降临吗？因此，圣神也这样降临；当祂降临时，父和子的完全临在就在祂内。\par

124. 但谁能将圣神与父及子分离呢？因为我们甚至无法不提圣神而称呼父及子。\Quote{除非受圣神感动，没有一个人能说\InnerQuote{耶稣是主}的}\BibleRef{格前 12:3}。因此，如果我们除非在圣神内，不能称耶稣为主，那么没有圣神，我们一定不能宣扬祂。但如果连天使也宣扬耶稣是主——除非在圣神内，没有人能这样宣扬——那么，圣神的职分也在他们内运作。\par

125. 因此，我们证明了父、子和圣神的临在与恩宠是一体的；这临在与恩宠是如此地属天且神圣，以至于子为此向父致谢说：\Quote{父啊！天地的主宰！我称谢你，因为你将这些事瞒住了智慧和明达的人，而启示给小孩子}\BibleRef{玛 11:25}。\par

% structure: p0154 source=h2 target=subsection reason=relative-heading-rank; base=h2

\subsection{第12章}

父、子和圣神的平安与恩宠合一，祂们的爱德也是合一的，这爱德尤其显示在救赎世人上。祂们与人的相通也是合一的。\par

126. \Emph{因此}，既然召唤是同一的，恩宠也是同一的。最后，经上写道：\Quote{愿恩宠与平安由我们的父天主，和我们的主耶稣基督赐给你们}\BibleRef{罗 1:7}。这样，你们看到，圣经告诉我们，父与子的恩宠是同一的，父与子的平安也是同一的；而这恩宠与平安乃是圣神的果实，正如宗徒亲自教导我们的：\Quote{圣神的效果却是仁爱、喜乐、平安、忍耐……}\BibleRef{迦 5:22}。平安是美好而必需的，使人不因可疑的争论而纷扰，也不被肉欲的风暴所动摇，却能使他的情感保持平静，以纯朴的信德和心灵的宁静来敬拜天主。\par

127. 关于平安，我们已经证明了这一点；至于恩宠，匝加利亚先知说，天主应许将恩宠与怜悯的神倾注在耶路撒冷\BibleRef{匝 12:10}；宗徒伯多禄也说：\Quote{你们悔改吧！你们每人要以耶稣基督的名字受洗，好赦免你们的罪过，并领受圣神的恩宠}\BibleRef{宗 2:38}。这样，恩宠也来自圣神，如同来自父和子一样。因为若没有圣神，怎能有恩宠？因为所有神圣的恩宠都在圣神内。\par

128. 我们也不仅读到父、子和圣神的和平与恩宠，忠信的皇帝啊，也读到爱与共融。因为论到爱，经上记载说：\BibleQuote{愿主耶稣基督的恩宠，和天主的爱} \BibleRef{格后 13:14}。我们听说过父的爱。这同样的爱既是父的，也是子的。因为祂自己说过：\BibleQuote{谁爱我，我父也必爱他，我也要爱他。} \BibleRef{若 14:21} 圣子的爱又是什么呢？岂不就是祂为我们牺牲自己，用自己的血救赎了我们吗？\BibleRef{弗 5:2} 然而同一的爱也在父内，因为经上记载说：\BibleQuote{天主竟这样爱了世界，甚至赐下了自己的独生子} \BibleRef{若 3:16}。\par

129. 所以，父赐下了子，而子也把自己赐给了我们。爱得以保全，应得的感情并未受损害，因为放弃时并无痛苦，感情就不会受伤害。祂献出了那甘心乐意的，献出了那自献的；父并没有将圣子交付惩罚，而是交付恩宠。如果你探询这行为的功德，就请查看这情感的描述。蒙选的器皿清楚地表明了这神圣之爱的合一，因为父赐下了子，而子也把自己赐给了我们。父赐下子，\BibleQuote{他没有怜惜自己的儿子，反而为我们众人把他交出了} \BibleRef{罗 8:32}。论到子，他又说：\BibleQuote{他为我舍弃了自己} \BibleRef{迦 2:20}。他说，\InnerQuote{他舍弃了自己}。如果是出于恩宠，我怎能挑剔？如果是祂受了错待，我亏欠得更多。\par

130. 但要明白，正如父赐下了子，子也把自己赐给了我们，同样，圣神也把他赐给了我们。因为经上记载说：\BibleQuote{那时，耶稣被圣神引入旷野，为受魔鬼的试探。} 这样，慈爱的圣神也赐下了天主圣子。因为父与子的爱是一体的，同样，我们也已指出，天主的这份爱藉着圣神已经倾注在我们心中，并且是圣神的效果，因为\BibleQuote{圣神的效果却是仁爱、喜乐、平安、忍耐。} \BibleRef{迦 5:22}\par

131. 父与子之间有共融，这是明显的，因为经上记载说：\BibleQuote{我们与父和他的儿子耶稣基督相通} \BibleRef{若一 1:3}。在另一处又说：\BibleQuote{愿圣神的相通常与你们众人相偕} \BibleRef{格后 13:14}。那么，如果父、子、圣神的和平是一体的，恩宠是一体的，爱是一体的，共融是一体的，那么作为也必然是一体的；作为既是一体的，能力就必然不能分割，实体也不能分离。因为，如果不然，同一作为的恩宠又怎能和谐一致呢？\par

% structure: p0162 source=h2 target=subsection reason=relative-heading-rank; base=h2

\subsection{第十三章。}

圣盎博罗削从圣经中指明，天主三位格的名号是一体的，首先指明圣子与圣神的名号的合一，因为二者都被称为护慰者和真理。\par

132. 那么，谁还敢否认名的合一呢，既然他看见了作为的合一？但我何必用论证来维护名的合一，因为有天主声音的明证表明父、子和圣神的名是一体的？因为经上记载说：\BibleQuote{你们要去使万民成为门徒，因父及子及圣神之名给他们授洗} \BibleRef{玛 28:19}。祂说，\Quote{因……之名}，而不是\Quote{因……诸名}。所以，父的名并不是一个，子的名另一个，圣神的名又一个，因为天主是唯一的；名并不是多个，因为没有两位天主，或三位天主。\par

132. 而且，为显示天主性是一体的，尊威也是一体的，因为父、子和圣神的名是一体的，圣子并非以一名号而来，圣神也非以另一名号而来，主亲自说：\BibleQuote{我因我父的名而来，你们却不接纳我；如果有人因自己的名而来，你们反而接纳他。} \BibleRef{若 5:43}\par

133. 圣经清楚表明，属于父的名号，同样也属于子，因为主在出谷纪中说：\BibleQuote{我要在你面前宣告我的名号；我要以我的名号在你们面前呼喊主。} \BibleRef{出 33:19} 所以，主说祂要以祂的名号呼喊主。因此，\Quote{主}是父与子的名号。\par

134. 但既然父与子的名号是一体的，那么也要明白，圣神的名号同样也是一体的，因为圣神是因子之名而来，如经上记载：\BibleQuote{但那护慰者，就是父因我的名所要派遣来的圣神，他必要教训你们一切} \BibleRef{若 14:26}。然而，那因子之名而来的，当然也因父的名而来，因为父与子的名是一体的。这样一来，父与子的名号也就是圣神的名号。因为在天下人间，没有赐下别的名字，使我们赖以得救的。\BibleRef{宗 4:12}\par

155. 同时祂也表明，必须教导天主名号的合一，而非差异，因为基督是以合一的名号而来，但假基督却要因自己的名而来，正如经上记载：\BibleQuote{我因我父的名而来，你们却不接纳我；如果有人因自己的名而来，你们反而接纳他。} \BibleRef{若 5:43}\par

156. 所以，这些经文清楚地教导我们，父、子及圣神的名号并无差别；那属于父的名号也属于子，同样那属于子的名号也属于圣神，因为子也被称为护慰者，正如圣神一样。因此主耶稣在福音中说：\BibleQuote{我也要求父，他必会赐给你们另一位护慰者，使他永远与你们同在，即真理之神。} \BibleRef{若 14:16} 祂说：\Quote{另一位}，说得好，以免你们以为子就是圣神，因为名号是合一的，并非子与圣神的\OriginalTerm{撒伯里乌式的混淆}{Sabellian confusion}。\par

因此，\Person{若望}称圣子为\OriginalTerm{护慰者}{Paraclete}，圣神是另一位护慰者；因为若望称圣子为护慰者，如你所见：\BibleQuote{若有人犯了罪，我们在父那里有护慰者，就是耶稣基督。}\BibleRef{若一 2:1}同样地，既然名字同一，权能也同一；因为护慰者圣神在哪里，圣子也在哪里。\par

因为正如主在此处说圣神将永远与信友同在，他在别处也表明他自己将永远与宗徒们同在，说：\BibleQuote{看，我同你们天天在一起，直到今世的终结。}\BibleRef{玛 28:20}因此圣子和圣神是一，天主圣三的名号是一，临在是一而不可分的。\par

但正如我们指明圣子被称为护慰者，同样我们也指明圣神被称为真理。基督是真理，圣神是真理，因为你在\BibleBook{若望}{一书}中找到：\BibleQuote{因为圣神是真理。}\BibleRef{若一 5:7}那么，圣神不仅被称为真理之神，也被称为真理，正如圣子也被宣认为真理，祂说：\BibleQuote{我是道路、真理、生命。}\BibleRef{若 14:6}\par

% structure: p0173 source=h2 target=subsection reason=relative-heading-rank; base=h2

\subsection{第十四章}

圣经中说，\OriginalTerm{天主圣三}{Trinity}的每位都是光。\Person{依撒意亚}称圣神为火，这火的预象曾在\Person{梅瑟}所见的荆棘丛中、在火舌中、在\Person{基德红}的水罐中出现。同一位圣神的\OriginalTerm{天主性}{Godhead}不可否认，因为祂的运作与父和子的运作相同，祂也被称为上主面容的光与火。\par

但我为何要论证父是光，子亦是光，圣神亦是光呢？这当然关乎天主的大能。因为天主是光，正如\Person{若望}所说：\BibleQuote{天主是光，在祂内没有黑暗。}\BibleRef{若一 1:5}\par

但圣子也是光，因为\BibleQuote{生命是人的光。}\BibleRef{若 1:8}而圣史为了表明他所论的是天主子，论及\Person{洗者若翰}说：\BibleQuote{他不是光，而是[被派遣]来为光作证。那普照每人的真光，正在进入这世界。}\BibleRef{若 1:9}因此，既然天主是光，天主子是真光，那么天主子毫无疑问是真天主。\par

你在别处也发现天主子是光：\BibleQuote{那坐在黑暗和死影中的人民，看见了一道皓光。}\BibleRef{依 9:2}但更清楚的是，经上说：\BibleQuote{因为在你那里有生命的泉源，藉你的光明，我们才能看到光明，}这意味着，天主全能父啊，你是生命之泉，在你的子——他是光——内，我们将看见圣神的光。正如主自己所表明的，说：\BibleQuote{领受圣神吧，}\BibleRef{若 20:22}在别处又说：\BibleQuote{能力从他身上出来。}\BibleRef{路 6:19}\par

但谁还能怀疑父是光呢？我们读到他的子是永恒光明的辉耀。确实，除了父，子还能是谁的辉耀呢？他永远与父同在，永远闪耀，不是以不同的，而是以同样的光辉。\par

\Person{依撒意亚}表明圣神不仅是光，也是火，说：\BibleQuote{以色列的光必将成为火。}\BibleRef{依 10:17}因此先知们称他为燃烧的火，因为在这三点上我们更强烈地看到天主性的尊威；因为圣化属于天主性，照亮是火与光的属性，而天主性常以火的形象被指示或显现：\BibleQuote{因为我们的天主是吞噬的烈火，}正如\Person{梅瑟}所说。\BibleRef{申 4:24}\par

因为他自己曾在荆棘丛中看见了火，并听到了天主的声音，当时有声音从火焰中对他说：\BibleQuote{我是\Person{亚巴郎}的天主，\Person{依撒格}的天主，\Person{雅各伯}的天主。}\BibleRef{出 3:6}那声音从火中发出，那声音在荆棘丛内，而火却毫无伤害。因为荆棘在燃烧却没有烧毁，因为在那项奥迹中，主显明他要来光照我们肉身的荆棘，他不是来消灭那些处于苦难中的人，而是要减轻他们的苦难；他要用圣神和火施洗，以赋予恩宠并消灭罪恶。\BibleRef{玛 3:11}所以，在火的象征中，天主保持他的用意。\par

在\WorkTitle{宗徒大事录}中，当圣神降临在信友身上时，也看见了火的形象，因为你读到：\BibleQuote{忽然从天上来了一阵响声，好像暴风刮来，充满了他们所在的全座房屋；有火舌般的形象显现给他们，分开落在他们每人头上。}\BibleRef{宗 2:2}\par

正是基于同样的理由，当\Person{基德红}快要战胜\Place{米德杨}人时，他命令三百人带上瓦罐，并将点燃的火把藏在瓦罐内，右手拿着号角。我们的前辈们保存了从宗徒们领受的解释：瓦罐象征我们泥土所造的身体，若它们因圣神的恩宠热忱而燃烧，并用响亮的口舌宣认主耶稣的受难，就无所畏惧。\par

那么，既然哪里有圣神的恩宠，哪里就有天主性的显现，谁还能怀疑圣神的天主性呢？从这些证据中，我们推断的不是分歧，而是天主能力的合一。因为，在所有行动的效果皆为同一的地方，能力又怎能被分割呢？\par

那么，这火是什么呢？它当然不是由普通的树枝堆成，也不是以燃烧林中的芦苇而呼啸作响的火，而是像炼金一样锻炼善行、像焚烧秸秆一样焚毁罪恶的火。这无疑就是圣神，他被称为天主面容的火与光；光，正如我们上面所说：\BibleQuote{上主，你面容的光已在我们身上留下印记。}那么，这被印记的光，不就是圣神的印记吗？相信他，\BibleQuote{你们就受了，}他说，\BibleQuote{预许的圣神的印证。}\BibleRef{弗 1:13}\par

170. 正如有天主圣容的光，同样，从天主的面容也发出火来，因为经上说：\BibleQuote{在祂眼前有烈火焚烧。}审判之日的恩宠预先照耀，好使宽恕随之而来，以赏报圣人们的服事。啊，圣经何等丰盛，无人能以人的才智领悟！啊，神圣合一的最大明证！因为这两节经文指出了多少事啊！\par

% structure: p0186 source=h2 target=subsection reason=relative-heading-rank; base=h2

\subsection{第十五章}

圣神与圣父和圣子同等是生命；事实上，无论是提到圣父——生命之泉与祂同在——或是圣子，那泉源都只能是圣神。\par

171. 我们已经说过，圣父是光，圣子是光，圣神是光；我们还应当知道，圣父是生命，圣子是生命，圣神也是生命。因为\Person{若望}说：\BibleQuote{论到那从起初就有的生命之言，就是我们听见过，亲眼看见过，瞻仰过，亲手触摸过的；这生命已显示出来，我们看见了，也为此作证，且把这与父同在，且已显示给我们的永生，传报给你们。} \BibleRef{若一 1:1} 他既说了生命之言，又说了生命，为要表明圣父和圣子都是生命。因为生命之言不就是天主之言吗？藉着这话，天主和天主之言都表明是生命。正如说生命之言是生命，同样，生命之神也是生命。\par

172. 现在你们要明白，正如圣父是生命之泉，同样，许多人也指出圣子意指生命之泉；以致有人说：全能的天主，和祢同在，祢的圣子就是生命之泉。那就是圣神的泉源，因为圣神就是生命，正如主所说的：\BibleQuote{我对你们所说的话，就是神，就是生命。} \BibleRef{若 6:64} 因为圣神在哪里，哪里就有生命；生命在哪里，圣神也在哪里。\par

173. 然而，许多人认为，在这段经文中，泉源只是指圣父。但是，让他们注意圣经所说的话：\BibleQuote{在祢那里有生命的泉源。}意即，圣子与圣父同在；因为圣言与天主同在，祂在起初就与天主同在，且与天主同在。\par

174. 然而，无论有人把这里的泉源理解为圣父还是圣子，我们当然不是指那受造之水的泉源，而是指那天主恩宠的泉源，即圣神的泉源，因为祂是活水。因此主说：\BibleQuote{若是你知道天主的恩赐，并知道对你说\InnerQuote{给我水喝}的是谁，你或许早求了祂，而祂也赐给了你活水。} \BibleRef{若 4:10}。\par

175. 这就是\Person{达味}的灵魂所渴慕的水。牝鹿渴望这些水的泉源，而不是渴慕毒蛇的毒液。因为圣神恩宠之水是活的，能洁净心灵的内在，洗除灵魂的一切罪恶，净化隐秘过犯的污秽。\par

% structure: p0193 source=h2 target=subsection reason=relative-heading-rank; base=h2

\subsection{第十六章}

圣神就是那浇灌奥秘的\Place{耶路撒冷}的大河。祂与祂的泉源——即圣父与圣子——是平等的，正如圣经所表明的。圣\Person{盎博罗削}自己渴慕那水，并告诫我们，要把它保存在我们内，就必须躲避魔鬼、贪欲和异端，因为我们的器皿是脆弱的，并且必须弃绝破裂的蓄水池，好能效法\Place{撒玛黎雅}妇人和古圣祖的榜样，寻得主的水。\par

176. 但或许有人会反对，仿佛圣神是微小的，并由此试图建立大小上的差异，争辩说水似乎只是泉源的一小部分；虽然取自受造物的例证似乎绝不适用于天主性体；然而，为使人们不致因这取自受造物的比较而做出不公正的判断，他们应当明白，圣神不仅被称为水，也被称为河，正如我们读到的：\BibleQuote{从他的腹中要流出活水的江河。但这话是指那相信祂的人将要领受的圣神说的。} \BibleRef{若 7:38}。\par

177. 因此，圣神就是那条河，且是丰沛的河；按希伯来人的记载，这河从\Person{耶稣}流到地上，正如我们领受了\Person{依撒意亚}口中预言的。\BibleRef{依 66:12} 这就是那条常流不息的伟大之河。不仅是一条河，更是一条水流充沛、满溢宏大的河，正如\Person{达味}所说：\BibleQuote{河流的支流使天主的城欢乐。}\par

178. 因为那座城，即天上的\Place{耶路撒冷}，并非由任何地上河流的河道灌溉，而是由那从生命之泉所发出的圣神灌溉；我们只需浅尝一口便得饱饫，而祂在那些天上的座天使、宰制天使、异能天使、天使和总领天使中间，似乎更丰沛地涌流，以圣神七德的全速奔流。因为如果一条涨溢出两岸的河流就能泛滥，那么，超越一切受造物之上的圣神，当祂触及我们心灵中仿佛是低洼的田地时，岂不更以祂圣化的大能丰产，使那受造物的天性欢乐吗？\par

179. 你们也不要因这里说\BibleQuote{江河}\BibleRef{若 7:38}，或别处说\BibleQuote{七神}\BibleRef{默 5:6}而感到困惑，因为这些圣神的七种恩赐，正如\Person{依撒意亚}所说的\BibleRef{依 11:2}，因着它们的圣化，意指一切德能的圆满；就是上智和明达之神，超见和刚毅之神，明识和孝爱之神，以及敬畏天主之神。因此，河是一条，但圣神恩赐的河道却有许多。这条河，乃是从生命之泉流出的。\par

180. 这里，你也不可让自己的思想转向低级的事物，因为泉源与河流之间似乎有些分别，然而神圣的圣经已经安排妥当，使人悟性的软弱不致因语言的卑微而受损。请在你面前设想任何一条河流，它从泉源涌出，却具有同一本性、同一光辉与美丽。你也要正确地断言，圣神与天主圣子、与天主圣父具有同一性体、光辉与荣耀。我将把所有一切都归结于属性的同一，并且将不畏惧任何关于伟大差异的疑问。因为在这一点上圣经也为我们作了准备；因为天主圣子说：\BibleQuote{谁若喝了我赐给他的水，他将永远不渴；并且我赐给他的水，将在他内成为涌到永生的水泉。}\BibleRef{若 4:14} 这井显然是圣神的恩宠，是从活泉源涌出的溪流。那么，圣神也是永生的泉源。\par

181. 那么，你从祂的话语中可以看出，神圣伟大的合一性已被指明，即使是异端者也不能否认基督是泉源，因为圣神也被称为泉源。正如圣神被称为河流，圣父也说过：\BibleQuote{看哪！我要在她身上广赐和平，有如河流一般；我要赐给她万国的荣耀，好似泛滥的溪流。}\BibleRef{依 66:12} 谁又能怀疑天主圣子是生命之河，从祂涌出了永生的溪流呢？\par

182. 那么，这水是好的，就是圣神的恩宠。谁会将这泉源赐予我的胸怀？愿它在我的内里涌出，愿那赐予永生的水流经我。愿那泉源溢出我们，却不流走。因为智慧说：\BibleQuote{你当饮你自己池里的水，饮你自己井里的活水；你的泉水岂可流到街市？}\BibleRef{箴 5:15} 我该如何守住这水，使它不流走，使它不滑失呢？我该如何保全我的器皿，以免任何罪恶的裂缝渗入，让永生之水渗出呢？主耶稣，请教诲我们，正如祢教诲祢的宗徒们那样，说：\BibleQuote{你们不要在地上为自己积蓄财宝，因为在地上有虫蛀，有锈蚀，在地上也有贼挖洞偷窃。}\BibleRef{玛 6:19}\par

182. 因为祂暗示，那贼就是不洁的灵，他无法进入那些行走在善工光明中的人，但若他发现有人陷于世俗欲望的黑暗中，享受世俗享乐，他便夺去他们永恒美德的一切花朵。因此，主说：\BibleQuote{你们该在天上为自己积蓄财宝，因为那里没有虫蛀，没有锈蚀，那里也没有贼挖洞偷窃。因为你的财宝在哪里，你的心也必在哪里。}\par

183. 我们的锈是放纵，我们的锈是贪欲，我们的锈是奢侈，它们以恶习的污秽蒙蔽心智的锐利眼光。再者，我们的蛀虫是亚略，我们的蛀虫是波提努，他们以不虔撕裂教会的圣衣，企图分离天主的能力不可分割的合一，以亵渎的牙齿啃咬信德的宝贵帕子。若亚略留下了他的齿痕，水便溅洒；若波提努在他的器皿上刺入毒刺，水便流走。我们不过是普通的泥土，很快就感受到罪恶。但没有人对陶匠说：\BibleQuote{你为什么这样造了我？}\BibleRef{罗 9:20} 因为，虽然我们的器皿不过是普通的，但有的受尊荣，有的受耻辱。\BibleRef{罗 9:21} 那么，不要敞开你的池子，不要以恶习和罪行挖掘，免得有人说：\BibleQuote{他掘了井又挖了坑，自己落在其中。}\par

184. 你若寻求耶稣，就当弃绝破裂的蓄水池，因为基督惯常不坐在池边，而是坐在井旁。在那里，那撒玛黎雅妇人找到了祂\BibleRef{若 4:6}，她信了，她想要汲水。虽然你本该清晨就来，但即使你来得晚，哪怕是第六时辰，你也会找到因行路疲倦的耶稣。祂疲倦了，但那是因为你，因为祂长久寻找你，你的不信长久使祂疲倦。但只要你前来，祂并不见怪；那将要施予的，却向人求饮。但祂所饮的，不是流过的溪水，而是你的救恩；祂饮的是你的善情，祂饮的是爵杯，即为你的罪作补赎的受难，好让你畅饮祂的圣血，以解此世的渴。\par

185. 所以，亚巴郎在掘了井之后赢得了天主。\BibleRef{创 21:30} 同样，依撒格在井旁散步时，迎接了那位前来归于他的妻子\BibleRef{创 24:62}，她是教会的预像。在井旁他是忠信的，在池旁他是不忠的。最后，如我们所读到的，黎贝加在井旁遇见了那寻找她的人，而妓女们则在依则贝耳的池中以血洗身。\par

\SourceRecord{Translated by H. de Romestin, E. de Romestin and H.T.F. Duckworth.；Nicene and Post-Nicene Fathers, Second Series；Vol. 10.；Edited by Philip Schaff and Henry Wace.；Buffalo, NY: Christian Literature Publishing Co.,；1896.；<http://www.newadvent.org/fathers/34021.htm>.}

% structure: p0001 source=h1 target=section reason=page-title

\section{论圣神，卷二}

% structure: p0002 source=h2 target=subsection reason=relative-heading-rank; base=h2

\subsection{引言。}

\OriginalTerm{天主性}{Godhead}的三\OriginalTerm{位格}{Persons}，无论对古代的\OriginalTerm{民长}{judges}还是对\Person{梅瑟}而言，都并非不为人知；因为\OriginalTerm{圣子}{Son}与\OriginalTerm{圣父}{Father}的同等，以及三位格彼此间的同等，不但在别处，也在他那里得到阐明。\OriginalTerm{三松}{Samson}也曾享有\OriginalTerm{圣神}{Holy Spirit}的助佑；他的事迹被提及，并显示出在某些方面是\OriginalTerm{教会}{Church}及其\OriginalTerm{奥秘}{mysteries}的预像。当圣神离开三松时，他便陷入种种灾祸；\Person{圣盎博罗削}解释了他那被剪去的发绺的灵性意义。\par

1. 即使在阅读古代史书的第一卷时，也已清楚显示出，\OriginalTerm{圣神}{Holy Spirit}的\OriginalTerm{七重恩宠}{sevenfold grace}在犹太人自己的\OriginalTerm{民长}{judges}身上闪耀，而且圣神也使人知晓了天上\OriginalTerm{圣事}{sacraments}的\OriginalTerm{奥秘}{mysteries}，\Person{梅瑟}并非不知道圣神的\OriginalTerm{永恒}{eternity}。此外，在世界之初，甚至在其开始之前，他就将至圣神与\OriginalTerm{天主}{God}相连，他知晓圣神在世界开始之前便已永恒存在。因为任何人若仔细留意，就会在起初认出\OriginalTerm{圣父}{Father}、\OriginalTerm{圣子}{Son}和圣神。关于圣父，经上记载：\BibleQuote{在起初天主创造了天地。}\BibleRef{创 1:1} 关于圣神，则说：\BibleQuote{天主的神在水面上运行。}\BibleRef{创 1:4} 在创造之初，便设立了\OriginalTerm{圣洗圣事}{baptism}的\OriginalTerm{预像}{figure}，受造物藉此得以净化，这安排十分美好。关于圣子，我们读到，正是祂将光与黑暗分开，因为有一位天主圣父发言，也有一位天主圣子行事。\par

2. 然而，为使你不致认为那位发言者的命令中有\OriginalTerm{擅取}{assumption}之意，或执行命令的那位有\OriginalTerm{次等}{inferiority}之意，圣父亲自承认圣子在执行这一工程时与祂同等，说：\BibleQuote{让我们照我们的肖像，按我们的模样造人。}\BibleRef{创 1:26} 因为共同的肖像、工程和模样，只能意味同等\OriginalTerm{尊威}{Majesty}的\OriginalTerm{同一性}{oneness}。\par

3. 但为使我能更充分认识圣父与圣子的同等：正如圣父发言，圣子创造；同样，圣父工作，圣子发言。圣父工作，如经上所载：\BibleQuote{我父到现在一直在工作。}\BibleRef{若 5:17} 论及圣子，你发现经上说：\BibleQuote{只要你说一句话，他就会好的。}\BibleRef{玛 8:8} 圣子又对圣父说：\BibleQuote{我愿我在哪里，他们也同我在一起。}\BibleRef{若 17:24} 圣父便按圣子所说的做了。\par

4. 但\Person{亚巴郎}也并非不认识\OriginalTerm{圣神}{Holy Spirit}；他看见了三位，却朝拜了一位，因为只有一位\OriginalTerm{天主}{God}、一位\OriginalTerm{主}{Lord}、一位圣神。因此，尊荣同一，因为能力同一。\par

5. 我又何必一一细说呢？因\OriginalTerm{天主的恩许}{divine promise}而生的\Person{三松}，有\OriginalTerm{圣神}{Holy Spirit}伴随他，因为经上记载：\BibleQuote{上主祝福了他，上主的神开始感动他，那时他在玛哈乃丹营。}\BibleRef{民 13:25} 就这样，他\OriginalTerm{预像}{foreshadowing}了未来的\OriginalTerm{奥秘}{mystery}，要求娶一位外邦女子为妻，这事，如经上所载，他的父母都不知道，因为这事出于上主。他理当被认为比他人更强壮，因为上主的神引导他；在这引导下，他独自一人击退了外邦的民众；另一次，面对狮子的撕咬，他毫无畏惧，凭借不可战胜的力量，徒手将狮子撕裂。但愿他保守\OriginalTerm{恩宠}{grace}能像战胜野兽那样谨慎！\par

6. 或许这不仅是\OriginalTerm{勇力的奇事}{prodigy of valour}，也是\OriginalTerm{智慧的奥秘}{mystery of wisdom}、\OriginalTerm{预言的声音}{utterance of prophecy}。因为当他前去迎娶时，遇见一只吼叫的狮子，他徒手将之撕裂，这看来并非毫无目的；在即将享受渴望的婚约时，他在狮子尸体内发现了一群蜜蜂，从它口中取了蜜，给他父母吃。相信的外邦子民拥有了蜜，从前野蛮的子民如今成了\OriginalTerm{基督的子民}{people of Christ}。\par

7. 他向同伴们提出的\OriginalTerm{谜语}{riddle}也并非没有\OriginalTerm{奥秘}{mystery}：\BibleQuote{吃的从吃者出来，甜的从强者出来。}\BibleRef{民 14:14} 这谜语在前三天中，人们徒然寻求答案，其中蕴含着奥秘；直到第七天，法律时期完结，主受难之后，这奥秘才能藉着教会的\OriginalTerm{信德}{faith}显明出来。因为你发现，宗徒们当时并不理解，\BibleQuote{因为耶稣尚未受到光荣。}\BibleRef{若 7:39}\par

8. 他们回答说：\Quote{有什么比蜜更甜？有什么比狮子更强？} \Person{三松}回答他们说：\Quote{如果你们没有用我的母牛耕地，你们绝不会猜出我的谜语。}\BibleRef{民 14:18} 哦，\OriginalTerm{神圣的奥秘}{divine mystery}！哦，\OriginalTerm{显明的圣事}{manifest sacrament}！我们逃脱了杀戮者，战胜了强者。\OriginalTerm{生命之粮}{food of life}现在所在之处，从前只有悲惨死亡的饥渴。危险化为安全，苦涩化为甘甜。\OriginalTerm{恩宠}{grace}从\OriginalTerm{过犯}{offense}中生出，能力从软弱中生出，生命从死亡中生出。\par

9. 然而，也有些人认为，若非\OriginalTerm{犹大支派的狮子}{lion of the tribe of Judah}被杀死，那婚约便无法成立；因此，在祂的身体，即\OriginalTerm{教会}{Church}内，找到了蜜蜂，它们储存智慧的蜜，因为主受难后，\OriginalTerm{宗徒}{apostles}们更充分地相信了。那么，\Person{三松}作为犹太人，杀死了这狮子，但在其中他找到了蜜，这正如那将要被赎回的产业的预像，好使剩余的人按\OriginalTerm{恩宠的拣选}{election of grace}得救。\BibleRef{罗 11:5}\par

10. 经上说：\BibleQuote{上主的神突然降在他身上，他便下到\Place{阿市刻隆}，击杀了他们中间三十个人。}\BibleRef{民 14:19} 因为他既看见了奥秘，就必定得胜。因此，那些解开并回答谜语的人，在衣服上获得了\OriginalTerm{智慧的酬报}{reward of wisdom}，作为\OriginalTerm{交往的标帜}{badge of intercourse}。\par

11. 在此，又出现别的奥秘：他的妻子被夺去，为此狐狸焚烧了外邦人的禾捆。因为外邦人的诡计往往欺骗那些反抗神圣奥秘的人。因此，在\BibleBook{}{雅歌}中又说：\BibleQuote{请给我们擒拿小狐狸，那毁坏葡萄园的，使我们的葡萄园得以茂盛。}\BibleRef{歌 2:15} 他说\Quote{小}说得好，因为大的不能毁坏葡萄园，而对于强壮者，就连魔鬼也是软弱的。\par

12. 因此，他（简略概括这故事，因为整个段落的考量留待其适当的时机）只要保有圣神的恩宠，就不可战胜，就像主所选的天主子民，那位法律之下的\OriginalTerm{纳齐尔}{Nazarite}人。参孙当时是不可战胜的，如此所向无敌，以至能用驴腮骨击杀一千人；\BibleRef{民 15:15} 他满溢天上的恩宠，干渴时竟在驴腮骨中找到了水。无论你视之为奇迹，还是翻转作奥秘，因为在谦卑的外邦人中，将有安息与凯旋，正如所记载的：\BibleQuote{有人打你的面颊，也把另一面转给他。}\BibleRef{玛 5:39} 因为藉着这忍受伤害——圣洗圣事所教导的——我们战胜了怒火的刺击，好使我们经过死亡后，得以抵达复活的安息。\par

13. 那么，那就是挣断了用皮带绞成的绳索、新绳如同细线的参孙吗？那就是只要保有圣神的恩宠，就感受不到头发被系在大梁上的束缚的参孙吗？我要说，自天主的神离他而去后，他大大改变，不再是那位穿着外邦人的战利品归来的参孙，而是从自己的伟大跌倒在女人的膝上，被爱抚并受骗，头发被剪掉了。\par

14. 那么，他头发的意义如此重大，只要留着，就能保持不可战胜的力量，一旦剃掉，那人就立刻失去所有力量吗？并非如此，我们也不能认为他头发的力量有这么大。有虔诚与信德的头发；有法律上成全的\OriginalTerm{纳齐尔}{Nazarite}人的头发，这头发因节俭和禁欲而被祝圣；她（教会的预像）曾用香液敷主脚，并以之擦拭天上圣言的脚——因为那时她也按血肉认识了基督。关于这头发有话说：\BibleQuote{你的头发如同山羊群，}\BibleRef{歌 4:1} 它长在那头上，论到这头有话说：\BibleQuote{男人的头是基督，}\BibleRef{格前 11:3} 而另一处：\BibleQuote{他的头像精金，他的发鬈如同黑色的松树。}\BibleRef{歌 5:11}\par

15. 同样，在福音中，我们的主指出有些头发是被看见且被知道的，说：\BibleQuote{就是你们的头发，也都被数过了。}\BibleRef{玛 10:30} 这暗指灵性德行的行为，因为天主并不挂虑我们的头发。但按字面相信也不荒谬，因为按照祂的天主尊威，没有什么能逃避祂。\par

16. 可是，即便天主自己知道我所有的头发，那对我有什么益处呢？那反而满溢并有益于我的，是善工的警醒见证酬报我永生的赏报。总之，参孙本人声明这些头发并非奥秘，说：\BibleQuote{如果我的头发被剃，我的力量就会离开我。}\BibleRef{民 16:17} 关于这奥秘，就说这么多，现在让我们考察这段经文的次序。\par

% structure: p0020 source=h2 target=subsection reason=relative-heading-rank; base=h2

\subsection{一}

圣神是主和权能；在这方面，祂并不次于圣父和圣子。\par

17. 前面你们读到：\BibleQuote{主祝福了他，主的神开始与他同行。}稍后又言：\BibleQuote{主的神降在他身上。}\BibleRef{民 14:6} 他又说：\BibleQuote{若我的头发被剃，我的力量就会离开我。}\BibleRef{民 16:17} 他被剃之后，请看圣经怎么说：\Quote{主，}他说，\Quote{离开了他。}\BibleRef{民 16:20}\par

18. 所以你们看到，那位与他同行的，正是那位离开他的。因此，这一位就是主，祂是主的神；也就是说，他将天主的神称为主，正如宗徒也说：\BibleQuote{主就是那神，主的神在哪里，那里就有自由。}\BibleRef{格后 3:17} 如此，你们发现圣神被称为主；因为圣神与圣子不是\OriginalTerm{一个位格}{unus}，而是\OriginalTerm{一个实体}{unum}。\par

19. 在此，他用了\Quote{权能}一词，指的是圣神。因为正如圣父是权能，同样，圣子也是权能，圣神也是权能。关于圣子，你们读过，基督是\Quote{天主的权能和天主的智慧}。\BibleRef{格前 1:24} 我们也读到，圣父是权能，如经上所载：\BibleQuote{你们将要看见人子坐在大能者的右边。}\BibleRef{玛 26:64} 他确实称圣父为权能，圣子就坐在祂的右边，如你们所读：\Quote{上主对我主说：你坐在我右边。}而且主亲自称圣神为权能，祂说：\BibleQuote{但当圣神降临于你们身上时，你们将领取权能。}\BibleRef{宗 1:8}\par

% structure: p0025 source=h2 target=subsection reason=relative-heading-rank; base=h2

\subsection{二}

圣父、圣子和圣神在旨意上同为一体。\par

20. 因为圣神本身就是权能，如你们所读：\BibleQuote{超见和刚毅（或大能）的神。}\BibleRef{依 11:2} 并且正如圣子是伟大旨意的天使，同样，圣神也是旨意的神，好使你们知道圣父、圣子、圣神的旨意是一。这旨意不是关乎任何不确定的事物，而是关乎那些预知和已定的事。\par

21. 而圣神是天主旨意的裁决者，你们从这一点就可知道。因为先前我们已指出圣神是圣洗的主，并且读到洗礼是天主的旨意，如你们所读：\BibleQuote{但法利塞人废弃了天主的旨意，没有受他的洗。}\BibleRef{路 7:30} 由此十分清楚：没有圣神就没有圣洗，同样，天主的旨意也少不了圣神。\par

22. 为使我们更完全地认识\OriginalTerm{圣神}{Spirit}就是\OriginalTerm{能力}{Power}，我们应当知道，当主说：\BibleQuote{我要将我的圣神倾注在所有有血肉的人身上。}\BibleRef{岳 2:28} 时，祂曾预许了圣神。那么，那预许给我们的，本身就是能力，正如在福音中，同一天主子声明说：\BibleQuote{我要把我父的恩许遣发到你们身上；至于你们，你们应当留在这城中，直到佩戴上来自高天的能力。}\BibleRef{路 24:49}\par

23. \OriginalTerm{圣史}{Evangelist}也如此表明圣神是能力：\Person{圣路加}叙述说，圣神带着大能降下，他说：\BibleQuote{忽然，从天上来了一阵响声，好像圣神带着大能降来。}\BibleRef{宗 2:2}\par

24. 再者，为使你不要认为这是指有形体的、可感觉的事物，要学习：圣神如此降来，正如基督将要降来一样，就如你所读到的：\BibleQuote{他们要看见人子，带着威能和大光荣，乘着天上的云彩降来。}\BibleRef{玛 24:30}\par

25. 因为，如果工作是同一个，审判是同一个，圣殿是同一个，\OriginalTerm{赋予生命者}{Life-giver}是同一个，\OriginalTerm{圣化}{sanctification}是同一个，而且父、子、圣神的国也是同一个，那么能力与大能怎会不是同一个呢？\par

% structure: p0033 source=h2 target=subsection reason=relative-heading-rank; base=h2

\subsection{第三章}

认识父与子就是生命，同样，认识圣神也是生命；因此，在\OriginalTerm{天主性}{Godhead}内，祂不可与父分离。\par

26. 那么，让他们说说，他们以为在天主的工程中有什么不同之处。因为正如认识父与子就是生命，如主亲自所说：\BibleQuote{永生就是：认识你唯一的真天主，和你所派遣来的耶稣基督。}\BibleRef{若 17:3} 同样，认识圣神也是生命。因为主说：\BibleQuote{如果你们爱我，就要遵守我的命令；我也要求父，他必会赐给你们另一位\OriginalTerm{护慰者}{Paraclete}，使他永远与你们同在：他是真理之神，世界不能领受，因为世界看不见他，也不认识他；你们却认识他，因为他与你们同在，并在你们内。}\BibleRef{若 17:14}\par

27. 因此，世界没有永生，因为它没有接受圣神；因为哪里有圣神，哪里就有永生；圣神自己就是产生永生的。因此，我奇怪为何阿里乌派提出关于唯一真天主的疑问。因为正如认识唯一真天主是永生，认识耶稣基督也是永生；同样，认识圣神也是永生；世界看不见圣神，如同看不见父，也不认识圣神，如同不认识子。但那不属于这世界的人，拥有永生，圣神——永生的光——与他永远同在。\par

28. 如果认识唯一真天主所带来的益处，与认识子和圣神的益处相同，那么你们为何在承认子与圣神带来益处的同时，却不敬拜真天主，却将子与圣神与真天主的尊荣割裂呢？因为你们必须要么相信，这是唯一真\OriginalTerm{天主性}{Godhead}的最大恩赐，并承认唯一真天主性属于父，同样也属于子和圣神；要么，如果你们说，那并非真天主的也能赐予永生，那么你们反而贬低了父，因为你们不把父的工程视为唯一真天主性的首要工程，反而视其为可与受造物的工程相比拟的。\par

% structure: p0038 source=h2 target=subsection reason=relative-heading-rank; base=h2

\subsection{第四章}

圣神赋予生命，与父和子无异，也并非经由不同的工程。\par

29. 圣神赐予生命，有何奇哉？祂与父和子一样，使人生活。谁能否认，赋予生命是永恒尊威的工程？因为经上记载：\BibleQuote{使你的仆人生活。}\par

30. 那么，我们看看，究竟是圣神被赋于生命，还是祂自己赋予生命。现今经上记载：\BibleQuote{文字叫人死，圣神却叫人活。}\BibleRef{罗 8:11} 所以，圣神赋予生命。\par

31. 为使你们明白，父、子及圣神的赋予生命并非分开的工程，要看到，赋予生命也是合一的，因为天主自己借着圣神赋予生命，如\Person{保禄}所说：\BibleQuote{那使基督从死者中复活的，也必要借那住在你们内的圣神，使你们有死的身体复活。}\par

% structure: p0043 source=h2 target=subsection reason=relative-heading-rank; base=h2

\subsection{第五章}

和父与子一样，圣神在圣经中被指明为\OriginalTerm{创造者}{Creator}，同样的真理甚至也由外教作家隐约预示，但这真理在\OriginalTerm{降生成人}{Incarnation}的奥迹中最为清楚地显示出来；在提及这奥迹之后，作者据只应归于创造者的崇拜也献给圣神这一事实，维护他的论证。\par

32. 但是，谁能怀疑圣神将生命赋予万物呢？因为祂和父与子一样，是万物的创造者；而且我们理解，全能父没有圣神便什么也不做；并且在创世之初，圣神运行在水面上。\par

33. 所以，当圣神运行在水面时，受造界尚无\OriginalTerm{恩宠}{grace}；但在这世界受造之后，经历了圣神的运作，它便获得了那恩宠的全部美善，世界因此蒙受光照。而宇宙的恩宠若没有圣神便不能存留，先知宣告说：\BibleQuote{你若收回你的圣神，他们便消灭，归于灰土。你遣发你的圣神，他们便受造，你使地面更新复兴。} 因此，他不仅教导说，没有圣神，任何受造物都不能存立，而且也教导说，圣神是整个受造界的创造者。\par

34. 谁能否认大地的创造是圣神的工作，而祂的工作正是使大地得以更新呢？因为如果那些人渴望否认大地是由圣神创造的，却又不能否认大地必须由圣神来更新，那么这些企图割裂三位格的人，就必然要主张圣神的运作超越圣父和圣子的运作，而这与真理相去甚远；因为毫无疑问，被更新的大地要比当初被创造时更为美好。或者，如果起初在没有圣神运作的情况下，圣父和圣子创造了大地，而圣神的运作是后来才加入的，那么这似乎暗示那受造物需要祂当时才添加的援助。但无论如何，谁也不该持有这种想法，即认为天主的工作会在创造者身上产生变化，这种错误是由\Person{摩尼}引进的。\par

35. 但我们是否认为，无圣神的运作，大地的实体依然存在，而无祂的工作连穹苍也无法维系呢？因为经上记着：\BibleQuote{因主的圣言，诸天已立定，万象都因祂口中的圣神而成。}请留意他所说的，诸天的一切力量都应归因于圣神。因为那位在\BibleRef{创 1:1}所记、大地尚未造成之前就已运行的，在大地被造时又岂会休息呢？\par

36. 外邦的作者们，由于无法汲取圣神的真理，仿佛是透过阴影追随我们的作者，在他们的诗句中指明，内在的灵滋养着天和地，以及月亮和星辰的灿烂天体。因此，他们不否认受造物的力量是藉由灵而存在，我们这些读到这话的人，难道要否认吗？但你们以为他们指的是由空气产生的灵。如果他们宣称一位空气的灵是万物的创造者，难道我们怀疑天主的圣神不是万物的创造者吗？\par

37. 但我为何在这些无关紧要的事上拖延呢？让他们接受一个明确的证据，即没有什么可以说不是圣神所造的；并且无可怀疑地，所有一切——无论是天使、总领天使、上座者，还是宰制者——都是藉着祂的运作而存在；因为连天使所事奉的主本身，也是由圣神降来于童贞女而受生的，正如根据玛窦所记，天使对\Person{若瑟}说：\BibleQuote{\Person{达味}之子\Person{若瑟}，不要怕娶你的妻子\Person{玛利亚}，因为那在她内受生的，是出于圣神。}\BibleRef{玛 1:20} 又根据路加所记，他对\Person{玛利亚}说：\BibleQuote{圣神要临于你。}\BibleRef{路 1:35}\par

38. 因此，由童贞女而来的诞生是圣神的工作。胎中的果实是圣神的工作，根据经上所记：\BibleQuote{你在妇女中受祝福，你的胎果亦受祝福。}\BibleRef{路 1:42} 由根而生的花同样是圣神的工作，我所说的那朵花，正是先知所预言的那朵：\BibleQuote{由\Person{叶瑟}的树干将生出一根枝条，由它的根上将发出花朵。}\BibleRef{依 11:1} 圣祖\Person{叶瑟}的根是犹太民族，\Person{玛利亚}是那枝条，基督是\Person{玛利亚}的花朵，祂准备向全世界传播信德的芬芳，从童贞的胎中绽放，正如祂亲口说的：\BibleQuote{我是原野的花，是谷中的百合。}\BibleRef{歌 2:1}\par

39. 花朵被剪下时，芬芳犹存；被压碎时，芬芳更浓；即使被撕裂，亦不失去其芬芳。同样，主耶稣在十字架的刑架上，被压碎时并未衰败，被撕裂时也并未气馁；当祂被刺入的矛所刺透时，因着倾流的圣血之色而变得更加美丽，祂仿佛重新变得俊美，自身不可能死亡，并向死者呼出永生的恩赐。因此，圣神憩息在这皇族枝条的花朵上。\par

40. 有人认为，好的枝条是主的肉躯，它从尘世的根中升到天上，遍及整个世界，结出了宗教的甘饴果实——降生成人的奥迹，将恩宠倾注在天上的祭台上。\par

41. 因此，我们不能怀疑圣神是创造者，因为我们确知祂是主降生成人的创造者。因为当你在福音的开端读到耶稣基督的诞生是这样发生的：\BibleQuote{当\Person{玛利亚}许配于\Person{若瑟}后，在同居前，她因\OriginalTerm{出于}{ex}圣神受孕。}\BibleRef{玛 1:18} 谁还能怀疑呢？\par

42. 因为虽然大多数权威抄本写作\Quote{de Spiritu}，但拉丁译者依据的希腊文却是\Quote{\GreekText{ἐχ} \GreekText{πνεύματος} \GreekText{ἁγίου}}，即\Quote{ex Spiritu Sancto}。因为那\Quote{\OriginalTerm{出于}{ex}}的，要么是指他的本质，要么是指他的能力。若指本质，则如圣子所说：\BibleQuote{我出自至高者之口；}\BibleRef{德 24:3} 又如圣神，\BibleQuote{祂由父所发；}\BibleRef{若 15:20} 关于祂，圣子说：\BibleQuote{祂要光荣我，因为祂要由我领受。}\BibleRef{若 16:14} 若指能力，则如经上所说：\Quote{只有一个天主圣父，万物都出于祂。}\par

43. 那么，\Person{玛利亚}是如何因圣神受孕的呢？如果是出于祂的本质，那么圣神就成了血肉和骨骼吗？当然不是。但如果童贞女是因祂的运作和能力而受孕，谁能否认圣神是创造者呢？\par

44. 同样，\Person{约伯}如何明确地宣称圣神是他的创造者，说：\BibleQuote{天主的圣神造了我}？\BibleRef{约 33:4} 他在短短一节经文中便指出圣神既是天主又是创造者。因此，如果圣神是创造者，祂当然就不是受造物，因为宗徒已将创造者与受造物分开，说：\BibleQuote{他们崇拜受造物，却不崇拜创造者。}\BibleRef{罗 1:25}\par

45. 他谴责那些崇拜受造物的人，教导我们应当服事创造者，因为我们应当服事创造者。而他既知圣神是创造者，便教导我们应当服事祂，说：\BibleQuote{慎防那些狗，慎防邪恶的工人，慎防自行割损的人，因为真受割损的人是我们，我们服事天主的圣神。}\BibleRef{斐 3:2}\par

46. 但若有人因拉丁抄本的不同而争论，其中有些已被异端者篡改，就让他查看希腊抄本，并注意那里写道：\Quote{\OriginalTerm{\GreekText{οἱ} \GreekText{πνεύματι} \GreekText{Θεοῦ} \GreekText{λατρεύοντες}}{\GreekText{οἱ} \GreekText{πνεύματι} \GreekText{Θεοῦ} \GreekText{λατρεύοντες}}}，翻译过来即是：\Quote{事奉天主之神的人}。\par

47. 因此，当同一位宗徒说我们应当事奉圣神，而他又声称我们不可事奉受造物，只当事奉造物主时；无疑地，他清楚表明圣神是造物主，应受与永恒天主相称的钦崇；因为经上记载：\BibleQuote{你要朝拜上主，你的天主，惟独事奉他。}\BibleRef{玛 4:10}\par

% structure: p0061 source=h2 target=subsection reason=relative-heading-rank; base=h2

\subsection{第六章}

对于那些根据\Person{亚毛斯}的话反对说圣神是受造物的人，答复是：那里所指的词是\Quote{风}，风是时常被造的，但这不能用以言说圣神，因为圣神是永恒的，不会消逝于死亡，也不会因异端所谓被圣父吸收而消灭。但他们若固执地认为这一段是指圣神而写，\Person{圣盎博罗削}指出，必须诉诸灵性解释，因为基督藉着祂的来临建立了\Quote{雷霆}，即神圣言语的力量，而\Quote{神}字实指人的灵魂以及基督所取的人性。既然此人性是天主圣三的每一位所创造的，因此可论定，那先前被确证为万有的创造者的圣神，亦即是主降生成人的根源。\par

48. 我也并非没有注意到，异端者惯于反对说，圣神看起来是个受造物，因为他们许多人利用\BibleBook{亚毛斯}{书}的那段记载来论证他们的不信，在那里他说到了风的吹动，正如先知的话清楚表明的。因为你读到这样：\BibleQuote{看，是我造成雷霆，创造风，向人宣示祂的基督，我形成光和雾，行走在大地的高处，祂的名字是\InnerQuote{上主，全能的天主}。}\BibleRef{亚 4:13}\par

49. 如果他们以此作为论证，说什么\Quote{神}是受造的，那么\Person{厄斯德拉}也教导我们说\Quote{神}是受造的，他在第四卷书中说：\Quote{在第二天，祢造成了穹苍中的风，}然而，为了不离题，在亚毛斯所言中，该段的次序岂不明示先知是在谈论此世的受造吗？\par

50. 他这样开始说：\Quote{我是上主，造雷霆，又造风。}话语的次序本身就教导了我们；因为如果他愿意谈论圣神，他肯定不会把雷霆放在首位。因为雷霆不比圣神更古老；尽管他们不虔敬，他们仍不敢那样说。而且我们看到随后论及光和雾的叙述，难道所说的不是对此世受造的理解吗？因为我们凭日常经验知道，当大地上有风暴时，先有雷霆，继之以风的吹动，天空因雾气而变黑，然后光又从黑暗中照耀出来。因为风的吹动也称为\Quote{风}，如经上所载：\Quote{火、硫磺和暴风。}\par

51. 为了让你知道他称之为\Quote{神}的那位，他说：\Quote{造雷霆，又造风。}因为这些东西在发生时，是经常被造的。但圣神是永恒的，如果有人竟敢称祂为受造物，他仍不能说祂像风一样每日被造。此外，智慧本身，按所取人性的奥秘而言，说：\BibleQuote{上主造了我。}\BibleRef{箴 8:22}虽然这是在预言将来的事，然而因为主的来临是预先定下的，所以不说\Quote{造}而说\Quote{造了我}；为使人相信耶稣的肉身是由童贞玛利亚所生，不是多次，而仅一次。\par

52. 因此，对于先知所说似乎是关于天主在雷霆和创造风中的日常作为，若将任何这类事理解为指圣神，就是大不敬；就连不虔敬的人自己也不能否认圣神在创世之前就已存在。因此，我们以虔诚的声明作证：祂永远存在，永恒常在。因为那位在世界未有之前就在水上运行的，不可能在世界受造之后才开始成为可见的；否则，就可以假定有许多圣神，他们是通过一种日常的生产过程而产生的。任何人都不应让自己污于这样的不虔敬中，竟说圣神是经常或曾经受造的。因为我不明白祂为何应经常受造；除非他们或许相信祂经常死去，因此经常受造。但生命之神如何能死去呢？如果祂不能死去，就没有理由说祂应经常受造。\par

53. 可是那些持不同想法的人陷于此亵渎之罪，他们不区分圣神；认为那被差遣的圣言回归圣父，那被差遣的圣神被重新吸收入天主内，从而出现一种重新吸收和某种相互变换，一位变为多种形式；然而，圣父、圣子与圣神之间的区别永远常存不变，保存了其权能的合一。\par

54. 但如果有人认为先知的话应就圣神解释，因为他说了\BibleQuote{向人宣示祂的基督}\BibleRef{亚 4:13}，那么他就会更容易将此解释为主的降生成人。因为若\Quote{神}字使你困惑，因而你以为这不能很好地解释为取得人性的奥秘，那么继续诵读圣经，你就会发现一切都完全与基督相符，思及祂完全适宜：祂藉着来临建立了雷霆，即天上圣经的力量与声响，我们的心灵犹如被雷霆震撼，充满惊异，从而学会敬畏，并敬重天上的神谕。\par

55. 最后，在福音中，主的兄弟们被称为\OriginalTerm{雷霆之子}{Sons of Thunder}；当圣父发出声音，说：\BibleQuote{我已经荣耀了它，还要再荣耀它。} \BibleRef{若 12:28} 犹太人便说这是向祂打雷。因为他们虽不能接受真理的恩宠，却不由自主地承认，并在无知中讲述奥迹，致使圣父为圣子作出了伟大的见证。在\BibleBook{约伯}{传}中，圣经也说：\BibleQuote{谁知道祂何时发出祂雷的威力？} 诚然，如果这些话指的是天上的雷，他便会说那力量已经发出，而非将要发出。\par

56. 因此，他将雷声归于主的话语，那话语的声音传遍普世；我们在此处将\Quote{神魂}理解为祂所取的那具有理智和完备的\OriginalTerm{灵魂}{soul}；因为圣经常用\Quote{神魂}一词指人的灵魂，如经上说：\BibleQuote{在人内造了人的神魂。} \BibleRef{匝 12:1} 同样，主也用\Quote{神魂}一词指祂自己的灵魂，当祂说：\BibleQuote{我把我的灵魂交托在祢手中。} \BibleRef{路 23:46}\par

57. 为叫你明白他指的是耶稣的降临，他又补充说，祂向世人宣示了祂的基督，因为在祂受洗时，祂宣示了祂，说：\BibleQuote{祢是我的爱子，我所喜悦的。} \BibleRef{玛 3:17} 在山上，祂宣示了祂，说：\BibleQuote{这是我的爱子，你们要听从祂！} \BibleRef{谷 9:7} 在祂受难时，当日头隐蔽、海洋和大地震动，祂宣示了祂。在百夫长身上，祂宣示了祂，百夫长说：\BibleQuote{这人真是天主子。} \BibleRef{谷 15:39}\par

58. 因此，我们应将这整段经文或单纯地理解为我们在此生活呼吸的境地，或理解为\OriginalTerm{主身体的奥迹}{mystery of the Lord's Body}；因为如果此处明言圣神是受造的，无疑圣经会在别处同样宣示，就如我们常常读到天主子，按肉身而言，祂既是受造的，也是被形成的。\par

59. 然而，我们理应在祂为我们取肉身的这一事实上，仰瞻祂的尊威，好能在祂摄取身体的这一事实上，看见祂的天主性能力。因为我们读到，圣父创造了主降生成人的奥迹，圣神也创造了它；同样我们也读到，基督自己创造了祂自己的身体。因为圣父创造了它，正如经上所载：\BibleQuote{上主创造了我，} \BibleRef{箴 8:12} 又在另一处说：\BibleQuote{天主派遣了自己的儿子，由女人所生，生于法律之下。} \BibleRef{迦 4:4} 而圣神则创造了整个奥迹，按我们所读到的，因为\BibleQuote{玛利亚因圣神受孕。} \BibleRef{玛 1:18}\par

60. 那么，你发现圣父创造，圣神也创造；也要学习，天主子也创造，当撒罗满说：\BibleQuote{\OriginalTerm{智慧}{Wisdom}为自己建造了房屋。} \BibleRef{箴 9:1} 那么，圣神创造了主降生成人的奥迹，这奥迹超越一切受造之物，祂自己又怎能是受造之物呢？\par

61. 我们在上文已一般地指出，圣神按肉身而言，在外在人方面是我们的创造者，现在让我们指出，祂按恩宠的奥迹也是我们的创造者。正如圣父创造，圣子也创造，圣神也创造，一如我们在保禄的话中所读到的：\BibleQuote{原来这是天主的\OriginalTerm{恩赐}{gift}，不是出于功行，免得有人自夸。我们是祂的化工，是在基督耶稣内受造的，为行各种善工。}\par

% structure: p0077 source=h2 target=subsection reason=relative-heading-rank; base=h2

\subsection{第七章}

圣神之于属灵的创造或重生，其为作者，不亚于圣父和圣子。此创造之卓越，及其所在。当圣经将身体或肢体归于天主时，我们当如何理解。\par

62. 所以，圣父在善工中创造，圣子亦然，因为经上记载：\BibleQuote{但是，凡接受祂的，祂给他们权能，好成为天主的子女，就是那些信祂名字的人；他们不是由血气，也不是由肉欲，也不是由人欲，而是由天主生的。} \BibleRef{若 1:12}\par

63. 同样，主自己也作证，我们按恩宠由圣神重生，说：\BibleQuote{由肉生的属于肉，因为是由肉生的；由圣神生的属于圣神，因为天主是神。你不要惊奇，因我给你说了：你们应该重生。风随意向哪里吹，你听到风的响声，却不知道风从哪里来，往哪里去：凡由圣神而生的就是这样。}\par

64. 那么，很清楚，圣神也是\OriginalTerm{圣神恩宠}{grace of the Spirit}的创作者，因为我们按天主的肖像受造，是为了成为天主的子女。所以，当祂藉着圣洁重生的义子名分，将我们接入祂的国时，难道我们要否认祂那属于祂的吗？祂使我们成为由上而来的新生的继承人，难道我们要宣称继承权而拒绝其创作者？然而，当创作者被排除在外时，恩惠便不能存留；创作者离不开恩赐，恩赐也离不开创作者。若你要恩宠，便当相信其能力；若你拒绝能力，就不要祈求恩宠。凡否认圣神的，同时也否认了恩赐。因为，若创作者被轻视，祂的恩赐又怎能宝贵？我们为何妒忌自己所领受的恩赐，减损我们的希望，弃绝我们的尊严，并否认我们的\OriginalTerm{施慰者}{Comforter}呢？\par

65. 但我们不能否认祂。我们绝不敢否认如此伟大的事，因为宗徒说：\BibleQuote{但是弟兄们，你们正如\Person{依撒格}一样，是恩许的子女；然而，正如那时那按血肉所生的，迫害了那按圣神所生的。} \BibleRef{迦 4:28} 从先前所说的可以清楚地知道，那按\OriginalTerm{圣神}{Holy Spirit}所生的。因此，那按圣神所生的，就是按天主所生的。现在，当我们更新我们内心的情感，脱去外在旧人时，我们就重生了。所以宗徒又说：\BibleQuote{但你们应在心神的精神上更新，穿上新人，就是按照天主所造的，具有真实的正义和圣善的新人。} \BibleRef{弗 4:23} 让他们听听圣经如何表达了天主化工的统一性。那在心神的精神上更新的人，就穿上了新人，就是按照天主所造的。\par

66. 因此，那更卓越的重生是圣神的工作；圣神是那按照天主的肖像所造的新人的创造者，没有人会怀疑这新人比我们这外在的人更好。因为宗徒指出，一个是属天的，另一个是属地的，当他这样说：\BibleQuote{那属天的怎样，所有属天的也怎样。} \BibleRef{格前 15:48}\par

67. 因此，既然圣神的恩宠能使那本可造成属土的东西成为属天的，我们就应该凭理性去观察，即使我们没有实例。因为在某处，圣\Person{约伯}说：\BibleQuote{我指着永生的上主起誓，祂如此判断我，并指着全能者，祂使我的灵魂愁苦（因为天主的神在我的鼻子内）。} \BibleRef{约 27:2} 他在这里用\Quote{祂的神}肯定不是指生命的气息和肉身的呼吸管道，而是指他内在之人的鼻子，借此他聚集永恒生命的芬芳，并以一种双重的感官吸取天上香膏的恩宠。\par

68. 因为据我们所读到的，有属灵的鼻子，那是\OriginalTerm{圣言}{Word}的净配所拥有的，对她说：\BibleQuote{你鼻子的气味；} \BibleRef{歌 7:8} 又在另一处：\BibleQuote{上主闻到了馨香。} \BibleRef{创 8:21} 因此，可以说，人有内在的肢体，他的手被视为在行动中，他的耳在聆听中，他的脚在某种善工的迈进中。于是，从行动中我们汇集起肢体的某种形像，因为我们不宜按肉身的方式去想像内在之人的任何事物。\par

69. 还有些人读到天主的手或手指时，就认为天主是按肉身的形像被塑造的，他们没有注意到这些记述不是为了暗示任何肉身的形状，因为在天主性内既无肢体也无部分，而是为了表达天主性的统一，好使我们相信圣子或圣神都不可能与天主圣父分离；因为天主性的圆满如同有形地居住在\OriginalTerm{天主圣三}{Trinity}的实体中。为此，圣子也被称为圣父的右手，正如我们读到：\BibleQuote{上主的右手行了大事，上主的右手高举了我。}\par

% structure: p0087 source=h2 target=subsection reason=relative-heading-rank; base=h2

\subsection{第八章。}

\Emph{圣\Person{安博}检视并驳斥异端的论点，即因为说天主在圣神内受光荣，而不是偕同圣神受光荣，所以圣神次於圣父。他指出\InnerQuote{在……内}这个小品词也可用于圣子，甚至用于圣父，另一方面，\InnerQuote{偕同}可用于受造物而不损害天主性的专有权利；事实上这些介词仅仅暗示天主三位的结合。}\par

70. 然而，如果愚妄的人甚至因音节而争论，他们因言词起争论又有何奇怪呢？因为有些人认为应当区分：天主应在圣神\Emph{内}受赞美，而非\Emph{偕同}圣神受赞美，并且他们认为天主性的伟大可由一个音节或某种习惯来估量，他们辩称，如果他们认为天主应在圣神\Emph{内}受光荣，那就是在指出圣神的某种职务，但如果说天主\Emph{偕同}圣神接受光荣或能力，他们似乎就在暗示圣父、圣子与圣神之间的某种结合与共融。\par

71. 但谁能分离那不可分离的呢？谁能分裂\Person{基督}所表明的不可分离的结合呢？\Person{祂}说：\BibleQuote{你们要去使万民成为门徒，因父及子及圣神之名给他们授洗。} \BibleRef{玛 28:19} \Person{祂}在这里关于圣父、圣子或圣神的言词或音节，有任何改变吗？当然没有。但祂说，因父及子及圣神之名。对圣神的表述与对圣父和祂自己的表述是相同的。由此推断出来的，不是圣神的任何职务，而是当我们说\InnerQuote{在圣神内}时，指的是尊荣或化工的分享。\par

72. 还要想想，你们的这种见解对圣父和圣子也是有害的，因为圣子并没有说\InnerQuote{偕同父及子及圣神之名}，而是说\InnerQuote{因…之名}，然而在这个音节中，所表达的并非任何职务，而是天主圣三的权能，\par

73. 最后，为使你们明白，并非音节有损于信德，而是信德使音节受到推崇，\Person{保禄}也讲论在\Person{基督}内。\Person{基督}并未因\Person{保禄}说在\Person{基督}内而有所减损，正如你们所看到的：\BibleQuote{我们在天主前，在基督内讲话。} \BibleRef{格后 2:17} 因此，正如宗徒说我们在\Person{基督}内讲话，照样，我们也是在圣神内讲话；正如宗徒自己所说：\BibleQuote{除非受圣神感动，没有一个能说\InnerQuote{耶稣是主}的。} \BibleRef{格前 12:3} 所以，在这里，所表示的并非圣神的任何隶属，而是恩宠的联结。\par

74. 为使你们知道，区别并不在于一个音节，他在另一处也说：\BibleQuote{你们中间从前也有这样的人，但现在你们已被洗净，已被祝圣，已因主耶稣基督之名，并在我们天主的圣神内成了义人。} \BibleRef{格前 6:11} 这样的例子我能举出多少！因为经上记载：\BibleQuote{你们众人在基督耶稣内已成了一个。} \BibleRef{迦 3:28} 又记载：\BibleQuote{致书给在基督耶稣内受祝圣的人，} \BibleRef{格前 1:2} 又记载：\BibleQuote{好叫我们在他内成为天主的正义。} \BibleRef{格后 5:21} 另一处说：\BibleQuote{恐怕你们失去在基督耶稣内的贞洁。} \BibleRef{格后 11:3}\par

75. 然而，我这是在做什么？因为当我说论及圣子的同样也是论及圣神的，我实际上是在引导出这一点：并非因为论及圣子，所以仿佛论及圣神也是恭敬的，而是因为论及圣神的同样也被记载，因此人们就断言，圣子的尊荣因圣神而减损了。他们说，论及天主圣父的，是这样记载的吗？\par

76. 但让他们了解，论及天主圣父也这样说：\Quote{在上主内我赞美圣言；} 另一处说：\Quote{靠天主我们愿成大事；} 又说：\Quote{我的记念常在你内；} 又说：\Quote{我们要在你的名内欢欣；} 另一处又说：\BibleQuote{为显示他的行为是在天主内完成的；} \BibleRef{若 3:21} 保禄也说：\BibleQuote{在天主内，他创造了万有；} \BibleRef{弗 3:9} 又说：\BibleQuote{保禄和息耳瓦诺及弟茂德，致书给在天主父和主耶稣基督内的得撒洛尼教会；} \BibleRef{得后 1:2} 在福音中：\BibleQuote{我在父内，父在我内，} 和 \BibleQuote{住在我内的父。} \BibleRef{若 14:10} 经上也记载：\BibleQuote{凡自夸的，应在主内自夸；} \BibleRef{格后 10:17} 另一处说：\BibleQuote{我们与基督一同藏在天主内。} \BibleRef{哥 3:3} 他在这里说我们与基督在天主内，是比父更多地归于子吗？或者，我们的地位比圣神的恩宠更有效，以至于我们能同基督在一起，而圣神却不能？当基督愿意同我们在一起时，正如他自己说的：\BibleQuote{父啊！你所赐给我的人，我愿我在那里，他们也同我在一起，} \BibleRef{若 17:24} 难道他会不屑于与圣神同在吗？因为经上记载：\BibleQuote{你们聚集在一起，我的神魂也同你们一起，因我们的主耶稣的权能，} \BibleRef{格前 5:4}\par

77. 所以宗徒认为，使用哪一个\OriginalTerm{连接虚词}{conjunctive particle}并没有区别。因为每一个都是连接虚词，而连接并不造成分离，因为如果它分开，就不会被称为连接了。\par

78. 那么，是什么促使你们说，光荣、生命、伟大或权能，是\Emph{在}圣神内归于天主父或他的基督，却拒绝说\Emph{与}圣神同在呢？是否你们害怕似乎将圣神与父和子并列？但请听论及圣神的记载：\BibleQuote{因为在基督耶稣内的生命之神的法律，} \BibleRef{罗 8:2} 另一处天主父说：\Quote{他们要敬拜你，并在你内祈求。} 天主父说，我们应在基督内祈祷；而你们认为，如果基督的光荣说是在圣神内，对圣神有任何减损吗？\par

79. 请听，你们害怕承认圣神的事，宗徒却不怕为自己申明；因为他说：\BibleQuote{死，与基督同在一起，实在是更好。} \BibleRef{斐 1:23} 难道你们否认，宗徒藉以配得与基督同在的那位圣神，是与基督同在的吗？\par

80. 那么，是什么理由使你们宁愿说天主或基督是在圣神内受光荣，而非与圣神同受光荣呢？是否因为若说在圣神内，圣神就被宣称为小于基督？虽然你们使主或大或小的说法是可以驳斥的，但既然我们读到：\BibleQuote{因为他曾使那不认识罪的，替我们成了罪，好叫我们在他内成为天主的正义，} \BibleRef{格后 5:21} 那么在我们被视为最低微的那一位内，我们却找到了最高者。同样，在另一处你们读到：\BibleQuote{因为在他内，万有得以维系，} \BibleRef{哥 1:17} 即在他权能内。凡在他内得以维系的都不能与他相比，因为它们从他权能接受了赖以维系的实体。\par

81. 那么，你们是否理解，天主在圣神内为王，以至于圣神的权能，如同一种实体之源，将统治的起源分赐给天主？但这是不虔敬的。因此，我们的前辈论及父、子及圣神的\OriginalTerm{权能的合一}{unity of power}，他们说基督的光荣是与圣神同在，为的是表明他们不可分离的结合。\par

82. 因为，圣神怎能与圣子分离呢？既然\BibleQuote{圣神亲自和我们的心神一同作证：我们是天主的子女，如果是子女，便是承继者，是天主的承继者，是基督的同承继者。} \BibleRef{罗 8:16} 那么，谁是如此愚蠢，竟想要割断圣神与基督的永恒结合，而那位使我们成为基督同承继者的圣神，甚至将分离的事物连接起来。\par

83. \BibleQuote{只要我们，} 他说，\BibleQuote{与他一同受苦，也必与他一同受光荣。} \BibleRef{罗 8:16} 那么，如果我们藉着圣神将要与基督一同受光荣，我们怎能拒绝承认圣神自己与基督一同受光荣呢？当圣神说我们将与天主子一同生活时，我们是否就将基督与圣神的生命割裂？因为宗徒说：\BibleQuote{如果我们与基督同死，我们相信也要与他同生；} 然后又说：\BibleQuote{如果我们与他一同受苦，也必与他同生，而且不但同生，也要与他一同受光荣，不但受光荣，还要与他一同为王。} \BibleRef{弟后 2:11}\par

84. 因此，这些小品词决不意味着分裂，因为它们每一个都是连接小品词。最后，我们常在圣经中发现一个插入而另一个隐含的情况，如经上记载：\BibleQuote{我要带着全燔祭进入你的殿宇}，意即\BibleQuote{带着全燔祭}；另一处说：\BibleQuote{他领他们携带金银出来}，意即\BibleQuote{带着金银}。圣咏作者在别处又说：\BibleQuote{不是你，天主，遗弃了我们，不再与我们的军队一同出征吗？}，其实意为\BibleQuote{与我们的军队一同出征}。所以，在使用这种表达方式时，并不意味着尊荣的减损，也不应推演出任何有损天主性尊荣的结论；人必须心里相信，得以成义，口里宣认，得以得救。但那些心里不信的人，却用口散布贬损的言论。\par

% structure: p0104 source=h2 target=subsection reason=relative-heading-rank; base=h2

\subsection{第九章}

\Emph{圣保禄宗徒的一段话被异端者歪曲，用以证明天主三位之间的区分，在此予以解释，并证明整段话可正确地用来说明每一位，尽管它特别指向圣子。随后证明这段话的每一成分都适用于每一位，正如说}，\OriginalTerm{万物出于他}{of Him are all things} \Emph{可适用于圣父，同样}，\OriginalTerm{万物借着他并在他内}{all things are through Him and in Him} \Emph{也可用于论及圣父。}\par

85. 另一段类似的话，他们说它暗示着区分，即经上记载的：\BibleQuote{可是为我们只有一个天主，就是父，万物都出于他，我们也归于他；只有一个主，就是耶稣基督，万物借他而有，我们也借他而有。}\BibleRef{格前8:6} 因为他们强辩说，当说\Quote{出于他}时，指的是质料；当说\Quote{借着他}时，指的是工作的工具或某种职能；但当说\Quote{在他内}时，则指的是受造万物得以显现的场所或时间。\par

86. 那么，他们力图证明在实体上存在某种差异，急切地在所谓的工具与真正的工匠或作者之间作出区分，也在时间或场所与工具之间作出区分。然而，难道圣子在性体上就与圣父相异，只因工具不同于工匠或作者吗？或者，圣子与圣神相异，只因时间或场所与工具不属于同类吗？\par

87. 现在对照一下我们的主张。他们认为质料出于天主，仿佛出于天主的本性，就像你们说一个箱子是由木头制成的，或一尊雕像是用石头制成的；质料就是以这种方式出于天主，而同一质料又由圣子如同某种工具所制造；因此他们声称，圣子与其说是工匠，不如说是工作的工具；并且万物都是在圣神内受造的，如同在某个场所或时间内；他们把每一部分分别归于每一位，并否认所有属性是共有的。\par

88. 但我们表明，万物都是如此地出于天主父，以至于父天主并无任何损失，因为万物或是借着他、或是在他内，而并非像出自质料那样出于他；同时，万物是借着主、圣子而有的，因此他不缺乏万物出于圣子并在圣子内的属性；而且，万物都在圣神内，因此我们也可以教导说：万物借着圣神而有，万物出于圣神。\par

89. 因为这些小品词，如我们之前所说的那些，是相互蕴含的。因为宗徒并非说\Quote{万物出于天主，万物借着圣子}，以表示父与子的实体可以割裂，而是要教导，通过一种没有混淆的区分，父是一位，子又是另一位。因此，这些小品词并非彼此对立，而是似乎彼此联合、和谐一致，以致常常适用于同一位，如经上记载：\BibleQuote{因为万物都出于他，借着他，并在他内。}\BibleRef{罗11:36}\par

90. 但如果你真正考虑这段话的出处，你便不会怀疑它是指圣子而言。因为宗徒按照依撒意亚先知的预言说：\BibleQuote{谁曾测度上主的心灵，或作过他的参谋？}\BibleRef{依40:13} 然后他又加上：\BibleQuote{因为万物都出于他，并在他内。} 而依撒意亚论及万物的工匠时，正如你所读到的：\BibleQuote{谁曾用手掌量过海水，用手掌测过高天，用升斗量过地上的尘土，用秤称过山岳，用天秤量过丘陵呢？谁曾测度上主的心灵，或作过他的参谋？}\BibleRef{依40:12-13}\par

91. 宗徒又加上说：\BibleQuote{因为万物都出于他，借着他，并在他内。} \Quote{出于他}是什么意思？即是说，万物的本性出于他的旨意，他是一切存在物的创造者。\Quote{借着他}是什么意思？即是说，万物的建立与持续是他的恩赐。\Quote{在他内}是什么意思？即是说，万物以一种奇妙的渴慕和不可言喻的爱仰望着自己生命的创造者，以及恩宠与功能的赐予者，正如经上所载：\BibleQuote{众生的眼睛都仰望你，} 以及 \BibleQuote{你张开手，使一切生物满享你的恩惠。}\par

92. 对于圣父，你也可以正确地说\Quote{出于他}，因为运作的智慧是出于他，这智慧凭着他自己与父的旨意，使一切不存在的得以存在。\Quote{借着他}，因为万物是借着他的智慧而受造的。\Quote{在他内}，因为他是实体的生命之泉源，我们在他内生活、行动、存在。\par

93. 对于圣神，因为我们由他所形成，由他所坚固，在他内得以建立，我们领受了永生的恩赐。\par

94. 因此，既然这些表述似乎适用于圣父、圣子或圣神，那么可以肯定，其中并无贬抑之处，因为我们既说许多事物属于圣子，又说许多事物藉着圣父，正如你发现论及圣子所说：\BibleQuote{我们应在各方面长进，而归于元首基督；整个身体都因祂而结构紧凑，借着各关节的供应，依照各肢体的功用，各尽其职，使身体增长，在爱德中建立自己。} \BibleRef{弗 4:15} 再者，他写信给哥罗森人，论到那些不认识天主圣子的人，说：\BibleQuote{因为他们不持守元首，全身由关节和筋络供给并连接，就因天主而增长。} \BibleRef{哥 2:19} 因为我们在前面已说过，基督是教会的头。在另一处你又读到：\BibleQuote{从祂的满盈中，我们都领受了。} 主自己也曾说：\BibleQuote{祂要由我领受，并传告给你们。} \BibleRef{若 16:14} 在这之前，祂又说：\BibleQuote{我觉得有能力从我身上出去了。} \BibleRef{路 8:46}\par

95. 同样，为使你认识合一性，论及圣神也说：\BibleQuote{那随从圣神撒种的，必由圣神收获永生。} \BibleRef{迦 6:8} 若望说：\BibleQuote{我们所以知道天主住在我们内，是由于祂赐给我们的圣神。} \BibleRef{若一 4:13} 天使说：\BibleQuote{那在她内受生的，是出于圣神。} \BibleRef{玛 1:20} 主说：\BibleQuote{由神生的属于神。} \BibleRef{若 3:6}\par

96. 因此，正如我们读到万物出于圣父，同样，也可以说万物出于圣子，因为万物是藉着祂；我们由证据得知，万物也出于圣神，因为万物都在祂内。\par

97. 现在让我们考虑，我们能否教导有什么事是藉着圣父的。圣经记载：\BibleQuote{保禄，基督的仆役，因天主的旨意；} \BibleRef{格前 1:1} 又在别处说：\BibleQuote{从此你不再是奴仆，而是儿子；既是儿子，也藉着天主成了承继人。} \BibleRef{迦 4:7} 又说：\BibleQuote{就如基督因父的光荣从死者中复活。} \BibleRef{罗 6:4} 别处，天主圣父向圣子说：\BibleQuote{看，外邦人将藉着我归向祢。}\par

98. 如果你寻找藉着圣父所行的事，你还会找到许多其他章节。那么，难道因为我们在圣经中读到许多事在子内及属于子，又发现许多事藉着圣父而行或赐予，圣父就较低吗？\par

99. 但同样，我们也读到许多事藉着圣神而行，如你所见：\BibleQuote{天主却藉着圣神将这一切启示了我们；} \BibleRef{格前 2:10} 在另一处：\BibleQuote{要借着圣神守护所受的美好寄托；} \BibleRef{弟前 6:20} 致厄弗所人书中：\BibleQuote{藉着祂的圣神，使你们内在的人得到坚固；} \BibleRef{弗 3:16} 致格林多人书中：\BibleQuote{有人由圣神蒙赐智慧的言语；} \BibleRef{格前 12:8} 又在别处：\BibleQuote{如果你们依赖圣神，致死肉性的妄动，必能生活；} \BibleRef{罗 8:13} 还有上文：\BibleQuote{那使基督从死者中复活的，也必要藉那住在你们内的圣神，使你们有死的身体复活。} \BibleRef{罗 8:11}\par

100. 但或许有人会说，请指出我们能明确地读到万物出于圣子，或万物出于圣神。但我答复说，让他们也指出圣经记载万物藉着圣父。既然我们已经证明这些表述适用于圣父、圣子或圣神，而且这类介词不会导致天主神力的任何区别，那么毫无疑问，万物出于那藉其而有万物的那位；万物也藉着那藉其而有一切的那位；我们必须理解，万物或藉着或由于那位万有都在其内的那一位。因为一切受造物，都是因\OriginalTerm{天主圣三}{Trinity}的旨意、透过其工程、在其能力内而存在，正如经上记载：\BibleQuote{让我们照我们的肖像，按我们的模样造人；} \BibleRef{创 1:26} 又在别处说：\BibleQuote{因主的一句话，诸天造成；因祂口中的气息，万象生成。}\par

% structure: p0122 source=h2 target=subsection reason=relative-heading-rank; base=h2

\subsection{第十章}

为证明圣三的旨意、召唤和命令是同一的，圣盎博罗削指出，圣神召唤教会的方式，与圣父和圣子完全相同，并借着拣选圣保禄和圣巴尔纳伯，尤其是借着派遣圣伯多禄往科尔乃略那里去，证明了这一点。在叙述中，他指明了宗徒的神视如何预示了外邦人的蒙召：他们从前如同野兽，如今因圣神的运作，却脱去了兽性。然后，他引用了其他支持这观点的章节，并指明耶肋米亚被犹太人投入坑中，又为阿贝得默肋客所救援，这预表了犹太人对圣神的轻慢，以及外邦人对祂的尊崇。\par

101. 不但圣父、圣子和圣神的工程在处处是同一的，而且他们的旨意、召唤和颁布命令也是同一的，这在教会伟大而救恩的奥迹中可以见到。因为正如圣父召唤外邦人进入教会，说：\BibleQuote{我要叫她\InnerQuote{我的民族}，她本不是我的民族；我要叫她\InnerQuote{蒙爱者}，她本未蒙爱；} \BibleRef{欧 2:23} 又在别处说：\BibleQuote{我的殿宇将称为万民的祈祷所，} \BibleRef{依 56:7} 同样，主耶稣也说，保禄蒙祂拣选，为召唤并聚集教会，正如你发现主耶稣对阿纳尼雅所说的：\BibleQuote{你去！因为这人是我所拣选的器皿，为把我的名字带到外邦人前。} \BibleRef{宗 9:15}\par

102. 正如天主圣父召叫了教会，基督也召叫了它，圣神也召叫了它，说：\BibleQuote{你们给我选拔出\Person{保禄}和\Person{巴尔纳伯}，去作我召叫他们担任的工作。} 於是，接着又记载：\BibleQuote{他们遂禁食祈祷，给他们覆了手，派他们走了。他们既被圣神派遣，就下到\Place{色娄基雅}。} 这样，\Person{保禄}不仅因基督的旨意，也因圣神的旨意领受了宗徒职务，并急忙去聚集\OriginalTerm{外邦人}{Gentiles}。\par

103. 不但\Person{保禄}，而且我们在\WorkTitle{宗徒大事录}中读到的\Person{伯多禄}也是如此。因为当他在祈祷中看见天开了，降下一个四角系着的器皿，好像一块布，里面有各种四足走兽、野兽和天空的飞鸟，\BibleQuote{有声音向他说：\Person{伯多禄}，起来，宰了吃！\Person{伯多禄}却说：主，绝不可！我从来没有吃过污秽和不洁之物。又有声音向他说：天主所洁净的，你不可称为污秽。这样一连三次，那器皿就被收回天上去了。} 当\Person{伯多禄}心里正在思量这异象时，\Person{科爾乃略}的仆人，由\OriginalTerm{天使}{Angel}指定，来到他那里，圣神对他说：\BibleQuote{看，有人找你，起来，下去同他们去罢！不要疑惑，因为是我派遣了他们。} \BibleRef{宗 10:19}\par

104. 圣神多么清楚地彰显了祂自己的能力！首先，祂启发了正在祈祷的人，并与恳求的人同在；接着，\Person{伯多禄}被召叫时，回答说：\Quote{主，}因此他配得上第二次的讯息，因为他认出了主。圣经却明确指出那位主是谁，因为他回答的那一位，在他回答时对他说话。接下来的话显示了圣神被清楚启示，因为那构成\OriginalTerm{奥迹}{mystery}的，将那奥迹启示出来。\par

105. 也要注意，这奥迹的显现重复了三次，表达了\OriginalTerm{天主圣三}{Trinity}的运作。因此，在奥迹中，提出了三重的询问，也作出三重的回答，除非经过三重的宣认，没有人能得到洁净。为此，\Person{伯多禄}在福音中三次被问及是否爱主，藉着这三重的回答，他因否认主而承担的罪债得以解开。\par

106. 还有，因为\OriginalTerm{天使}{Angel}被派到\Person{科爾乃略}那里，圣神对\Person{伯多禄}说：\BibleQuote{主的眼目眷顾地上的信友。} 这并不是没有目的的，当祂先前说过：\BibleQuote{天主所洁净的，你不可称为污秽，} \BibleRef{宗 10:15} 圣神降临到\OriginalTerm{外邦人}{Gentiles}身上，为净化他们，这表明圣神的行动是天主的行动。但是，\Person{伯多禄}受圣神派遣时，并没有等候天主圣父的命令，而是承认那讯息是出于圣神自己，那\OriginalTerm{恩宠}{grace}也是出于圣神自己，当他说：\BibleQuote{那么，既然天主赐给了他们同样的恩宠，如同赐给我们这些信主耶稣基督的人一样，我是谁，竟能阻挡天主呢？}\par

107. 所以，是圣神把我们从\OriginalTerm{外邦人}{Gentiles}的不洁中解救了出来。因为在那些四足走兽、野兽和飞鸟之中，含有人的状况的预像，这状况除非因圣神的\OriginalTerm{圣化}{sanctification}而变得温和，否则便披着野兽的凶残。因此，那改变兽类的狂暴为圣神的纯朴的\OriginalTerm{恩宠}{grace}，是何等美妙：\BibleQuote{因为我们从前也是昏愚的，悖逆的，迷途的，受各种贪欲和逸乐的奴役；但现在，藉着圣神的更新，我们开始成为基督的继承人，并成为与天使同承继业者。} \BibleRef{铎 3:3}\par

108. 所以，圣先知\Person{达味}在圣神内预见我们应从野兽变成相似天上的居民，便说：\BibleQuote{求你叱责林中的野兽，} 这显然不是指那因野兽奔跑而震动、因动物的吼叫而摇晃的树林，而是指经上记载的树林：\BibleQuote{我们在树林的田野中找到了它。} 在这树林中，正如先知所说的：\BibleQuote{义人将兴盛如棕榈树，繁茂如\Place{黎巴嫩}的香柏。} 这树林在先知所预言的树梢上摇动，倾泻出天上\OriginalTerm{圣言}{Word}的滋养。\Person{保禄}进入这树林时，确实像一只残暴的狼，但出来时却成了牧者，因为\BibleQuote{它们的声音传遍了整个大地。}\par

109. 我们曾经是野兽，因此主说：\BibleQuote{你们要提防假先知，他们外面披着羊皮，里面却是残暴的狼。} \BibleRef{玛 7:15} 但现在，藉着圣神，狮子的狂暴、豹子的斑点、狐狸的狡猾、狼的贪婪，都已从我们的感受中消失；因此，那使大地变成了天堂的\OriginalTerm{恩宠}{grace}是何等伟大，我们这些从前像野兽在林中游荡的人，如今却在天上生活。 \BibleRef{斐 3:20}\par

110. 不仅在此处，而且在同书的其他地方，宗徒\Person{伯多禄}也宣称教会是由圣神建立的。因为你读到他说：\BibleQuote{洞察人心的天主，为他们作证，赐给了他们圣神，如同赐给我们一样；在我们和他们之间，祂没有作任何区别，因着\OriginalTerm{信德}{faith}净化了他们的心。} \BibleRef{宗 15:8}\par

111. 因此，你不要像犹太人那样，轻视先知们所预言的圣子；因为那样你就也轻视了圣神，轻视了\Person{依撒意亚}，轻视了\Person{耶肋米亚}，就是那位蒙主拣选的人，用破布和绳子从犹太人寓所的深坑中将他拉上来的那位。 \BibleRef{耶 38:11} 因为犹太人民轻视了先知的预言，把他抛入深坑。没有一个犹太人出来将先知拉出，只有一个\Place{厄提约丕雅}人\Person{厄贝得默客肋}，正如圣经所证实的。\par

在该叙述中有一个极美的预像，就是说，我们这些外邦人中的罪人，因过犯而预先成为黑暗，先前不结果实，从深处举起了犹太人因心思和肉欲，如同投入泥沼中一般所抛弃的先知之言。因此经上记载：\Quote{埃塞俄比亚将向天主伸出她的手。}这预表着圣教会的形象，她在\WorkTitle{雅歌}中说：\BibleQuote{耶路撒冷的众女子啊，我虽黑，却是秀美}\BibleRef{歌 1:5}，因罪而黑，因恩宠而秀美；因本性状况而黑，因救赎而秀美，或确切地说，因劳苦的尘土而黑。所以她在战斗中时是黑的，在戴上胜利的冠冕时是秀美的。\par

先知被用绳索提起是合宜的，因为忠信的作者说：\Quote{绳尺落在我的佳美之处。}用破布也是合宜的；因为主自己，当那些首先被邀请赴婚宴的人推辞时，就打发人到路口去，凡遇到的，无论好人坏人，都请来赴婚宴。于是，随着这些破布，祂从泥沼中举起了先知之言。\par

% structure: p0137 source=h2 target=subsection reason=relative-heading-rank; base=h2

\subsection{第十一章}

如果我们相信子与圣神知道万事，我们就效法\Person{厄贝得默肋客}的榜样。圣经将这知识归于圣神，也归于子。子因圣神受光荣，圣神也因子受光荣。此外，当我们读到父、子和圣神所说和启示的都相同，我们就必须承认在祂们内有着本性和知识的合一。最后，圣神洞察天主的深奥事理，这并不是无知的标记，因为父与子同样被说成洞察，而\Person{保禄}虽被基督拣选，却仍由圣神教导。\par

而你，也将成为\Person{厄贝得默肋客}，即被主所拣选者，如果你从外邦人无知的深处举起天主的\OriginalTerm{圣言}{Word of God}；如果你相信天主子不受欺骗，没有什么能逃过祂的认知，祂对将要发生的事并不无知。圣神也不受欺骗，论到祂，主说：\Quote{当那\OriginalTerm{真理之神}{Spirit of Truth}来到时，祂要引领你们进入一切真理。}\BibleRef{若 16:13}说\Emph{一切}的祂，丝毫不遗漏，无论是时日，无论过去的事或未来的事。\par

为使你知道祂既知道万事，又预告未来的事，且祂的知识与父和子的知识是一体的，请听天主的真理论到祂所说的话：\Quote{因为祂不是凭自己说话，而是祂所听见的，祂要说出来，并将未来的事告诉你们。}\BibleRef{若 16:13}\par

所以，为叫你看出祂知道万事，当子说：\Quote{至于那日子和时辰，没有人知道，连天上的天使也不知道}\BibleRef{谷 13:32}时，祂将圣神除外了。但若圣神被免于无知，为何天主子不被除外呢？\par

但你说祂也将天主子列于天使之中。祂确实将子列入，却没有将圣神列入。那么，请你承认，要么圣神大于天主子，这样你现在不仅是像\OriginalTerm{亚略派}{Arian}，甚至是像\OriginalTerm{福提尼安派}{Photinian}在说话；要么就承认你应当把那句话理解为：祂说子不知道，是因为作为人的祂（就人性而言）可以被列于受造物之中。\par

但如果你愿意相信天主子知道万事，并预知一切，就看看那些你以为子不知道的事，圣神却从子领受了。然而，祂是借着\OriginalTerm{本体的合一}{Unity of Substance}领受，就如子从父领受一样。祂说：\Quote{祂要光荣我，因为祂要把由我所领受的传告给你们。凡父所有的一切，都是我的；为此我说：祂要把由我领受的传告给你们。}\BibleRef{若 16:14}那么，还有什么比这合一更清楚的呢？凡父所有的都属于子；凡子所有的，圣神也领受了。\par

但你还要知道，子知道审判的日子。我们在\WorkTitle{匝加利亚}书中读到：\Quote{上主我的天主必将来临，众圣者与祂同来。在那一天，没有光，只有寒冷和冰霜，那将是一个独特的日子，为上主所知。}因此，这日子为上主所知，祂将与祂的圣者们一同来临，以祂的\OriginalTerm{第二次来临}{second Advent}光照我们。\par

但让我们继续前面关于圣神的话题。因为在我们引用的经文中，你发现子论圣神说：\Quote{祂要光荣我。}这样，圣神光荣子，就如父也光荣祂，而天主子也光荣圣神，如我们以上所说。因此，谁借着\OriginalTerm{永恒之光的合一}{Unity of the Eternal Light}成为彼此光荣的原因，谁就不是软弱的，祂也不低于圣神——祂被圣神光荣是确实的。\par

而你也要被拣选，如果你相信圣神所说的，就是父所说的，也是子所说的。的确，\Person{保禄}之所以被拣选，正是因为他如此相信并如此教导，因为经上记载，天主\Quote{借着圣神，将眼所未见，耳所未闻，人心所未想到的，天主为爱祂的人所预备的，都启示给我们了。}\BibleRef{格前 2:9}因此，祂被称为\OriginalTerm{启示之神}{Spirit of revelation}，正如你读到的：\Quote{因为天主将\OriginalTerm{智慧和启示之神}{Spirit of wisdom and revelation}赐给那些如此预备自己的人，好使人认识祂。}\BibleRef{依 64:4}\par

因此，存在着\OriginalTerm{知识的合一}{Unity of knowledge}，因为如同父赐下启示之神而启示，子也启示，如经上所载：\Quote{除了父，没有人认识子；除了子和子所愿意启示的人，也没有人认识父。}\BibleRef{玛 11:27}经上论到子说得更多，不是因为祂比父有更多，而是为避免人以为祂比父少。父这样被子启示，也并无不当，因为子认识父，正如同父认识子一样。\par

124. 现在要知道，圣神也认识天主圣父，因为经上记载说：\BibleQuote{除了在人内的灵，没有人知道人的事；同样，除了天主的圣神，也没有人知道天主的事。}\BibleRef{格前 2:11} 他说：\Quote{没有人，}\Quote{知道，除了天主的圣神。} 那么，天主子被排除了吗？当然不，因为当说：\Quote{除了子，没有人认识父。}的时候，圣神也没有被排除。\par

125. 因此，圣父、圣子和圣神具有同一性体和同一知识。圣神不应被列入由圣子所造的一切事物之中，因为祂认识圣父，而（如经上所载）除了圣子，谁能认识圣父呢？但圣神也认识。那么，这意味着什么？当论及万物的总和时，结论便是圣神不包含在内。\par

126. 现在，我希望他们回答：在人内认识人的事的是什么。那必定是超出灵魂其他能力、并用以衡量人之最高性体的理性能力。那么，圣神是谁？祂认识天主的深奥事理，全能的天主也藉祂而启示。祂在天主性的圆满上逊于圣父吗？而这一实例甚至证明祂与圣父是同一实体。或者，认识天主的筹谋及自始隐藏的奥迹的那一位，还会有任何不知之事吗？那认识所有属于天主之事者，还会有什么不知道的呢？因为\BibleQuote{圣神甚至洞察天主的深奥事理。}\BibleRef{格前 2:10}\par

127. 为使你们不致认为圣神洞察未知的事，既洞察而后方知所不知者，这里首先说明：天主藉祂的圣神将这些事启示给我们了；同时，为使你们明白，那藉圣神亲自启示给我们的事，圣神原是知道的，随后便说：\BibleQuote{因为除了在人内的灵，谁晓得人的事？同样，除了天主的圣神，谁也不知道天主的事。}\BibleRef{格前 2:11} 那么，如果人的灵知道人的事，且在洞察之前就已知道，天主的事中难道还有什么是天主的圣神所不知的吗？宗徒论及祂时并非无意地说：\Quote{除了天主的圣神，谁也不知道天主的事；} 这并非说祂是藉着洞察而知道，而是由于本性而知道；对天主事的认知，在祂并非是偶性的，而是祂本性的知识。\par

128. 若你们因祂说过\Quote{洞察}而困惑，便应知道，论及天主时也如此说，因为祂是洞察人心和肺腑的天主。祂亲口说过：\BibleQuote{我是洞察人心和肺腑的那一位。}\BibleRef{耶 17:10} 论及天主子，你们在\BibleBook{}{希伯来书}中也读到：\BibleQuote{祂是洞察心思和意念的。}\BibleRef{希 4:12} 由此可见，低位的不会洞察其高位的内里，因为能知道隐秘事只属于天主的权能。所以，圣神作为洞察者，与圣父一样；圣子也是洞察者，一样如此，这说法正隐含了：那无从遁形者，显然没有祂不知道的事。\par

129. 最后，他由基督所拣选，并由圣神所教导。正如他自己所见证的，他藉圣神获得了对神圣奥迹的知识，既显示圣神认识天主，也显示圣神将天主的事启示给了我们，正如圣子也曾启示了一样。他接着又说：\BibleQuote{我们所领受的，不是这世界的精神，而是那出于天主的圣神，为使我们能明了天主所赐予我们的一切；我们讲论这些事，也不是用人智慧之说服的言辞，而是用圣神的彰显和天主的德能。}\BibleRef{格前 2:12}\par

% structure: p0154 source=h2 target=subsection reason=relative-heading-rank; base=h2

\subsection{第十二章。}

在证明圣神与圣父、圣子同样是启示的赐予者之后，现在要解释这同一圣神如何不从自己发言；并指明在祂内不可设想任何身体器官，也不可从我们读到祂听见这一事实推想出任何低劣性，因为同样的说法也须归于圣子，甚至归于圣父，因为圣父也听从圣子。于是，圣神听见并光荣圣子，其意义是祂将圣子启示给先知们和宗徒们，由此推断出三位的工作行动是合一的；并且，既然圣神与圣父行同样的事工，那么每一位的实体也被宣认为同一。\par

130. 至此已经证明，正如天主将祂的事启示给我们，圣子也如此，圣神也如此，将天主的事启示了。因为我们的知识是从一位圣神，藉一位圣子，而归于一位圣父；并且从一位圣父，藉一位圣子，而归于一位圣神，善与圣化，及永恒德能的至高权柄被通传出来。因此，哪里圣神彰显，那里就有天主的德能，也不能有任何区分，因为工作是同一的。所以，圣子所说的，圣父也说了；圣父所说的，圣子也说了；圣父和圣子所说的，圣神也说了。\par

131. 因此，天主子论及圣神也说过：\BibleQuote{祂不凭自己发言，}\BibleRef{若 16:13} 意思是，并非没有圣父和我自己的参与。因为圣神不是分裂且分离的，而是说出祂所听到的。祂听见，就是说，借着实体的合一和知识的秉性。因为祂不是借着身体上的任何孔窍来听取，也不是以任何血肉的尺度来发出神圣的声音，祂也不会听到祂不知道的事；因为在人的事务中，通常听觉产生知识，然而，即使在人的本身内，也不常有肉身的言语或肉欲的听闻。因为经上说：\BibleQuote{那说语言的，不是对人说，而是对天主说，因为没有人听得懂，他是在圣神内讲说奥秘。}\BibleRef{格前 14:2}\par

132. 因此，如果人的听觉并非总是属于肉体，那么当天主被说成听见，是为了使我们相信祂知晓，难道你要求天主有人的软弱之声，以及某种属肉的听觉器官吗？因为我们听见了，我们知道；我们预先听见，以便能知晓；但在知晓万有的天主内，知晓先于听见。所以，为了说明圣子并非不知圣父的意愿，我们说祂听见了；但在天主内，没有声音或音节，这类东西通常表示意愿的指示；而是意志的合一在天主内以隐秘的方式被领悟，在我们内却藉着标记显示出来。\par

133. 那么，\Quote{祂不凭自己说话}是什么意思？意思是，祂不会没有我而说话；因为祂讲真理，祂呼出智慧。祂不凭自己说话，因为祂是天主的圣神；祂不凭自己听见，因为万有都出于天主。\par

134. 圣子从圣父领受了一切，因为祂亲自说过：\Quote{一切都由我父交给了我。}\BibleRef{玛 11:27} 凡父所有，子也有，因为祂又说：\Quote{凡父所有的一切，都是我的。}\BibleRef{若 16:15}\EditorialNote{此处经文在原稿中误标为若15:15，据上下文更正为若16:15。} 而祂自己因本性的合一所领受的，圣神也因本性的同一合一从祂领受了，正如主耶稣亲自论及祂的圣神时所声明的：\Quote{所以我说过，祂要领受我的，而传告给你们。}\BibleRef{若 16:15} 因此，圣神所说的，是圣子的；圣子所赐予的，是圣父的。这样，圣子和圣神都不凭自己说什么。因为天主圣三不说任何在自身之外的事。\par

135. 但是，如果你坚持认为这是论证圣神的软弱，并且与肉身的卑微有某种相似，那么你也会以此论证来损害圣子，因为圣子论自己说：\Quote{我怎样听见，就怎样审判，}\BibleRef{若 5:30} 以及 \Quote{子由自己什么也不能作，除非他看见父所作的。}\BibleRef{若 5:19} 因为，正如圣子所说的，那是真实的：\Quote{凡父所有的一切，都是我的，}\BibleRef{若 16:15} 而圣子按\OriginalTerm{天主性}{Godhead}与圣父为一，因自然的本体而为一，并非依据\OriginalTerm{萨培利派}{Sabellian}的虚妄；凡因本体的\OriginalTerm{特性}{property}而为一的，当然不能分开，所以圣子除非听见圣父的，什么也不能作，因为天主的圣言永远长存，圣父也从不与圣子的行动分离；圣子所作的，祂知道是圣父所愿意的，圣父所愿意的，圣子知道怎样去作。\par

136. 最后，为了使人不认为圣父与圣子之间在工作上有时间或次序上的差别，而相信同一行动的合一性，祂说：\Quote{我所做的，祂也做。} 再者，为了使人不认为工作的区分上有任何差异，而判断圣父与圣子的意愿、工作和能力是相同的，智慧曾论及圣父说：\Quote{因为凡祂所做的，子也照样做。}\BibleRef{若 5:19} 这样，无论哪一位的行动不在另一位之先或之后，而是同一运作的同样结果。为此，圣子说自己由自己什么也不能作，因为祂的行动不能与圣父的行动分离。同样，圣神的行动也不被分离。因此，祂所说的一切，也被说是从圣父听见的。\par

137. 如果我证明圣父也听见圣子，正如圣子也听见圣父，那又如何？因为你在福音中记载着圣子说：\Quote{父啊！我感谢你，因为你俯听了我。}\BibleRef{若 11:41} 圣父怎样听见了圣子？因为在先前关于拉匝禄的记载中，圣子并未向圣父说什么。为了使我们不认为圣子只被圣父听见一次，祂又加上说：\Quote{我早已知道你常俯听我。}\BibleRef{若 11:42} 所以，这听觉并非隶属的服从，而是永恒的合一。\par

138. 同样地，圣神被说成从圣父听见，并光荣圣子。光荣，因为圣神教导我们，圣子是那不可见的天主的肖像，\BibleRef{哥 1:15} 是祂光荣的光辉，是祂本体的印象。\BibleRef{希 1:3} 圣神也曾在圣祖和先知们内发言，最后，宗徒们在领受圣神之后，便开始成为更成全的。因此，没有天主性能力和恩宠的分离，因为虽然有\Quote{神恩虽有区别，却是同一的圣神；职分虽有区别，却是同一的主；功效虽有区别，却是同一的天主，在一切人身上行一切事。} 职务虽有区别，但不是天主圣三的分割。\par

139. 最后，那在一切人身上行一切事的，是同一的天主，为使你知道天主父和圣神之间的运作没有区别；因为圣神所做的，天主父也做，即那\Quote{在一切人身上行一切事}的。因为当天主父在一切人身上行一切事时，却\Quote{有人藉圣神蒙赐智慧的言语；另一人却由同一圣神蒙赐知识的言语；同一圣神赐人信心；同一圣神赐人治病的奇恩；有人能行奇迹；有人能说先知话；有人能辨别神类；有人能说各种语言；有人能解释语言：可是，这一切都是这同一的圣神所作的，随祂的心愿，个别分配与人。}\par

140. 所以毫无疑问，圣父所做的，圣神也做。祂行事不依据命令，不像人以肉体的方式听从，而是自愿地，因为祂在自己的意志上是自由的，不是别人能力的仆役。因为祂不服从，如奉命行事，而是作为赐予者，祂是自己恩赐的主宰。\par

141. 此外，试想你能否说圣神成就圣父所成就的一切事；因为你不能否认圣父成就圣神所成就的事；否则，若圣父不成就圣神也成就的事，圣父就不是成就万事。但如果圣父也成就圣神所成就的事，既然圣神是按自己的旨意分施祂的运作，你就必须说，要么圣神所分施的是按自己的旨意，违背天主圣父的旨意；要么你若说圣父与圣神的旨意相同，你就必须承认神性意志与运作的合一，即使你不情愿，若不从心里，至少从口里。\par

142. 然而，若圣神与天主圣父有同一意志与运作，祂也就有同一性体，因为造物主由祂的工程而被认识。因此，\BibleQuote{圣神是同一的，主是同一的，天主是同一的}\BibleRef{格前 12:5}。你称圣神，祂仍是同一的；你称主，祂仍是同一的；你称天主，祂仍是同一的。并非同一到祂自己同时是圣父、圣子、圣神[同一而独有的位格]；而是因为圣父与圣子有同一的德能。所以，祂在性体与德能上都是同一的，因为在天主性内既没有\Person{撒伯流}的混同，也没有\Person{亚略}的分割，也没有任何尘世与肉身的变迁。\par

% structure: p0169 source=h2 target=subsection reason=relative-heading-rank; base=h2

\subsection{第十三章}

预言不仅来自圣父和圣子，也来自圣神；圣神在宗徒们身上的权柄与运作，与圣父和圣子的权柄与运作同等；因此我们应当明白，在权柄、治理与恩惠这三点上存在着合一；然而，这种分享并不需担心有什么不利，正如在人类的友谊中也不会有这样的不利。最后，宗徒们被描述为服从圣神，由此事实可确证，这乃是宗徒信仰的传承。\par

143. 可敬的皇帝，请接受这个问题上的另一个有力例证，也是你所知道的：\BibleQuote{天主在古时，曾多次并以多种方式，藉着先知对我们的祖先说过话。}\BibleRef{希 1:1} 天主的智慧曾说：\BibleQuote{我要派遣先知和宗徒。}\BibleRef{路 11:49} 经上也记着说：\BibleQuote{这人从圣神蒙受了智慧的言语；另一人却由同一圣神蒙受了知识的言语；有人在同一圣神内蒙受了信心；另有人在同一圣神内却蒙受了治病的奇恩；有的能行奇迹；有的能说预言。} 因此，按宗徒的话，预言不仅通过圣父和圣子，也通过圣神；所以，职分是同一的，恩宠也是同一的。由此你便见到，圣神也是预言的创始者。\par

144. 宗徒们也曾说：\BibleQuote{圣神和我们决定，不再加给你们什么重担。}\BibleRef{宗 15:28} 他们说\Emph{决定}，不仅指这恩宠的施行者，也指这命令得以执行的创始者。因为正如我们读到论天主的话：\BibleQuote{这竟使天主喜悦，} 同样，当说\BibleQuote{圣神决定}时，也描绘出一位自己权柄的主宰。\par

145. 说话按自己的旨意，命令按自己的旨意，如同圣父命令和圣子命令一样，祂怎么会不是主宰呢？因为正如\Person{保禄}听见有声音向他说：\BibleQuote{我就是你所迫害的耶稣，}\BibleRef{宗 9:5} 同样，圣神也禁止保禄和\Person{息拉}到\Place{彼提尼雅}去。又正如圣父藉先知们说话，同样，\Person{阿加波}论圣神说：\BibleQuote{圣神这样说：犹太人要在\Place{耶路撒冷}这样捆绑这条腰带的主人。}\BibleRef{宗 21:11} 正如智慧派遣宗徒们，说：\BibleQuote{你们往普天下去，向一切受造物宣传福音，}\BibleRef{谷 16:15} 同样，圣神也说：\BibleQuote{你们给我分出\Person{巴尔纳伯}和\Person{扫禄}，去做我叫他们做的工作。}\BibleRef{宗 13:2} 于是他们为圣神所派遣，正如圣经后来所指明的，他们与其他宗徒没有分别，仿佛他们是由天主圣父以某种方式派遣的，由圣神以另一种方式派遣的。\par

146. 最后，\Person{保禄}被圣神所派遣，既是基督所拣选的器皿，他自己也述说天主在他身上工作，说：\BibleQuote{因为那在\Person{伯多禄}身上工作，使他为受割损者尽宗徒之职的，也在我身上工作，使我为外邦人效力。}\BibleRef{迦 2:8} 既然在保禄身上工作的，与在伯多禄身上工作的是同一者，那么很明显，既然圣神在保禄身上工作，圣神也在伯多禄身上工作。但伯多禄自己作证，圣父在他身上工作，正如\WorkTitle{宗徒大事录}所记载的，伯多禄站起来对他们说：\BibleQuote{诸位仁人弟兄！你们都知道，多时以前，天主就在你们中间拣选了，由我的口叫外邦人听到福音的宣讲。} 所以，在伯多禄身上，天主工作了宣讲的恩宠。谁又敢否认基督在他身上的工作呢？因为他确实是由基督所拣选所指定的，当主说：\BibleQuote{你喂养我的羔羊。}\BibleRef{若 21:15}\par

147. 因此，圣父、圣子和圣神的运作是同一的，除非你也许否认这同一的运作在宗徒身上，而认为：圣父和圣神在伯多禄身上工作，而圣子曾在他身上工作，好像圣子的运作单靠他自己不足以使伯多禄获得这恩宠似的。这样，圣父、圣子和圣神的能力仿佛被联合起来，以使这工作变得多重化，免得基督单独的运作过于软弱，不足以坚固伯多禄。\par

148. 不仅在伯多禄身上可发现圣父、圣子和圣神的同一行动，而且在所有宗徒身上也显示出天主行动的合一性，以及对天上恩赐分施的某种权柄。因为天主的行动是凭借命令的德能而行事，而非通过行使职务；因为天主行动时，并不用劳苦或技艺来制造什么，而是\BibleQuote{祂一发言，万有造成。}祂说：\BibleQuote{要有光，就有了光，}\BibleRef{创 1:3}因为工程的成就包含在天主的命令之中。\par

149. 我们若愿意思考，就能轻易发现，圣经的见证将这种王权归于圣神；并且清楚显明所有宗徒不仅是基督的门徒，也是圣父、圣子和圣神的仆役。正如外邦人的导师告诉我们说：\BibleQuote{天主在教会内设立了：第一是宗徒，第二是先知，第三是教师；其次是行奇迹的，再次是治病的奇恩、救助人的、治理人的、说各种语言的恩赐。}\BibleRef{格前 12:28}\par

150. 你看，天主设立了宗徒，也设立了先知和教师，赐予了治病的奇恩——论到这些，你在前面已看到是由圣神赐予的；又赐予说各种语言的恩赐。但并非所有人都是宗徒，并非所有人都是先知，并非所有人都是教师。他说，并非所有人都有治病的奇恩，他说，并非所有人都说各种语言。\BibleRef{格前 12:30}因为全部神恩不可能集中在每一个人身上；各人按自己的能力，领受他所渴望或所配得的。但天主圣三的德能，慷慨赐予一切恩宠，却不像这种软弱。\par

151. 最后，天主设立了宗徒。那些天主在教会中设立的人，基督拣选并任命为宗徒，派遣他们进入世界，说：\BibleQuote{你们往普天下去，向一切受造物宣传福音。信而受洗的必要得救；但不信的必被判罪。信的人必有这些奇迹随着他们：因我的名驱逐魔鬼，说新语言，手拿毒蛇，甚或喝了什么致死的毒物，也决不受害；按手在病人身上，可使人痊愈。}你看，圣父和基督也在教会中设立了教师；正如圣父赐予治病的奇恩，圣子也同样赐予；正如圣父赐予说语言的恩赐，圣子也同样授予了。\par

152. 同样地，我们在上文也听到关于圣神的事，祂也赐予同样的恩宠。因为经上说：\BibleQuote{有人通过圣神领受了治病的奇恩，另一人领受了说各种语言的恩赐，另一人领受了说先知话的恩赐。}\BibleRef{格前 12:8}所以，圣神赐予与圣父同样的恩赐，圣子也同样赐予。现在让我们更清楚地认识我们在上文所提及的，即圣神托付与圣父和圣子相同的职务，并指派相同的人；因为保禄说：\BibleQuote{你们务要留意自己和整个羊群，圣神在其中委派你们作监督，牧养天主的教会。}\BibleRef{宗 20:28}\par

153. 因此，权柄是合一的，委派是合一的，赐予也是合一的。因为如果你将委派与权能分开，那么有什么理由说，那些基督立为宗徒的人，也是天主圣父所立，并为圣神所立呢？除非，也许像人一样，分享财产或权利，他们害怕法律上的侵害，于是将行动分割，权柄分散。\par

154. 这些事即使是世人之间也是狭隘而卑微的，他们大多数人，虽然行动不一致，但意志却能一致。因此有人问朋友是什么，回答说：\Quote{第二自己。}那么，如果有人将朋友定义为第二自己，也就是说，通过爱与善意的合一，那么，我们岂不更该珍视圣父、圣子、圣神之间的至尊合一吗？当藉由同一行动和神圣德能，所表达的不仅是合一，或更确切地说，是希腊文所谓的\OriginalTerm{同一性}{\GreekText{ταυτότης}}，因为\OriginalTerm{相同}{\GreekText{ταύτο}}意指\Quote{相同}；如此，圣父、圣子、圣神具有相同性。因此，拥有相同的意志和相同的德能，并非出于意志的情感，而是内在于天主圣三的本体内。\par

155. 这是宗徒信仰与虔诚的遗产，我们在\BibleBook{宗徒}{大事录}中也能看到。因此，保禄与巴尔纳伯服从了圣神的命令。所有宗徒都服从了，并立即任命了圣神吩咐要选拔出来的人：祂说：\BibleQuote{给我选拔出巴尔纳伯和扫禄。}\BibleRef{宗 13:2}你看见命令者的权柄了吗？细想服从者的功德。\par

156. 保禄信了，因为他信了，便弃绝了迫害者的狂热，获得了正义的冠冕。他曾经残害各教会，如今却信了；他归向信仰之后，便在圣神内宣讲圣神所吩咐的。\BibleRef{宗 9:20}圣神给祂的斗士傅了油，抖落了不信的尘土，将他作为不可战胜的征服者，置于各不敬虔者的集会前，并用许多苦难磨练他，为赢得在基督耶稣内崇高召叫的奖赏。\par

157. 巴尔纳伯也信了，因为信了而服从。因此，他被圣神的权柄所拣选，圣神丰沛地临于他，作为他功德的特殊标记，他当得起如此伟大的共融。因为同一恩宠照耀在这些同一圣神所拣选的人身上。\par

158. 保禄也不次于伯多禄，虽然伯多禄是教会的基石，而保禄是明智的建筑师，懂得如何使信从者的脚步稳固。我说，保禄并非不配与宗徒们共融，而是极易与首要者相比，不逊于任何人。因为不知道自己次一等的人，便使自己与别人同等了。\par

\SourceRecord{Translated by H. de Romestin, E. de Romestin and H.T.F. Duckworth.；Nicene and Post-Nicene Fathers, Second Series；Vol. 10.；Edited by Philip Schaff and Henry Wace.；Buffalo, NY: Christian Literature Publishing Co.,；1896.；<http://www.newadvent.org/fathers/34022.htm>.}

% structure: p0001 source=h1 target=section reason=page-title

\section{论圣神，第三卷}

% structure: p0002 source=h2 target=subsection reason=relative-heading-rank; base=h2

\subsection{第一章}

不仅先知和宗徒是由圣神所派遣，天主圣子亦然。这可由\Person{依撒意亚}和圣史们证明，并且解释了\Person{圣路加}为何记载同一圣神如同鸽子降下，停留在基督身上。然后，在确立了基督的这项派遣后，作者推断出圣子是由圣父和圣神所派遣，正如圣神是由圣父和圣子所派遣。\par

1. 我们在前一卷中，藉着圣经的明确证据已指出，宗徒和先知是由圣神所委派，后者是为说预言，前者是为宣讲福音，与由圣父及圣子委派的方式相同；现在我们加上一件众人都会合理地惊奇、且不容怀疑的事：圣神曾在基督身上；并且，正如基督派遣圣神，圣神也派遣了天主圣子。因为天主圣子说：\BibleQuote{上主的神临于我身，因为祂给我傅了油，派遣我向贫穷人传报喜讯，向俘虏宣告释放，向盲者宣告复明。} 祂从\BibleBook{依撒意亚}{书}中读了这段话后，在\BibleBook{}{福音}中说：\BibleQuote{今天这段圣经在你们耳中应验了；} \BibleRef{路 4:21} 祂是要指出这话是指祂自己说的。\par

2. 那么，既然基督说过：\BibleQuote{上主的神临于我身}，我们还会惊奇圣神派遣了先知和宗徒吗？祂说\BibleQuote{临于我身}说得恰当，因为祂是作为\OriginalTerm{人子}{Son of Man}在说话。因为作为人子，祂受了傅油并被派遣去宣讲福音。\par

3. 但如果他们不信圣子，且听圣父也说主的神临在基督身上。因为圣父对\Person{若望}说：\BibleQuote{你看见圣神降下，停在谁身上，谁就是那要以圣神付洗的。} \BibleRef{若 1:33} 天主圣父对若望说了这话，若望听见了，看见了，就信了。他听见了从天主来的话，在主身上看见了，就相信是圣神从天降下。因为降下的不是一只鸽子，而是如同鸽子般的圣神；因为经上这样记载：\BibleQuote{我看见圣神像鸽子一样从天降下。} \BibleRef{若 1:32}\par

4. 正如若望说他看见了，\Person{马尔谷}也这样记载；\Person{路加}还补充说圣神以形体的形状如同鸽子降下；你不要以为这是\OriginalTerm{降生成人}{incarnation}，而是一种显现。那么，那一位将显现带给他，为藉着这显现，使那看不见圣神的人得以相信，并藉着这显现显明祂与圣父及圣子共享同一尊威与权柄，同一奥迹中的同一工程，同一圣洗中的同一恩赐；免得我们以为主在其中受洗的那一位太软弱，以致仆人不能在祂内受洗。\par

5. 并且他说得恰当：\BibleQuote{停在祂身上} \BibleRef{若 1:33}，因为圣神随意感动先知发言或行事，却永远存留在基督内。\par

6. 再者，不要因他说\BibleQuote{在祂身上}而困惑，因为他是在论及人子，因为祂是作为人子受洗的。因为按\OriginalTerm{天主性}{Godhead}而言，圣神不是在基督之上，而是在基督之内；因为正如圣父在圣子内，圣子在圣父内，同样，天主的神和基督的神既在圣父内也在圣子内，因为祂就是祂口中的神。因为那出于天主的，居住在天主内，如经上所载：\BibleQuote{我们领受的，不是这世界的精神，而是出于天主的圣神。} \BibleRef{格前 2:12} 祂也居住在基督内，因为祂由基督领受；因为经上又写道：\BibleQuote{祂要由我领受：} \BibleRef{若 16:14} 在另一处：\BibleQuote{在基督耶稣内生命之神的律法，已使我脱离罪恶与死亡的律法。} \BibleRef{罗 8:2} 所以，按基督的天主性而言，祂并不在基督之上，因为天主圣三不在自己之上，而超越万有：圣三不在自己之上，而在自己之内。\par

7. 那么，既然天主圣子说：\BibleQuote{上主的神临于我身} \BibleRef{路 4:18}，谁还能怀疑圣神派遣了先知和宗徒呢？在另一处又说：\BibleQuote{我就是那一位，我是元始，我是终末。我的手奠定了大地，我的右手铺张了诸天；我一叫它们，它们就一同出现。你们都集合静听！他们中谁曾说过这事呢？上主爱了他，他必实行上主的意思，攻击\Place{巴比伦}和加色丁人后裔。我，我亲自说过，又召叫了他，我领他出来，使他的道路亨通。你们走近我前，听我说吧！我从起初就没有在暗处说过话，从事情发生之时我就在那里；现在，上主天主派遣了我与他的神。} 是谁说\BibleQuote{上主天主派遣了我与他的神}？除了那由圣父而来、为拯救罪人的那一位，还有谁呢？而且，正如你听见的，圣神派遣了祂，免得你听见圣子派遣圣神时，会以为圣神的能力较弱。\par

8. 因此，圣父和圣神都派遣了圣子；圣父派遣了祂，因为经上记载：\BibleQuote{但那\OriginalTerm{护慰者}{Paraclete}，就是父因我的名所要派遣来的圣神，} \BibleRef{若 14:26} 圣子派遣了祂，因为祂说过：\BibleQuote{当护慰者来到时，就是我从父那里给你们派遣的，那发于父的神，真理之神，} \BibleRef{若 15:26} 如果圣子和圣神彼此派遣，正如圣父派遣一样，那么，就没有隶属的次等，而是能力的共融。\par

% structure: p0012 source=h2 target=subsection reason=relative-heading-rank; base=h2

\subsection{第二章}

圣子和圣神同样被赐下；因此，这工程所显示的，不是隶属，而是同一的\OriginalTerm{天主性}{Godhead}。\par

9. 不是只有父派遣了子，也赐下了子，如同子自己舍下了自己一样。经上记载：\BibleQuote{愿恩宠与平安由天主我们的父和主耶稣基督赐与你们，祂为我们的罪舍弃了自己。}\BibleRef{迦 1:3}如果他们以为子因被派遣而从属于父，他们就不能否认祂被赐下是出于恩宠。但祂是由父赐下的，正如依撒意亚所说：\BibleQuote{有一个婴孩为我们诞生了，有一个儿子赐给了我们；}\BibleRef{依 9:6}但我敢说，祂也是由圣神赐下的，因为祂也被圣神所派遣。因为先知并未指明祂是由谁赐下的，这表明祂是由\OriginalTerm{天主圣三}{Trinity}的恩宠而赐下的；而且，既然子自己舍下了自己，就祂的\OriginalTerm{天主性}{Godhead}而言，祂不能从属于自己。因此，被赐下不能是天主性内从属的标记。\par

10. 但圣神也被赐下，因为经上记载：\BibleQuote{我要请求父，祂必会赐给你们另一位\OriginalTerm{护慰者}{Paraclete}。}\BibleRef{若 14:16}宗徒也说：\BibleQuote{为此，凡轻视这些事的，不是轻视人，而是轻视那将祂的圣神赐予我们的天主。}\BibleRef{得前 4:8}依撒意亚也指出圣神和圣子都是被赐下的：\BibleQuote{上主天主这样说：祂创造了并铺张了诸天，奠定了大地，及大地上的万物，祂将气息赐给地上的民众，把灵赐给在地上行走的人。}\BibleRef{依 42:5}对圣子则说：\BibleQuote{我是上主，我因仁义召叫了祢，我必提携祢的手臂，保护祢，立祢作人民的盟约，作万民的光明，开启盲人的眼目，从狱中领出被囚禁的人，从牢房里领出住在黑暗中的人。}\BibleRef{依 42:6}既然圣子既被派遣又被赐下，圣神也既被派遣又被赐下，那么，在行动上既有一体性，在\OriginalTerm{天主性}{Godhead}上也必然一致。\par

% structure: p0016 source=h2 target=subsection reason=relative-heading-rank; base=h2

\subsection{第三章}

同样的合一也可从圣神被称为\OriginalTerm{手指}{Finger}，圣子被称为\OriginalTerm{右手}{Right Hand}这一事实中识别出来；因为借着人类语言的运用，有助于理解神圣之事。律法的石板是由这手指写成的，后来它们被摔碎了，原因如下。最后，基督用同一手指书写；然而，我们不可因这肉身的比喻而承认圣神有任何次等性。\par

11. 同样，圣神也被称为\OriginalTerm{天主的手指}{Finger of God}，因为父、子及圣神之间有着不可分割、不可分离的共融。因为正如圣经称圣子为天主的\OriginalTerm{右手}{Right Hand}，如经上说：\BibleQuote{上主，祢的右手大显神能；上主，祢的右手击碎了仇敌；}\BibleRef{出 15:6}同样，圣神被称为天主的手指，如主亲自说的：\BibleQuote{但如果我是靠着天主的手指驱魔。}\BibleRef{路 11:20}因为在另一部福音书的同一处，祂称其为天主的神，正如你所见：\BibleQuote{但如果我是靠着天主的神驱魔。}\BibleRef{玛 12:28}\par

12. 那么，还有什么话能更明确地表明\OriginalTerm{天主性}{Godhead}及其作为的合一呢？这合一是按照父、子及圣神的天主性。我们应当理解，如果有人分裂\OriginalTerm{本体}{Substance}的合一，并将其能力增多，那么永恒天主性的圆满似乎比我们的身体更容易被分割，然而同一的天主性的永恒却是单一的。\par

13. 因为往往用我们自己的语言来估量那超越我们的事物是合适的，并且因为我们不能看见那些事，我们就从所能见的事物中推论出来。宗徒说：\BibleQuote{因为祂那看不见的美善，即祂永远的大能和祂的天主性，自从创世以来，都可凭祂所造的万物，辨认洞察出来。}\BibleRef{罗 1:20}他又加上：\BibleQuote{祂永远的大能和天主性。}\BibleRef{罗 1:20}其中，似乎一件事论及圣子，另一件事论及圣神；即如同圣子被称为父的永远大能，同样，圣神因为是天主，也应被相信是祂的永远天主性。因为圣子也因永远生活而是永恒生命。那么，天主的手指也是永恒的、天主性的。因为天主所有的，哪一样不是永恒和天主性的呢？\par

14. 正如我们所读到的，天主就是用这手指在梅瑟所接受的那些石版上写字。因为天主并非用血肉的手指写出我们所读的字母的形式和笔画，而是借着祂的圣神颁布了法律。所以宗徒说：\BibleQuote{法律原是属神的，确实不是用墨水写的，而是用生活天主的圣神写的；不是写在石版上，而是写在血肉的心版上。}\BibleRef{格后 3:3}因为如果宗徒的书信是在圣神内写成的，那么，我们有什么理由不相信天主的法律不是用墨水，而是用天主的圣神写成的呢？这圣神当然不会玷污，而是照亮我们心灵和思想的隐密处。\par

14. 法律原是写在石版上的，因为是作为一种预像而写的；但那些石版先被摔碎，从梅瑟手中抛出，因为犹太人背离了先知的事工。适当的是，石版被摔碎，而非文字被抹去。你要留心，不要让你的石版破碎，不要让你的心神和灵魂分裂。基督是被分裂的吗？祂没有被分裂，而是与父原为一体；谁也不能使你与祂分离。如果你的信德失落，你心中的石版就碎了。你灵魂的连贯性就减弱了，如果你不相信在\OriginalTerm{天主圣三}{Trinity}内的天主性的合一。你的信德被记录，你的罪也被记录，正如耶肋米亚所说：\BibleQuote{犹大的罪是用铁笔记录下来的，是用钻石尖刻在他们的心版和他们的祭坛角上的。}\BibleRef{耶 17:1}因此，哪里是恩宠，罪也在那里；但罪是用铁笔写下的，恩宠则是由圣神标明的。\par

15. 主耶稣也用这手指，低着头，神秘地在地上写字，那是在犹太人把一个淫妇带到祂面前的时候，以这幅图像表明，当我们审判别人的罪时，应当想起自己的罪。\par

16. 再者，为避免我们因天主藉祂的圣神颁布了法律，就认为圣神的职分有某种低微之处，或者从我们自身肉体的角度，以为圣神是天主的一小部分，宗徒在别处说，他不是用人的智慧言语，而是用圣神所教的言语，用属神的话解释属神的事；但属血气的人，不能领受天主圣神的事。\BibleRef{格前 2:13} 因为他知道，凡以属神的与属血肉的相比的人，便处于属血气的境界，不能列为属神的人；\BibleQuote{因为在他看来，这些都是愚妄，}他说，\BibleQuote{他不能领受。}\BibleRef{格前 2:13} 所以，他既知道这些问题会出现在属血气的人中，便预见未来而说：\BibleQuote{谁曾知道主的心意，而去教导祂呢？可是我们却有基督的心意。}\BibleRef{格前 2:16}\par

% structure: p0025 source=h2 target=subsection reason=relative-heading-rank; base=h2

\subsection{第四章}

对于那些因为圣神被称为\Quote{天主的手指}而认为祂次于圣父的人，\Person{圣盎博罗削}答复说，这同样也会贬低圣子，因为圣子被称为\Quote{天主的右手}。这些名称只应归因于天主性的合一，为此\Person{梅瑟}宣告，整个\OriginalTerm{天主圣三}{Trinity}在通过\Place{红海}时一同工作。事实上，圣神的工在那预像圣洗之处施展，并不足为奇，因为圣经教导说，三位同样地圣化并在此圣事中运作。\par

17. 但如果有人仍陷于属血肉的疑虑中，并因形体的形象而犹豫，他应当想想，凡对圣神抱有错误想法的人，对圣子也不会有正确的认识。因为若有人认为圣神是天主的一小部分，因为祂被称为\Quote{天主的手指}，那么这同一批人必定也会坚持说，圣子也只是天主的一小部分，因为祂被称为\Quote{天主的右手}。\par

18. 但圣子既被称为天主的右手，又被称为天主的德能；那么，如果我们考量自己的言语，没有德能就不可能有完美；因此，他们应当小心，不要以为那种说不出口的恶念：即圣父在自己的性体中只是半完美的，藉着圣子才得以圆满；他们也当停止否认圣子与圣父同为永恒。因为天主的德能何时不存在呢？但他们若认为天主的德能曾有一时不存，那他们就是说，天主圣父在某个时候缺乏完美，因他们认为那时德能有所欠缺。\par

19. 但如我所说，这些事被记载下来，是要我们将其归于\OriginalTerm{天主性的合一}{Unity of the Godhead}，并相信宗徒所说的话：\OriginalTerm{天主性的圆满}{fullness of the Godhead}有形有体地居住在\Person{基督}内，\BibleRef{哥 2:9} 这圆满也居住在圣父内，并居住在圣神内；而且，正如有天主性的合一，也同样有运作的合一。\par

20. 这也可从\OriginalTerm{梅瑟的歌}{Song of Moses}中得知，因为他带领犹太人民过海之后，承认了父、子、圣神的工，说：\BibleQuote{上主，你的右手大显神能；上主，你的右手粉碎了仇敌。}\BibleRef{出 15:6} 在此他宣认了圣子和圣父，因为圣子是圣父的右手。再往后，为了不忽略圣神，他补充说：\BibleQuote{你发出你的气，海涛涌起；水便聚起成堆，洪涛便凝立如垒。}\BibleRef{出 15:10} 由此表明了天主性的合一，而非天主圣三的不平等。\par

21. 由此可见，圣神也与圣父和圣子一同工作，以致波浪仿佛在海中凝结，水墙竖起让\Place{犹太人}通过，随后圣神使水回流，淹没了\Place{埃及人}。许多人还认为，出于同一根源，白天有\OriginalTerm{云柱}{pillar of cloud}引领犹太人，夜间有\OriginalTerm{火柱}{pillar of fire}，使圣神的恩宠保护祂的子民。\par

22. 而且，这件全世界都正确惊叹的天主的工程，并非没有圣神的工而成就，宗徒也声明说，其中预表了\OriginalTerm{属神的奥迹}{spiritual mystery}的真理，因为我们读到：\BibleQuote{我们的先祖都曾在云柱下，都从海中走过，都曾在云中和海中受了洗而归於梅瑟，都吃过同样的属神食粮，都饮过同样的属神饮料。}\par

23. 因为若没有圣神的运作，怎能有\OriginalTerm{圣事的预像}{type of a sacrament}呢？这圣事的全部真理都在圣神内。正如宗徒也阐明的，说：\BibleQuote{你们因着我们的主耶稣基督之名，并因我们天主的圣神，已经洗净了，已经圣化了，已经成了义人。}\BibleRef{格前 6:11}\par

24. 由此你们看到，圣父在圣子内工作，而圣子在圣神内工作。因此，不要怀疑，按圣经的次序，在预像中存在着真理本身所宣称的真实奥迹。因为谁能否认祂在\OriginalTerm{圣洗池}{Font}中的运作呢？我们在其中感受到祂的运作和恩宠。\par

25. 因为圣父如何圣化，圣子也如何圣化，圣神也如何圣化。圣父圣化，正如经上所载：\BibleQuote{愿赐平安的天主亲自完全圣化你们，将你们整个的神魂、灵魂和肉身，在我们的主耶稣基督来临时，保持得无瑕可指。}\BibleRef{得前 5:23} 而在别处圣子说：\BibleQuote{父啊，求你以真理圣化他们。}\BibleRef{若 17:17}\par

26. 但论到圣子，同一宗徒说：\BibleQuote{祂成了我们的智慧、正义、圣化者和救赎者。}\BibleRef{格前 1:30} 你们看到祂成了\OriginalTerm{圣化}{sanctification}者吗？祂是为了我们而成为这样的，并不是要改变祂原有的本质，而是要藉着血肉圣化我们。\par

27. 宗徒也教导说，圣神圣化人。因为他这样说：\BibleQuote{主所爱的弟兄们，我们该时常为你们感谢天主，因为天主从起初就拣选了你们，藉着圣神的圣化和信从真理而获得救恩。}\BibleRef{得后 2:13}\par

28. 那么，圣父圣化，圣子也圣化，圣神也圣化；但圣化是同一的，因为圣洗是同一的，圣事的恩宠也是同一的。\par

% structure: p0039 source=h2 target=subsection reason=relative-heading-rank; base=h2

\subsection{第五章}

作者总结了他已开始的论证，并从\Quote{手指}和\Quote{右手}等术语所意指的事实证实了统一性，即天主的工作与手的工作相同；而手的工作与手指的工作相同；最后，\Quote{手}这个词同样适用于圣子和圣神，而\Quote{手指}这个词适用于圣神和圣子。\par

29. 不过，那本身无需圣化、反而丰盈圣化的那位，若圣化每个人，又有什么惊奇的呢？因为，如我所说，我们已受教导，祂的尊威如此伟大，圣神似乎不可分离于天主圣父，就像手指不可分离于身体一样？\par

30. 但若有人认为这不应归于能力的同一性，而应归于能力的减损，那么他确实会陷入如此疯狂，以至于似乎把圣父、圣子和圣神塑造成一个身体形态，并为自己想象其肢体的某些区分。\par

31. 但是，让他们学习吧，正如我常说的，这证据所表明的不是不平等，而是能力的同一性；因为那些是天主的工程，也就是手的工程，而我们读到，这些也就是手指的工程。因为经上记载：\BibleQuote{高天陈述天主的光荣，穹苍宣扬祂手的化工；} 另一处又说：\BibleQuote{在起初你奠定了大地，主啊，诸天也是你手创造的工程。} 所以，手的工程与天主的工程是相同的。因此，并非按身体肢体的种类而有工程的区分，而是能力的同一性。\par

32. 然而，手所做的工程也是手指所做的工程，因为经上同样记载：\BibleQuote{当我仰望你手指创造的穹苍，和你所布置的月亮星辰，} 而又记载：\BibleQuote{主啊，你以你的工程令我欢乐，我要因你手中的工程而欢跃。} 那么，手指在这里被说成所做，岂比手所少呢？既然它们所做的与手所做相同，正如经上所说。\par

33. 然而，既然我们读到圣子是\Quote{手}（因为经上记载：\BibleQuote{这一切不都是我手所造的吗？}\BibleRef{依 66:2} 另一处：\BibleQuote{我要将你放在磐石缝中，用我的手遮掩你，我已将我的手放在磐石的遮掩之下，}\BibleRef{出 33:22} 这指的是\OriginalTerm{降生成人}{Incarnation}的奥迹，因为天主永恒的大能给自己披上了肉身的遮掩），显然圣经使用了\Quote{手}这个词既指圣子，也指圣神。\par

34. 再者，既然我们读到圣神是天主的手指，我们认为复数\Quote{手指}是指圣子和圣神。最后，为表明他既领受了圣子也领受了圣神的圣化，某位圣人说：\BibleQuote{你的双手创造了我，塑造了我。}\par

% structure: p0047 source=h2 target=subsection reason=relative-heading-rank; base=h2

\subsection{第六章}

圣神斥责，就像圣父和圣子斥责一样；事实上，没有祂，判官们也无法判断，正如撒罗满和达尼尔的判断所表明的，这将在后文略作解释；且启示达尼尔的，不是别人，正是圣神。\par

35. 当我们主张能力同一性时，为何要拒绝类似的说法呢？既然能力的同一性延伸至如此地步，圣神斥责，如同圣父斥责，圣子斥责。因为经上记载：\BibleQuote{主啊，求你不要在震怒中斥责我，也不要在气愤中惩戒我。} 然后，在第四十九篇[即第五十篇]圣咏中，上主这样说：\BibleQuote{我要斥责你，要把你的罪恶摆在你面前。} 同样地，圣子论及圣神说：\BibleQuote{当我离去时，我将派遣护慰者到你们这里来。祂来时，要指证世界关于罪恶、正义和审判。}\BibleRef{若 16:7}\par

36. 但是，不信者的疯狂要将我们引向何方，以至于我们似乎要证明圣神是否斥责，好像这尚属可疑，而事实上判官们若非藉着圣神，自己也不能判断。最后，撒罗满那著名的判断，在争执者所造成的困境中，一个因压死自己所生的孩子而想要夺去别人的孩子，另一个则在保护自己的儿子，他既发现了隐密思想中的欺诈，也发现了母亲心中的慈爱，这确实只是因着圣神的恩赐才如此令人赞叹。因为除了圣神的剑，没有什么剑能刺透那些女人的隐密情感；主论及这剑说：\BibleQuote{我来不是为带平安，而是带刀剑。}\BibleRef{玛 10:34} 因为内心深处不能被钢铁刺透，惟有藉着圣神：\BibleQuote{因为智慧之神是圣善的、唯一的、多样的、敏锐的、活泼的，} 稍后又提到：\BibleQuote{无所不察。}\BibleRef{智 7:22}\par

37. 考虑一下先知所说的，祂无所不察。因此，撒罗满也洞察一切，便命人取刀来，因为他一面假装要劈开婴儿，一面考虑到真正的母亲会更顾惜自己的儿子，胜于自己的安逸，会将慈爱置于权利之前，而非权利置于慈爱之前。而那伪装母亲情感的人，却因贪图达到自己目的的渴望而瞎眼，对她所没有柔情的人之毁灭毫不在意。于是，那位属神的人，为能判断一切（因为属神的人能审断一切），\BibleRef{格前 2:15} 在情感中寻找那隐藏在言语中的自然禀性，并追问慈爱，以彰显真理。因此，那位母亲因爱的情感而得胜，这爱是圣神的果实。\par

38. 祂在先知身上施行审判，因为智慧的言语是由圣神赐予的；\BibleRef{格前 12:8} 那么，人们怎能否认圣神能指证世界关于审判呢？祂从判断中除去疑惑，并赐予成功的结局。\par

39. \Person{达尼尔}也是一样，若非领受了\OriginalTerm{天主的圣神}{Spirit of God}，他决不可能揭穿那淫乱的奸情和欺诈的谎言。因为当\Person{苏撒纳}受到长老们的阴谋攻击，看到民众的心因顾及老人们而动摇，她孤立无助，在人中间孑然一身，意识到自己的贞洁，便祈求天主审判；经上记载：\BibleQuote{主听见了她的声音，当她将被带去受死的时候，主兴起了一个年轻青年的圣神，他的名字叫\Person{达尼尔}。} 于是，按照他所领受的圣神的恩宠，他揭穿了奸诈者自相矛盾的证词，因为这实在是天主大能的运作，使他的声音将那些内心隐藏的意念显明出来。\par

41. 那么，请领悟圣神这神圣而属天的奇迹。她宁愿自身保持贞洁，而不在乎民众的评价；她宁愿冒险牺牲自己清白的名声，而不愿牺牲她的端庄；当她被控告时，她缄默不语，当被定罪时，她保持静默，满足于自己良心的判断，甚至在危难中仍保守端庄，为使他们虽不能强迫她的贞洁，也不至于看似强迫她变得轻浮；当她呼求主的时候，她便获得了圣神，那圣神使长老们隐藏的心意显露出来。\par

42. 贞洁的人应学会不怕诽谤。因为那位宁要贞洁不要性命的人，并未失去生命，且保住了贞洁的光荣。同样，\Person{亚巴郎}曾奉命前往异域，既未因妻子端庄可能遭遇危险而却步，也未因摆在眼前的死亡恐惧而退缩，他保全了自己的生命和妻子的贞洁。因此，凡信赖天主的，从未有人后悔；在\Person{撒辣}身上，贞洁增进了虔诚，而虔诚也增进了贞洁。\par

43. 为免有人或许会想，正如经上说的\BibleQuote{天主兴起了一个年轻青年的圣神}，他里面的那圣神是人的精神，而非圣神，就让他往下读，便会发现\Person{达尼尔}领受了圣神，并且因此说了预言。最后，国王也擢升了他，因为他拥有圣神的恩宠。因为国王这样说：\BibleQuote{达尼尔阿，你能，因为\OriginalTerm{天主的圣神}{Holy Spirit of God}在你内。}\BibleRef{达 5:14} 再往后又记载：\BibleQuote{达尼尔被委任管理他们，因为在他内有一个\OriginalTerm{卓越的圣神}{excellent Spirit}。}\BibleRef{达 6:3} 而\OriginalTerm{\Person{梅瑟}的圣神}{Spirit of Moses}也分施给那些要担任民长的人。\BibleRef{户 11:25}\par

% structure: p0057 source=h2 target=subsection reason=relative-heading-rank; base=h2

\subsection{第七章。}

圣子自己不凭圣神审判或惩罚，因此，那同一圣神被称为\OriginalTerm{圣言的剑}{Sword of the Word}。但反过来，圣言也被称为\OriginalTerm{圣神的剑}{Sword of the Spirit}，由此可知在二者内有着最高的权能合一。\par

44. 但对于其他方面又该怎么说呢？我们听到\Person{主耶稣}不仅在圣神内审判，也加以惩罚。因为若非事先已判定其恶行应得的报应，他就不会惩罚那假基督，正如我们读到的：\BibleQuote{主耶稣要以\OriginalTerm{自己口中的圣神}{Spirit of His mouth}击杀他。}\BibleRef{得后 2:8} 然而在此，并不是领受的恩宠，而是合一保持不可分，因为基督不能没有圣神，圣神也不能没有基督。因为天主性体的合一不可分割。\par

45. 既然我们面前有这样的实例，主耶稣要以自己口中的圣神击杀，这圣神就被理解为好似\OriginalTerm{圣言的剑}{Sword of the Word}。最后，在福音中，主耶稣自己也说：\BibleQuote{我来不是为送平安，而是送刀剑。}\BibleRef{玛 10:34} 因为他来，是要赐下圣神；所以在他口中有一把双刃的剑，\BibleRef{默 19:15} 那其实就是圣神的恩宠。因此，圣神就是圣言的剑。\par

46. 为使你知道并没有不平等，而是性体的合一，圣言也是\OriginalTerm{圣神的剑}{Sword of the Holy Spirit}，因为经上记载：\BibleQuote{拿起信德作盾牌，使你们能以此扑灭恶者的一切火箭。并戴上救恩的盔，拿着圣神的剑，即\OriginalTerm{天主的圣言}{Word of God}。}\BibleRef{弗 6:16}\par

47. 那么，\OriginalTerm{圣言的剑}{Sword of the Word}既是圣神，\OriginalTerm{圣神的剑}{Sword of the Holy Spirit}既是\OriginalTerm{天主的圣言}{Word of God}，它们中间就确实有着权能的同一。\par

% structure: p0063 source=h2 target=subsection reason=relative-heading-rank; base=h2

\subsection{第八章。}

上述合一由此得到证明：正如圣经上说圣父被忧伤被试探，圣子也同样。圣子也在旷野中受了试探，在那里，铜蛇被竖立，作为十字架的预像；但宗徒说，圣神也在那里受了试探。\Person{圣盎博罗削}由此推论，以色列子民是由同一圣神引导进入预许的福地，而他的旨意和权能与圣父和圣子的旨意和权能是一体。\par

48. 我们也能在其他圣经章节中看到这种合一。因为\Person{厄则克耳}对犹太民众说：\BibleQuote{在这一切事上，你们使我忧伤，主说；}\BibleRef{则 16:43}\Person{保禄}在书信中对新子民说：\BibleQuote{不要使\OriginalTerm{天主的圣神}{Holy Spirit of God}忧伤，因为你们是在他内受了印记。}\BibleRef{弗 4:30} 同样，\Person{依撒意亚}论到犹太人自己说：\BibleQuote{但他们不信，反而使圣神忧伤；}\BibleRef{依 63:10}\Person{达味}论到天主说：\BibleQuote{他们在旷野中使至高者忧伤，在内心试探了天主。}\par

49. 也要明白，圣经在其他地方既说圣神被试探，又说天主被试探，也说基督被试探；因为你看见宗徒对格林多人说：\BibleQuote{我们也不要试探基督，如同他们中有些人试探了，就被蛇所灭。}\BibleRef{格前 10:9} 这惩罚是公义的：让那些不敬拜造物主的人遭受毒害。\par

50. 主确实命定了，藉着举起铜蛇，使那些被咬者的伤得到医治；因为铜蛇是十字架的预像；因为虽然\Person{基督}在肉身中被举起来，但宗徒也在祂内被钉在十字架上，向世界死，世界也向他死；因为他说：\BibleQuote{世界于我已被钉在十字架上了，我于世界也被钉在十字架上了。}\BibleRef{迦 6:14} 因此，世界在其诱惑中被钉在十字架上，所以挂起来的不是真蛇，而是铜蛇；因为主取了罪人的形状，的确有真实的肉身，却没有罪的真实，藉着人类软弱的欺骗外表模仿蛇，脱去肉身的皮，以摧毁那真蛇的狡诈。因此，在主为报复试探而前来援助人的十字架上，我接受了\OriginalTerm{天主圣三}{Trinity}的良药，在恶人中认出了反对\OriginalTerm{天主圣三}{Trinity}的冒犯。\par

51. 因此，当你在梅瑟的书中发现，主受试探而给犹太人百姓派遣了蛇时，你必须要么宣认圣父、圣子、圣神在神圣尊威中的合一，要么当宗徒的书信说圣神受试探时，无疑是以主之名指示圣神。但\BibleBook{致希伯来人}{书}的宗徒说圣神受试探，因为你可以读到：\BibleQuote{因此，圣神说：\InnerQuote{今天你们如果听从他的声音，不要再心硬了，像在叛乱之时，像在旷野中试探的那一天，在那里，你们的祖宗试探我，考验我，也看了我的事业，共四十年之久。所以我厌恶了那一世代，说：他们心中常迷惑，他们不认识我的道路；所以我在怒中起誓说：他们决不得进入我的安息。}}\BibleRef{希 3:7}\par

52. 因此，根据宗徒，圣神受试探。如果祂受试探，祂也必定在引导犹太人进入应许之地，正如经上记载：\BibleQuote{领他们由深渊经过，如马行于旷野而不颠仆，又如牛行于平原。圣神由上主那里降下，引导了他们。}\BibleRef{依 63:13} 祂确实给他们赐下了天上食粮的和缓之雨，祂以丰沛的甘霖使那并非出于土地、亦非农夫所种的每日收成果实累累。\par

53. 现在让我们逐一审视这些要点。天主曾应许犹太人安息；圣神称那安息为祂的。天主圣父述说自己被不信者试探，而圣神也说祂被同样的不信者试探，因为这试探是一个，是同一个不信者试探了天主圣三的同一\OriginalTerm{天主性}{Godhead}。天主断定犹太人不能进入流奶流蜜之地，即不能进入复活的安息；圣神也以同样的裁定断定他们：\BibleQuote{他们决不得进入我的安息。} 因此，这是一个意志的裁定，一个能力的卓绝。\par

% structure: p0071 source=h2 target=subsection reason=relative-heading-rank; base=h2

\subsection{第九章}

圣神被激怒，这事由圣\Person{伯多禄}的话语所证明，其中表明，天主的圣神就是主的圣神，既通过其他经文，也通过同一位宗徒对\Person{阿纳尼雅}和\Person{撒斐辣}的判决来证明；由此可论证，圣神与圣父及圣子的结合，以及祂自己的\OriginalTerm{天主性}{Godhead}，都得到了证明。\par

54. 然而，也许有人会说，如果同一位宗徒圣\Person{伯多禄}没有在另一处教导我们，圣神能因我们的罪受试探，那么这段经文就不能特别应用于圣神；因为你可以读到，\Person{阿纳尼雅}的妻子被这样问到：\BibleQuote{你们为什么共谋试探主的圣神呢？}\BibleRef{宗 5:9} 因为主的圣神就是天主的圣神；因为只有一个圣神，正如宗徒\Person{保禄}也教导说：\BibleQuote{你们已不属于肉性，而是属于圣神，只要天主的圣神住在你们内。谁若没有基督的圣神，就不属于基督。}\BibleRef{罗 8:9} 他先提及天主的圣神，随即又补充说，那同一位也就是基督的圣神。而在谈到圣神时，为了让我们明白，有圣神所在的地方就有基督，他又加上了：\BibleQuote{如果基督在你们内，}\BibleRef{罗 8:10}\par

55. 因此，正如此处我们理解，有圣神所在的地方就有基督；同样，在其他地方，他表明，有基督所在的地方也有圣神。因为他曾说过：\BibleQuote{你们既然寻找基督在我内说话的凭据，}\BibleRef{格后 13:3} 在另一处他又说：\BibleQuote{我想我也有天主的圣神。}\BibleRef{格前 7:40} 那么，这一结合是不可分离的，因为根据圣经的证词，凡圣父、圣子或圣神被指出的地方，就有天主圣三的圆满无缺。\par

56. 但\Person{伯多禄}本人，在我们所引述的事例中，首先提到了圣神，然后称祂为主的圣神，因为你可以读到如下：\BibleQuote{阿纳尼雅！为什么撒殚充满了你的心，使你欺骗圣神，扣留了田地的价钱呢？田地留着不是你的吗？卖了以后，价钱不是还由你支配吗？你为什么心中打算了这事？你不是欺骗人，而是欺骗天主！}\BibleRef{宗 5:3} 然后他下面对那妻子说：\BibleQuote{你们为什么共谋试探主的圣神呢？}\BibleRef{宗 5:9}\par

57. 首先，我们明白他称圣神为主的圣神。然后，由于他首先提及圣神，又加上：\BibleQuote{你不是欺骗人，而是欺骗天主，} 你必须要么理解圣神的\OriginalTerm{天主性}{Godhead}的同一性，因为当圣神受试探时，就是对天主说谎；要么，如果你试图排除天主性的同一性，根据圣经的话，你必定自己相信祂是天主。\par

如果我们理解这些表达既用于圣神也用于圣父，我们确实观察到在天主圣父和圣神内的真理与知识的合一，因为谬误同样被圣神和天主圣父所揭露。但如果我们已领受了关于圣神的每个真理，你，无信德的人，为什么试图否认你所读的呢？那么，你要么承认圣父、圣子及圣神的天主性的合一，要么承认圣神的天主性。无论你说哪一种，你都是在天主内说的，因为合一维系着天主性，天主性也维系着合一。\par

% structure: p0078 source=h2 target=subsection reason=relative-heading-rank; base=h2

\subsection{第十章。}

圣神的天主性由\Person{圣若望}的一段经文所支持。这节经文确实曾被异端者删除，但那是徒劳的尝试，因为他们的无信德因此更容易被定罪。我们考虑上下文的顺序，为了表明这节经文是指圣神的。他由圣神而生，就是由同一圣神而重生的人；基督自己也被相信是由圣神所生并重生的。再者，圣神的天主性可从\Person{圣若望}的两项见证推断出来；最后，解释圣神、水及血如何被称为见证者。\par

圣经不仅仅在这一处为\OriginalTerm{天主性}{\GreekText{θεότης}}，即圣神的天主性作证；而且主自己在福音中说：\Quote{圣神是天主。} 亚略派啊，你们如此明确地作证这节经文是指圣神说的，以至于你们从你们的抄本中删除了它，但愿这只是从你们的抄本中删除，而不是也从教会的抄本中删除！因为在\Person{奥森提乌斯}以不虔不信的武器和势力占领了\Place{米兰}教会的时候，\Place{色米姆}教会也遭到了\Person{瓦伦斯}和\Person{乌尔撒齐乌斯}的攻击，那时他们的主教们在信德上失足了；在教会的书籍中发现了你们的这个谎言和亵渎。你们可能在从前也做过同样的事。\par

你们确实能够涂抹掉文字，却不能除掉信德。那涂抹更暴露了你们，那涂抹更定了你们的罪；你们不能抹去真理，但那涂抹却将你们的名字从生命册上涂抹了。如果\Quote{因为天主是神}这节经文不是指圣神，为什么要删除呢？因为如果你们认为这话是指天主圣父说的，那么你们这些认为应该删除的人，就因此否认了天主圣父。任凭你们选择，在每一种情况下，你们自己不虔的罗网都会捆绑你们，如果你们因为否认圣父或圣神是天主而承认自己是异教徒。因此，你们虽涂抹了天主圣言，你们的宣认却依然留存，而你们却惧怕原文。\par

你们确实将经文涂抹在你们的心里和思想上了，但天主圣言没有被涂抹，圣神没有被涂抹，反而远离不虔的心灵；被涂抹的不是恩宠，而是邪恶；因为经上记载：\BibleQuote{是我，是我抹去你的过犯。}\BibleRef{依 43:25} 最后，\Person{梅瑟}为百姓祈求时说：\BibleQuote{求你从你所记录的册子上涂抹我，如果你不宽恕这百姓。}\BibleRef{出 32:32} 然而他并没有被涂抹，因为他没有邪恶，而是恩宠涌流。\par

这样，你们就被自己的宣认所定罪，你们不能说自己这样做是出于智慧，而是出于狡猾。因为你们狡猾地知道，你们被那节经文的证据所定罪，而你们的论证无法反驳那见证。否则，那段经文的意义能从何而出呢？因为整个经文的脉络都是关于圣神的。\par

\Person{尼苛德摩}询问重生的事，主回答说：\BibleQuote{我实实在在告诉你，人除非由水和圣神重生，不能进入天主的国。}\BibleRef{若 3:5} 为要显示有一种生是按照肉身的，另一种是按照圣神的，他接着说：\BibleQuote{由肉生的属于肉，因为是由肉生的；由圣神生的属于圣神，因为圣神是天主。} 他说：\BibleQuote{你不要惊奇，我给你们说：你们应该重生。风随意向哪里吹，你听到风的响声，却不知道风从哪里来，往哪里去：凡由圣神而生的就是这样。}\BibleRef{若 3:7} 跟随整段经文的脉络，你将发现天主以他陈述的圆满隔绝了你的不虔。\par

谁是由圣神而生，并成为神的人呢？岂不是那在心志上因圣神而更新的人吗？\BibleRef{弗 4:23} 这当然就是那藉着水和圣神重生的人，因为我们藉着重生的洗礼和圣神的更新而获得永生的盼望。\BibleRef{铎 3:5} 宗徒\Person{伯多禄}在另一处说：\BibleQuote{你们要领受圣神的洗礼。}\BibleRef{宗 11:16} 领受圣神的洗礼的人，岂不是那藉着水和圣神而重生的人吗？因此主论圣神说，我实实在在告诉你，人除非由水和圣神重生，不能进入天主的国。所以他宣称，我们在后者是在他内受生的，而他又曾说我们是在前者受生的。这是主的言论；我依据的是所记载的，不是凭论证。\par

然而我要问，如果我们毫无疑问是由圣神重生的，那么为何却要怀疑我们是由圣神而生的呢？因为主耶稣自己既是由圣神所生，又是由圣神重生的。如果你们承认他是由圣神所生的，因为你们不能否认这点，但却否认他重生了，那么承认唯独属于天主的事，却否认人所共有的，这是极大的愚昧。因此那对犹太人说的话，也可恰当地对你们说：\BibleQuote{若我给你们说地上的事，你们尚且不信， 若我给你们说天上的事，你们怎么会信呢？}\BibleRef{若 3:12}\par

66. 然而我们发现希腊原文中如此写道，祂说的不是\Emph{藉着}圣神，而是\Emph{由}圣神。因为经文是这样写的：\OriginalTerm{由水和圣神}{\GreekText{ἀ} \GreekText{μήν}, \GreekText{ἀμήν}, \GreekText{λέγω} \GreekText{σοι}, \GreekText{ἐὰν} \GreekText{μή} \GreekText{τις} \GreekText{γεννηθῇ} \GreekText{ἐξ} \GreekText{ὓδατος} \GreekText{και} \GreekText{Πνεύματος}}，也就是\Emph{由水和圣神}。因此，既然不可怀疑\BibleQuote{那由圣神所生的}指的是圣神；那么毫无疑问，圣神也是天主，正如经上所说：\BibleQuote{圣神是天主。}\par

67. 但这位圣史为使人们明白他所写的是关于圣神，又在别处说：\BibleQuote{耶稣基督藉着水和血而来，不单藉着水，而是藉着水和血。并且圣神作证，因为圣神是真理。原来作证的有三个：圣神、水及血；而这三者是一致的。}\par

68. 请听他们如何作证：圣神更新心灵，水用于洗濯，血指向赎价。因为圣神使我们成为子女，领养我们；圣洗的水洗涤了我们；主的血救赎了我们。这样，我们在属灵的圣事中获得了一个无形和一个有形的见证，因为\BibleQuote{圣神亲自与我们的心神一同作证。} \BibleRef{罗 8:16} 虽然圣事的圆满具于每一要素，但职分却有分别；既有职分的分别，当然没有同等的见证。\par

% structure: p0090 source=h2 target=subsection reason=relative-heading-rank; base=h2

\subsection{第十一章}

有人提出反对说，圣若望所说的\BibleQuote{圣神是天主}那句话，应指着天主圣父；因为基督后来宣称，当以心神和真理朝拜天主。答复如下：首先，\Emph{神}一词有时指属灵的恩宠；其次，要指出的是，如果他们坚持\Quote{以心神}一词是指圣神的位格，并因此否认应朝拜祂，那么这种论证同样不利于圣子；既然有无数经文证明圣子应受朝拜，我们由此推知，关于圣神也应立定同样的规则。为什么我们受命向祂的脚凳下拜？因为这脚凳象征主的身体，而圣神是这身体的创造者，因此祂应受朝拜，但这并不因此推论玛利亚应受朝拜。所以，对圣神的朝拜并未被废除，而当说\Quote{当以心神朝拜父}时，所表达的是祂与圣父的合一，这一点也有类似表述的支持。\par

69. 但也许有人会提及，在同一卷书的稍后部分，主再次说天主是神，但指的是天主圣父。因为福音中有这段经文：\BibleQuote{时候要到，且现在就是，那些真正朝拜的人，将以心神和真理朝拜父，因为父就是寻找这样朝拜祂的人。天主是神，朝拜祂的人，必须以心神和真理朝拜。} \BibleRef{若 4:23} 你试图用这段话不仅否认圣神的天主性，而且还从天主在心神内受朝拜，推论出圣神的低一等地位。\par

70. 对此我要简短地回答说，\Emph{神}常被用来指圣神的恩宠，正如宗徒也说：\BibleQuote{圣神亲自以无可言喻的叹息，为我们转求；} \BibleRef{罗 8:26} 也就是圣神的恩宠，除非你能够听到圣神亲口的叹息。因此，这里也是指天主受朝拜，不是出于心灵的邪恶，而是在圣神的恩宠中。\BibleQuote{因为智慧不进入存心不良的灵魂里；} \BibleRef{智 1:4} 因为 \BibleQuote{除非受圣神感动，也没有一个能说\Quote{耶稣是主}的。} \BibleRef{格前 12:3} 接着他立刻加上：\BibleQuote{神恩虽有区别。} \BibleRef{格前 12:4}\par

71. 但这不能涉及圣神的圆满或分割；因为人的心灵既无法领悟祂的圆满，祂也不会被分割成任何部分；祂乃是倾注[于灵魂]属神灵性恩宠的礼物，藉此，天主受朝拜，如同也在真理中受朝拜一样，因为除非有人以虔诚的情感汲取祂天主性的真理，没有人能朝拜祂。而且，他当然不是亲自把握基督，也不是亲自把握圣神。\par

72. 或者，如果你认为这是好像亲自论及基督和圣神，那么天主在真理中受朝拜，与祂在圣神中受朝拜的方式相同。因此，要么有同等的隶属关系——但愿你万万不要相信这种说法——那么圣子也不受朝拜；要么，正确的是，有同等合一的恩宠，而圣神也受朝拜。\par

73. 那么让我们在此得出结论，结束亚略异端派的不虔质疑。因为如果他们因天主在圣神内受朝拜，就说圣神不应受朝拜，那么让他们也因天主在真理中受朝拜，就说真理不应受朝拜。因为虽然有许多真理，正如经上所说：\BibleQuote{世人的忠诚信实已绝迹；} 但这些都是由神圣的真理——也就是基督——所赐的，祂说：\BibleQuote{我是道路、真理、生命。} \BibleRef{若 14:6} 因此，如果他们按常理理解这段经文中的真理，那么也让他们按常理理解圣神的恩宠，这样就没有绊脚石；或者，如果他们接受基督就是真理，那么让他们否认祂应受朝拜。\par

74. 但他们被虔敬者的行动和圣经的教导所驳斥。因为玛利亚曾朝拜基督，因此她被任命为向宗徒们报告复活的信使，\BibleRef{若 20:17}，解开了那世代相传的桎梏，以及女人那巨大的罪过。因为这正是主奥妙地所行的，\BibleQuote{好使罪恶在哪里越多，恩宠在哪里也越格外丰富。} \BibleRef{罗 5:20} 女人被委任为[向男子的信使]也是合宜的；好使那首先将罪的信息传给男人的，也首先将主恩宠的信息传报。\par

75. 宗徒们也朝拜了；因此，那些做信德见证的人，领受了关于信仰的权柄。天使也朝拜了，经上记载着：\BibleQuote{天主的众天使都要朝拜祂。} \BibleRef{希 1:6}\par

76. 但他们不仅崇拜祂的天主性，也崇拜祂的脚凳，正如经上所载：\BibleQuote{当崇拜祂的脚凳，因为脚凳是圣的。} 或者，如果他们否认在基督内，祂降生成人的奥迹也当受崇拜——在这些奥迹中，我们仿佛看到祂天主性的某些清晰痕迹，以及属天圣言的某些途径——那就让他们读读，就连宗徒们在祂带着肉身的荣耀复活时也朝拜了祂。\BibleRef{玛 28:17}\par

77. 因此，既然说天主在基督内受崇拜，这一点完全无损于基督，因为基督也受崇拜；那么，天主在圣神内受崇拜，也绝无损于圣神，因为圣神也受崇拜，正如宗徒所说：\BibleQuote{我们事奉天主的圣神，} \BibleRef{斐 3:3} 因为，事奉的也崇拜，正如之前经上说：\BibleQuote{你要朝拜上主你的天主，惟独事奉祂。} \BibleRef{申 6:13}\par

78. 但或许有人会曲解我们所举的例子，那么，让我们想一想，先知所说的\BibleQuote{朝拜祂的脚凳}这句，如何看来是指向天主降生成人的奥迹，因为我们不可从人间的习俗去评断脚凳。因为天主既没有形体，也不是可度量的，我们不应以为放下脚凳是为支撑祂的脚。而且经上记载，除了天主以外，没有什么当受崇拜，因为经上写着：\BibleQuote{你要朝拜上主你的天主，惟独事奉祂。} 那么，先知是在法律下成长并受教于法律的，怎会给出违背法律的诫命呢？因此，这个问题并非无关紧要，让我们更仔细地思考脚凳是什么。因为我们另处读到：\BibleQuote{天是我的宝座，地是我的脚凳。} \BibleRef{依 66:1} 但地并不当受我们崇拜，因为地是天主的受造物。\par

79. 但是，让我们看看，先知是否说的是：主耶稣摄取肉身时所取的那地当受崇拜。因此，脚凳被理解为地，而地又指基督的肉身，我们今日也在奥迹中朝拜这肉身，而宗徒们，如上所述，也在主耶稣内朝拜了这肉身；因为基督并非分裂的，而是一体的；当祂作为天主子受崇拜时，并不否定祂由童贞女所生。那么，既然降生成人的奥迹当受崇拜，而降生成人是圣神的工程，正如经上记载：\BibleQuote{圣神要临于你，至高者的能力要庇荫你，因此，那要诞生的圣者，将称为天主子，} \BibleRef{路 1:35} 无疑地，圣神也当受崇拜，因为按肉身由圣神所生的那位是受崇拜的。\par

80. 不要有人把这话扭曲到童贞玛利亚身上；玛利亚是天主的圣殿，而非圣殿的天主。因此，独有那在祂圣殿内行事的那位当受崇拜。\par

81. 那么，说天主在圣神内受崇拜，这一点并不相反我们的论证，因为圣神也受崇拜。不过，如果我们细审这些话本身，那么在圣父、圣子与圣神内，除了同一能力的合一性之外，我们还能理解出什么呢？因为所谓\BibleQuote{必须用心神和真理去崇拜}是什么意思呢？然而，如果你不把这话理解为圣神的恩宠，也不理解为良心的真诚信德，而是像我所说的那样，按\OriginalTerm{位格}{person}来理解（如果位格一词确实适合表达天主的尊威），你就必须把它理解为指基督和圣神。\par

82. 那么，说父在基督内受崇拜，这是什么意思？除非是说父在基督内，父在基督内说话，父居住在基督内。确实，这不是像一个物体在另一个物体内，因为天主并非物体；也不是像一种\OriginalTerm{混乱的混合}{confusus in confuso}，而是像真实的在真实的之内，天主在天主之内，光在光之内；如同永恒的父在共永恒的子之内。因此，所指的并非身体的结合，而是能力的合一。所以，凭着能力的合一，当天主父在基督内受崇拜时，基督也一同在父内受崇拜。同样地，凭着同一能力的合一，当天主在圣神内受崇拜时，圣神也一同在天主内受崇拜。\par

83. 让我们更勤奋地探究这一措辞与表达的力量，并从其它章节推导其本义。经上说：\BibleQuote{你以智慧造成了他们一切。} 我们在这里理解的是，智慧在被造的事物中毫无份量吗？但\BibleQuote{万有是借着祂而造成的。} \BibleRef{若 1:3} 达味也说：\BibleQuote{因上主之言，诸天造成。} 这样，那位称呼天主子为天上万物的创造者的他，也明确地说万有都是在子内造成的，好使在祂工程的更新中，绝不将子与父分离，反而将子结合于父。\par

84. 保禄也说：\BibleQuote{因为在天上和在地上的一切，可见的与不可见的，都是在祂内受造的。} \BibleRef{哥 1:16} 那么，当他说\Quote{在祂内}时，是否否认万有是借着祂而造成的呢？当然没有否认，而是加以肯定。而且他在另一处又说：\BibleQuote{只有一个主耶稣，万有都是借着祂而有。} \BibleRef{格前 8:6} 那么，当他说\Quote{借着祂}时，他是否否认万有是在祂内造成的呢？——他正是说万有都是借着祂有的。这些话，\Quote{在祂内}和\Quote{与祂同在}，具有这样的力量：它们被理解为在各方面都是一样的、相似的，而非相反的。他在后面也澄清了这一点，说：\BibleQuote{万有都是借着祂并在祂内受造的；} \BibleRef{哥 1:16} 因为，如我们上面所说，圣经证明这三个表述：\Quote{与祂同在}、\Quote{借着祂}和\Quote{在祂内}，在基督内是等同的。因为你读到，万有都是借着祂并在祂内造成的。\par

85. 也要明白，当万物受造时，圣父与祂同在，祂也与圣父同在。智慧说：\BibleQuote{当祂预备高天时，我已与祂同在；当祂划定水泉的界限时。}\BibleRef{箴 8:27} 在旧约中，圣父说 \BibleQuote{让我们造}，\BibleRef{创 1:26} 表明圣子乃万物创造者，当与祂自己同受钦崇。那么，既然万物被称为在圣子内受造，圣子就被视为创造者；同样，当天主被说成是在真理中受钦崇时，根据这词本身的正当含义，常以同样方式表达，应理解为圣子也受钦崇。同样，圣神也以同样方式受钦崇，因为天主在圣神内受钦崇。因此，圣父与圣子、圣神同受钦崇，因为天主圣三受钦崇。\par

% structure: p0109 source=h2 target=subsection reason=relative-heading-rank; base=h2

\subsection{第十二章}

圣\Person{保禄}已指出，三位宗徒在基督内所钦崇的\OriginalTerm{天主性}{Godhead}的光，就在天主圣三内；由此清楚显明，圣神也当受钦崇。从那些话本身可以看出，宗徒们所指的是圣神。这位圣神的\OriginalTerm{天主性}{Godhead}，由祂拥有一座圣殿这一点得到证明；祂居住其中，不是像司祭，而是像天主：并且与圣父和圣子同受钦崇；由此可知祂们性体的合一。\par

86. 然而，有人会否认永恒天主圣三的\OriginalTerm{天主性}{Godhead}应受钦崇吗？况且圣经也表述了天主圣三的不可言喻的尊威，正如宗徒在别处所说：\BibleQuote{因为那吩咐光从黑暗中照耀的天主，曾经照耀我们心中，为使我们以那在耶稣基督的面貌上，所闪耀的天主的光荣的知识，来光照我们。}\BibleRef{格后 4:6}\par

87. 当主耶稣在山上以其\OriginalTerm{天主性}{Godhead}的光闪耀时，宗徒们确实看见了这光荣：经上说：\BibleQuote{宗徒们看见了，就俯伏在地。}\BibleRef{玛 17:6} 难道你不认为，当他们俯伏时，他们是在钦崇吗？因为他们肉身的眼睛无法承受神圣光辉的明亮，永恒之光的荣耀使凡人的目力昏眩？那时，看见祂光荣的人还能说什么别的，除了 \Quote{请大家前来，一齐伏地朝拜，向祂俯伏}？因为 \BibleQuote{天主曾经照耀我们心中，为使我们以那在耶稣基督的面貌上，所闪耀的天主的光荣的知识，来光照我们。}\BibleRef{格后 4:6}\par

88. 那么，那照耀我们，使我们得以在耶稣基督的面貌上认识天主的是谁呢？因为他曾说 \BibleQuote{天主照耀}，为使天主的光荣在耶稣基督的面貌上被认识。我们还能想到谁，除了那显明出来的圣神？除了圣神之外，还有谁能被认为具有\OriginalTerm{天主性}{Godhead}的德能呢？因为那些排斥圣神的人，必须引进另一位，与圣父和圣子同受\OriginalTerm{天主性}{Godhead}的光荣。\par

89. 那么，让我们回到同样的话：\BibleQuote{那照耀我们心中，为使我们以那在耶稣基督的面貌上，所闪耀的天主的光荣的知识，来光照我们的，就是天主。} 这里明确指出了基督。因为那被说成是光照的光荣，若非圣神的光荣，又是谁的呢？所以，既然他论到天主的光荣，就是他提出了天主本身；如果是论圣父，那么 \BibleQuote{那吩咐光从黑暗中照耀的，并照耀在我们心中} 的这一位，必须被理解为圣神，因为我们不能钦崇圣父和圣子以外的另一位。所以，如果你明白了圣神，那么宗徒也称圣神为天主；因此，你现在也必须承认圣神的\OriginalTerm{天主性}{Godhead}，虽然你现在否认。\par

90. 但你怎能如此厚颜否认这一点，既然你曾读过圣神拥有一座圣殿。因为经上记载：\BibleQuote{你们是天主的圣殿，天主圣神住在你们内。}\BibleRef{格前 3:16} 然而，天主有圣殿，受造物没有真正的圣殿。但那住在我们内的圣神，拥有一座圣殿。因为经上记载：\BibleQuote{你们的身体是圣神的圣殿。}\BibleRef{格前 6:19}\par

91. 但祂住在圣殿里，不是像一位司祭，也不是像一位臣仆，而是像天主，因为主耶稣自己曾说：\BibleQuote{我要在他们中居住，在他们中往来；我要做他们的天主，他们要做我的百姓。}\BibleRef{肋 26:12} 达味也说：\Quote{主在自己的圣殿内。} 因此，圣神住在自己的圣殿里，正如圣父居住，正如圣子居住一样，圣子曾说：\BibleQuote{我父和我，我们要到他那里去，并在他那里作我们的住所。}\BibleRef{若 14:23}\par

92. 但圣父藉着所赐给我们的圣神，住在我们内。那么，不同性体怎能一同居住呢？这当然是不可能的。但圣神与圣父和圣子同住。因此，宗徒也将圣神的共融与耶稣基督的恩宠及天主的爱联系在一起，说：\BibleQuote{愿主耶稣基督的恩宠，和天主的爱，以及圣神的共融，与你们众人同在。}\BibleRef{格后 13:14}\par

91. 因此，我们观察到，圣父、圣子和圣神藉着同一性体的合一，居住在同一个〔主体〕内。所以，那住在圣殿内的拥有天主性的大能，因为正如我们是圣父和圣子的圣殿，我们也是圣神的圣殿；不是许多圣殿，而是一座圣殿，因为它是同一大能的圣殿。\par

% structure: p0119 source=h2 target=subsection reason=relative-heading-rank; base=h2

\subsection{第十三章}

对于那些反对者，他们说，当天主教徒将\OriginalTerm{天主性}{Godhead}归于圣神时，就是引进了三位神。回答是：依照同样的论证，他们自己便引进了两位神，除非他们否认圣子的\OriginalTerm{天主性}{Godhead}；随后，正统的道理被陈述出来。\par

92. 但你害怕什么呢？是否你惯于吹毛求疵的那事？怕你弄出三位神来。天主不许！因为何处\OriginalTerm{天主性}{Godhead}被理解为一，便只有一位天主被说出。因为即使我们称圣子为天主，我们也并不说有两位神。因为如果你在承认圣神的\OriginalTerm{天主性}{Godhead}时，认为被说成是三位神，那么，当你论及圣子的\OriginalTerm{天主性}{Godhead}时，由于你无法否认，你便引进了两位神。因为按你的意见，如果你认为天主是指一个位格的名称，而非一个性体的名称，那么你必然或者说有两位神，或者否认圣子是天主。\par

93. 然而，我们虽不原谅你们的过错，也要解除你们无知的罪名。因为按照我们的看法，由于只有一位天主，我们便领悟到一位\OriginalTerm{天主性}{Godhead}及能力的合一。因为正如我们说只有一位天主，既宣认圣父，又不以天主性的真名否认圣子；同样，我们也不把圣神排除在天主性的合一之外，而且不肯定、反否认有三位天主，因为造成多数的是能力的分裂，而不是合一。因为多属于数目，但\OriginalTerm{天主性体}{Divine Nature}不容许数目，天主性的合一又如何能承纳多呢？\par

% structure: p0123 source=h2 target=subsection reason=relative-heading-rank; base=h2

\subsection{第十四章}

除了以上所引的证据外，还可提出其他章节以证明三个\OriginalTerm{位格}{Persons}的至高权能。这里从致\Place{得撒洛尼}人书中引用两段，并对照\WorkTitle{圣经}的其他证词，以显示在这些章节中，圣神与其它位格一样被赋予了权能。然后，再引用致\Place{格林多}人后书中一段更为明确的经文，既可推断出圣神是主，又可推断出主在哪里，圣神也在哪里。\par

94. 所以，天主是唯一的，这并不损害永恒\OriginalTerm{天主圣三}{Trinity}的尊威，正如我们面前所引的实例所宣明的。我们不仅在这一处看到天主圣三在天主性之名中被表达出来；而且在许多地方，如我们上面所说，尤其在宗徒致\Place{得撒洛尼}人书信中，他极其清楚地阐明了圣父、圣子和圣神的天主性及主权。因为你读到如下的话：\BibleQuote{愿主使你们彼此相爱，并爱众人，日益增长且满溢，像我们待你们一样，好坚固你们的心，使你们在我们的父天主面前，在我们的主耶稣同祂的众圣者来临时，在圣德上无可指摘。}\BibleRef{得前 3:12}\par

95. 那么，使我们日益增长、满溢，在我们的主耶稣来临时，在父天主面前的主是谁呢？他提到了圣父，也提到了圣子；那么，除了圣神外，他还把谁与圣父和圣子并列呢？是谁在主坚固我们的心，使之成为圣的？因为圣德是圣神的恩宠，正如后面所说：\BibleQuote{借着圣神的圣德和对真理的信从。}\BibleRef{得后 2:13}\par

96. 那么，你以为这里被称为主的是谁，若非圣神？天主圣父岂没有教导过你吗？祂说：\BibleQuote{你看见圣神降下，停在谁身上，谁就是那在圣神内施洗的。}\BibleRef{若 1:33} 因为圣神以鸽子的形状降下，\BibleRef{路 3:22} 好能见证祂的智慧，成全那\OriginalTerm{圣洗圣事}{sacrament of the spiritual laver}，并显示祂的作为与圣父和圣子的作为是一体的。\par

97. 为了使你不至于以为宗徒出于疏忽而有所遗漏，而是他明知、有意并受圣神感动，将他所认为的天主称为主，他在致得撒洛尼人后书中再次重申说：\BibleQuote{愿主引导你们的心，使你们爱慕天主并耐心等待基督。}\BibleRef{得后 3:5}\par

98. 但我们不能否认，因为主论到祂说：\BibleQuote{我还有许多事要告诉你们，然而你们现在不能担负。当那一位真理之神来时，祂要把你们引入一切真理。}\BibleRef{若 16:12} \Person{达味}论到祂说：\Quote{愿祢的圣神引导我走上正路。}\par

99. 请看主的声音关于圣神说了什么。天主圣子来了，因为祂还没有倾注圣神，祂宣称我们像小孩子一样活着，没有圣神。祂说那要来的圣神，借着属灵年龄的增长，将使这些小孩子成为更坚强的人。祂这样说，并非要将圣神的能力置于首位，而是为显示，力量的圆满在于认识\OriginalTerm{天主圣三}{Trinity}。\par

100. 因此，你们必须或者提出除圣神以外的某个第四位，你们应当意识到祂；或者确定地，你们不要以为另有主，除非那已指明的圣神。\par

101. 但如果你要求\WorkTitle{圣经}明言圣神是主，你不可能没有注意到经上写着说：\BibleQuote{主就是那圣神。}\BibleRef{格后 3:17} 整段话的上下文表明这确是指圣神。那么，让我们思考宗徒的话：他说：\BibleQuote{几时诵读\Person{梅瑟}，他们的心上就蒙上了一层帕子；但当他们归向主时，帕子就会被除去。主就是那圣神；主的圣神在哪里，那里就有自由。}\BibleRef{格后 3:15}\par

102. 因此，他不仅称圣神为主，还加上说：\BibleQuote{因为主的圣神在哪里，那里就有自由。我们众人以揭开的脸，反映主的光荣，渐渐地光荣上加光荣，都变成了与主同样的肖像，正如由主、即圣神所变成的；}\BibleRef{格后 3:17} 也就是说，我们先前归向了主，借着属灵的理解，仿佛在\WorkTitle{圣经}的镜子中看见了主的光荣，如今正从使我们归向主的光荣，转变为天上的光荣。所以，既然我们归向的是主，而主就是那位圣神，我们借着祂得以更新且归向主，那么显然所指的就是圣神，因为那更新者接纳归向者。祂若不接纳，又怎能更新呢？\par

103. 然而，既然我们看见了合一的表达，为什么还要寻求字句的表达呢？因为即使你将主和圣神加以区分，你也不能否认主在哪里，圣神也在哪里，而归向主的人也就是归向了圣神。如果你纠缠字句，也不能损害合一；如果你要分离合一，你就得承认圣神自己是能力之主。\par

% structure: p0135 source=h2 target=subsection reason=relative-heading-rank; base=h2

\subsection{第十五章}

虽然圣神被称为主，但由此并不暗示有三位主；正如圣经许多章节中，圣子以与圣父相同的方式被称为主，这也不暗示有两位主；因为主性存在于天主性内，天主性亦存在于主性内，而这二者在三位格内不可分割地重合。\par

104. 但也许你会再次说：如果我称圣神为主，我就是在提出三位主。那么，当你称圣子为主时，你是要否认圣子，还是要承认两位主呢？绝对不可！因为圣子亲自说过：\BibleQuote{不要事奉两个主。} \BibleRef{玛 6:24} 但当然，他既不否认自己，也不否认圣父是主；因为他称圣父为主，正如你读到的：\BibleQuote{父啊！天地的主宰！我称谢你。} \BibleRef{玛 11:25} 并且主论及自己，正如我们在福音中读到的：\BibleQuote{你们称我为师傅和主，说得不错，我本来就是。} \BibleRef{若 13:13} 但他并没有说两位主；实际上，当他警告他们：\BibleQuote{不要事奉两个主。} 时，他也表明他没有说过两位主。因为主性只有一位的地方，就没有两位主，因为父在子内，子在父内，所以只有一位主。\par

105. 律法也是如此教导：\BibleQuote{以色列，你听着：主我们的天主是唯一的主。} \BibleRef{申 6:4} 也就是说，他是不可改变的，常处于能力的统一中，始终如一，不因任何增加或减少而改变。因此，\Person{梅瑟}称他为\Emph{一}，但同时也记载了主从主那里降下火来。\BibleRef{创 19:24} 宗徒也说：\BibleQuote{愿主赐他那天在主那里获得怜悯。} \BibleRef{弟后 1:18} 主从主那里降下；主从主那里赐予怜悯。当主从主那里降下时，他不被分割；当他从主那里赐予怜悯时，也没有分离，反而在每种情况中都表达了主性的唯一性。\par

106. 在圣咏中你也能看到：\BibleQuote{主对我主说。} 他并没有因为他称圣子为他的主，就否认圣父是他的主；但他称圣子为他的主，是为了让你不要认为他只是儿子，而是先知的主，正如主自己在福音中显示的那样，当他说：\BibleQuote{\Person{达味}因圣神称他为主，他怎么会是达味的儿子呢？} \BibleRef{玛 22:43} 是\Person{达味}，而非圣神，在圣神内称他为主。或者，如果他们由此错误地推论说，是圣神称他为主，那么，他们也就必须以同样的亵渎主张，天主之子也是圣神之子了。\par

107. 因此，正如当我们称圣父和圣子为主时，我们不说有两位主，同样，当我们承认圣神是主时，我们也不说有三位主。因为说有三个主或三个天主是亵渎的，同样，说有两个主或两个天主也是十足的亵渎；因为只有一位天主，一位主，一位圣神；那有天主性的就是主，那有主性的就是天主，因为天主性在主性内，而主性也在天主性内。\par

108. 最后，你已读到圣父既是主又是天主：\BibleQuote{主，我的天主，我要呼求你，求你应允我。} 你发现圣子既是主又是天主，正如你在福音中读到的，当\Person{多默}触摸了基督的肋旁时，他说：\BibleQuote{我主！我天主！} \BibleRef{若 20:28} 所以，正如圣父是天主，圣子是主，同样，圣子也是天主，圣父也是主。神圣的称谓从这一位转移到那一位，但天主性不变，尊位却永恒不变。因为他们并非[仿佛是]由恩惠汇集而成的贡献，而是自然之爱的自愿赠与；因为合一既有其特质，而特质也都在合一中联结。\par

% structure: p0142 source=h2 target=subsection reason=relative-heading-rank; base=h2

\subsection{第十六章。}

圣父是圣的，同样，圣子和圣神也是圣的，因此他们在同一个\OriginalTerm{三圣颂}{Trisagion}中受尊崇：我们无法用比称天主为圣更相称的方式来谈论天主；由此清楚可见，我们决不能贬低圣神的尊位。在他内具备一切属于天主的，因为在圣洗中，他与圣父和圣子一同被呼号，而且圣父已赐予他超越万有之上，任何人都不能剥夺他的这一地位。这样，从异端者用以攻击他尊位的\Person{圣若望}的那段经文本身，反而确立了三位的平等和天主性的合一。最后，在解释子如何从父领受之后，\Person{圣安博}显示了所引用的那段经文如何驳斥了各种异端。\par

109. 因此，圣父是圣的，圣子是圣的，圣神是圣的，但他们并不是三个圣者；因为只有一位圣的天主，一位主。因为真正的圣性是一，正如真正的天主性是一，正如那属于天主性体的真正圣性是一。\par

110. 所以，我们尊为圣的一切，都宣示那唯一的圣性。\OriginalTerm{革鲁宾}{Cherubim}和\OriginalTerm{色辣芬}{Seraphim}以不知疲倦的声音赞美天主说：\BibleQuote{圣，圣，圣，万军的主天主！} \BibleRef{依 6:3} 他们如此说，不是只说一次，免得你相信只有一位；不是两次，免得你排除圣神；他们不说\OriginalTerm{多个圣}{holies}，免得你以为有多位，而是重复三次，说同一个词，好使你在颂歌中也能明白三位格的区别，以及天主性的合一，并且当他们这样说时，他们是在宣扬天主。\par

111. 我们也发现，没有比称天主为圣更相称的方式，使我们能够宣扬天主。对天主而言，一切都太过低下，对主而言，一切都太过低下。因此，也从这一事实考虑，我们究竟是否应该贬低圣神，因为圣神的名是赞美天主。因为圣父如此受赞美，圣子也如此受赞美，圣神也以同样的方式被呼号和赞美。色辣芬发出赞美，所有蒙福者的团体也发出赞美，因为他们称呼天主是圣的，圣子是圣的，圣神是圣的。\par

112. 那么，祂既被司铎在圣洗中与圣父及圣子一同呼名，又在献祭中被呼求，在天上被\OriginalTerm{色辣芬}{Seraphim}与圣父及圣子一同宣认，与圣父及圣子一同居住于圣人之中，倾注于义人身上，并赐予先知们作为灵感之源，祂怎会不拥有属于天主的一切呢？为此，在圣经中，一切都称为\OriginalTerm{天主所默示的}{\GreekText{θεόπνευστος}}，因为天主默示了圣神所说的话。\par

113. 假若他们不愿意承认圣神拥有属于天主的一切，且能行万事，就请他们说说，圣神没有什么，不能行什么吧。因为正如圣子拥有一切，圣父也并不吝惜按照祂的本性将一切赐给圣子，已将那比一切都大的赐给了祂，正如圣经所证明的，说：\BibleQuote{我父所赐给我的，比一切都大。}\BibleRef{若 10:29} 同样，圣神也从基督领受了那比一切都大的，因为义德不知吝惜。\par

114. 所以，我们若用心细察，便也能在此领悟到天主能力的合一。祂说：\BibleQuote{我父所赐给我的，比一切都大，谁也不能从我父手里将他们夺去。我与父原是一体。}\BibleRef{若 10:29} 因为，倘若我们之前已正确地指明圣神是圣父的手，那么，这圣父的手无疑也就是圣子的手，因为同一个圣父的圣神也就是圣子的圣神。因此，我们中凡以天主圣三之名领受永生的人，既不会被夺离圣父，也不会被夺离圣子，同样也不会被夺离圣神。\par

115. 再者，正因圣经记载说：\BibleQuote{他要光荣我，因为他要把由我所领受的，传告给你们。}\BibleRef{若 16:14} （这话看似是指分施的职分，而非天主能力的神权，因为圣子所救赎的人，圣神也领受了他们，为圣化他们），我正是从他们用以构造诡辩的这些话语中，领悟出天主性的合一，而非恩赐的匮乏。\par

116. 圣父是通过生，而非通过收养赐予的；祂所赐的，仿佛是那蕴含在天主性本有的神权之内的事物，而不是仿佛因祂丰厚的恩惠而有所欠缺的事物。因此，因为圣子像圣父一样为自己赢取人灵，也像圣父一样赐予生命，祂便在大能的合一中表明了自己与圣父的平等，说：\Quote{我与父原是一体。} 因为当祂说 \Quote{我与父} 时，显示的是平等；当祂说 \Quote{原是一体} 时，宣告的是合一。平等排除混乱，合一排除分离。平等区分了圣父与圣子，合一则不分离圣父与圣子。\par

117. 因此，当祂说 \Quote{我与父} 时，就驳斥了撒伯里乌派，因为祂表明祂是一位，圣父是另一位；祂驳斥了福蒂尼安派，因为祂将自己与天主圣父结合。以头一句话，祂驳斥了那些人，因为祂说：\Quote{我与父}；以后一句话，祂驳斥了亚略派，因为祂说：\Quote{原是一体。} 然而，以这前后两句话，祂都驳斥了（1）撒伯里乌派的异端强暴，因为祂说：\Quote{我们是一体} \OriginalTerm{实体}{Substance}，而非 \Quote{我们是一个} \OriginalTerm{位格}{Person}。也驳斥了（2）亚略派的异端，因为祂说：\Quote{我与父}，而非 \Quote{父与我}。这当然不是粗鲁的标志，而是孝爱与预知的表现，以免我们从语序上产生错误的想法。因为合一不知次序，平等不知等级；也不能将粗鲁之名加于天主之子，祂本人乃是孝爱之师，岂会以粗鲁冒犯孝爱？\par

% structure: p0153 source=h2 target=subsection reason=relative-heading-rank; base=h2

\subsection{第十七章}

\Person{圣安博}通过实例说明，主发言时的地点对理解主的话语有所助益；他指出，基督在\Place{撒罗满廊}说出了引自\Person{圣若望}的那段话，这\Place{撒罗满廊}象征智者的心灵，因为他认为基督绝不会在愚昧或好争辩者的心中说出这话。他接着说，凡不信这些话的人，就用石头砸基督；正如\Person{伯多禄}因宣认这些话而蒙赐天国的钥匙，\Person{依斯加略}却因不信，而惨遭丧亡。他借此机会痛斥那买了天主圣子并出卖\Person{若瑟}的犹太人。他神秘地解释了各人所付的代价；同样地，在阐释了叛徒对\Person{玛达肋纳}香膏的抱怨之后，他补充说，异端者以某种方式购买基督，公教徒以另一种方式购买基督，而那些将圣神与圣父分离的人，是徒然自称为基督徒的。\par

118. 值得注意的是，主在何地进行了这番讨论，因为祂的言论常常要\Emph{更恰当地}通过祂交谈的地点类型来评估。当祂即将禁食时，（我们读到）祂被领入旷野，以使魔鬼的试探归于徒然。因为，虽然在富足中能节制度日固然值得赞许，但财帛与逸乐中的诱惑更为频繁。那时，试探者为了试探祂，应许祂丰饶；主却为了克胜，抱守饥饿。我并不否认在富足之中亦能节制；然而，正如航行海上的人虽常脱险，却比不肯出海的人更易遭遇危险。\par

119. 让我们思考其他几点。当耶稣即将应许天国时，祂上了一座山。另一次，祂带着门徒经过麦田，那时祂要在他们心中播下天上训诲的种子，好使灵魂的丰硕庄稼得以成熟。当祂即将完成祂所摄取的肉身的工程时，那时祂已在门徒身上看见成全——祂已将他们建立在祂圣言的基础之上——祂进入了一座园子，为能在上主的殿宇内栽植稚嫩的橄榄树，为能以祂圣血的溪流灌溉茂盛如棕榈的义人和多产的葡萄树。\par

120. 在这段经文中，我们也读到，祂在重建节那天，正在撒罗满廊下行走，也就是说，基督正在智慧和明达之人的胸怀中行走，好将他们的善情献给祂自己。什么是那廊，先知教导说：\Quote{我愿以心中的纯洁，行走于祢的殿中。} 因此，我们自身内就有天主的殿宇，有厅堂，有廊房，也有庭院，因为经上记载：\BibleQuote{让你的水在庭院中畅流。} \BibleRef{箴 5:16} 那么，敞开你心中的这廊房吧，向着圣言，祂对你说：\Quote{大大地张开你的口，我就要充满它。}\par

121. 所以，让我们听一听，在智慧与和平之人心中行走的圣言所说的话：\BibleQuote{我与父原是一体。} \BibleRef{若 10:30} 祂不会在不安和愚昧之人的胸怀中说出这话，因为\BibleQuote{属血气的人不能领受天主圣神的事，对他来说这是愚妄。} \BibleRef{格前 2:14} 罪人狭隘的胸怀容纳不了信德的伟大。最后，犹太人听了\BibleQuote{我与父原是一体，} \BibleRef{若 10:31} 就拿起石头来要砸祂。\par

122. 他若无法聆听这话，就是犹太人；他若无法聆听这话，便以自己背信弃义的石头，比任何岩石更粗糙的石头，砸向基督，若你相信我，他就是在伤害基督。因为虽然祂如今再也不能感受伤害：\BibleQuote{从今以后，我们不再按肉情认识基督；} \BibleRef{格后 5:16} 然而那位以教会的爱为乐的，却被亚略派的邪恶所石击。\par

123. \Quote{上主，你口中的法律对我是好的，我遵守你的诫命。} 你亲自说过，你与父原是一体。因为伯多禄相信了这话，他就领受了天国的钥匙，并毫无焦虑地赦免了罪过。犹达斯因为不信这话，便用自己邪恶的绳索吊死了自己。啊，不信之话的坚硬石头！啊，可耻的卖主者的绳索，以及犹太人更丑恶的买价！啊，可憎的金钱，义人竟被用来买死，或被卖！若瑟被卖了，耶稣基督被买了，一个为奴，一个受死。啊，可憎的遗产，啊，致命的卖价，它或者把自己的兄弟卖去受苦，或者给主定价为要毁灭祂——那位万民救恩的购买者。\par

124. 犹太人对两件最为首要的东西施以暴力，即信德与责任，而且在这两件事上，他们都对信德与责任的创立者基督施以暴力。因为在圣祖若瑟身上已有了基督的预像，而基督自己带着祂身体的真实性来临了，\BibleQuote{祂虽具有天主的形体，并没有以自己与天主同等，为应当把持不舍的，却使自己空虚，取了奴仆的形体，} \BibleRef{斐 2:6} 因着我们的堕落，也就是说，祂自愿取了奴仆之形，不避苦难。\par

125. 一处是以二十块银钱卖，另一处是以三十块银钱卖。因为祂的真实价值怎能被人领会呢？那价值是无法限量的。价格上出了错，是因为询问上有了错。在旧约中是卖二十块银钱，在新约中是卖三十块；因为真理比预像更有价值，恩宠比训练更慷慨，亲临比法律更优越，因为法律许下了来临，而来临则完成了法律。\par

126. 依市玛耳人用二十块银钱买，犹太人用三十块银钱买。这并非无关紧要的象征。不信的人为罪行更为慷慨，胜过信者为救恩。然而，值得思考每一次买卖的性质。二十块银钱是被卖为奴者的价格，三十块银钱是被交付于十字架者的价格。因为虽然\OriginalTerm{降生成人}{Incarnation}与\OriginalTerm{受难}{Passion}的奥迹必定同样令人惊异，可是信德的实现却在于受难的奥迹。我不曾对生于圣洁童贞女之事评价略低，但我更感谢地领受圣体的奥迹。还有什么比祂宽恕我对祂所做的不义更充满仁慈呢？然而，祂赐给我们如此大的恩赐，这更是更丰满的尺度：那本不会死因为祂是天主的，竟因我们的死而死，好使我们因祂的圣神而活。\par

127. 最后，犹达斯依斯加略将那香液估价为三百块银钱，并非没有意义，这似乎确实通过价格本身表达出了主的十字架。因此，主也说：\BibleQuote{她把这香液倒在我身上，是为我的安葬而做的。} \BibleRef{玛 26:12} 那么，为什么犹达斯对此估以如此高价呢？因为罪过的赦免对罪人更有价值，宽恕显得更为珍贵。最后，你会发现经上记载：\BibleQuote{那得赦免多的，也爱得更多。} \BibleRef{路 7:47} 因此，罪人自己也承认那已失去的主的受难之恩宠，并且为他们所迫害的基督作见证。\par

128. 或者因为，\BibleQuote{智慧不进入存心不良的灵魂，} \BibleRef{智 1:4} 叛徒邪恶的性情说出了这话，他以更高的价格估定主身体的苦难，是为了用这巨额价钱将所有人从信德引开。因此，主无偿地奉献了自己，好使贫困的需求无法阻止任何人归向基督。圣祖们以低价卖了他，好让所有人都能买。依撒意亚说：\BibleQuote{你们没有钱的，都来购买、来喝吧！不用花钱，也来吃吧！} \BibleRef{依 55:1} 为要赢得那没有钱的。啊，叛徒犹达斯，你将祂受难的香液估值为三百块银钱，却将祂的受难卖作三十块银钱。估价时慷慨，卖出时吝啬。\par

129. 所以，并非所有人都以同样的价格买卖基督；福提努斯，为置他于死地而买他，出的是某一种价格；亚略派，为侮辱他而买他，出的是另一种价格；公教徒，为光荣他而买他，出的又是另一种价格。但是，按照经上所记载的，没有钱的，可以不花代价来买他：\BibleQuote{没有钱的，可以不花代价来购买。} \BibleRef{依 55:1}\par

130. \BibleQuote{不是凡向我说\InnerQuote{主啊！主啊！}的人，就能进天国；} \BibleRef{玛 7:21} 虽然许多人自称基督徒，使用这名字，但并非所有人都将获得赏报。加音献了祭品，犹达斯领受了亲吻，但却对他说：\BibleQuote{犹达斯，你用口亲来负卖人子吗？} \BibleRef{路 22:48} 也就是说，你用爱情的标志充满你的邪恶，用和平的工具播种仇恨，用爱的外在标记施加死亡。\par

131. 所以，亚略派的人不要因使用这名字而自鸣得意，因为他们自称基督徒。主将回答他们说：你们高举我的名，却否认我的本体，我不承认我的名在哪里没有我永恒的天主性。那与圣父分离、与圣神隔绝的名，不是我的名；我不承认我的名在哪里不认识我的教义；我不承认我的名在哪里不认识我的圣神。因为他不知道，他将圣父的圣神与祂所创造的仆人相提并论。关于这一点，我们已经详细地讲过了。\par

% structure: p0169 source=h2 target=subsection reason=relative-heading-rank; base=h2

\subsection{第十八章}

在意图通过以上讨论的要点确立圣神的天主性时，圣盎博罗削再次触及其中一些要点；例如，祂不是犯罪而是赦罪；祂不是受造物而是造物主；最后，祂不是献上敬拜而是接受敬拜。\par

132. 但为总结，为在末尾更清楚地汇集散见于各处的论证，天主性的明显荣耀既由其他论证，也尤其由以下四点证明。天主通过这些标志被认识：或者祂没有罪；或者祂赦罪；或者祂不是受造物，而是造物主；或者祂不是给予，而是接受敬拜。\par

133. 因此，只有天主没有罪，因为除了天主一个外，没有谁是善的。\BibleRef{玛 19:17} 同样，只有天主才能赦罪，因为经上记载：\BibleQuote{除了天主一个外，谁能赦罪呢？} \BibleRef{路 5:21} 一个人若非受造物，就不能是万物的造物主；而非受造物者无疑就是天主；因为经上记载：\BibleQuote{去崇拜事奉受造物，以代替造物主；他是永远可赞美的。} \BibleRef{罗 1:25} 天主也不敬拜，而是受敬拜，因为经上记载：\BibleQuote{你要敬畏上主你的天主，只事奉他。} \BibleRef{申 6:13}\par

134. 因此，让我们考虑圣神是否具有任何这些标志，可以为祂的天主性作证。首先，让我们论述这一点：只有天主没有罪，并要求他们证明圣神有罪。\par

135. 但他们无法向我们说明这一点，却要求我们提供权威，即我们应通过经文证明圣神没有犯过罪，正如论及圣子时所说明的，祂没有犯过罪。\BibleRef{伯前 2:22} 让他们明白，我们是依据圣经的权威教导的；因为经上记载：\BibleQuote{在她内的神，原是聪明的，至圣的，唯一的，多样的，微妙的，敏捷的，明澈的，无玷的，} \BibleRef{智 7:22} 圣经说祂是无玷的；难道圣经关于圣子撒了谎，以致你们信它关于圣神也撒了谎吗？因为先知在同一地方论智慧说，任何玷污的东西都不能进入她内。她自己是无玷的，她的神也是无玷的。因此，如果圣神没有罪，祂就是天主。\par

136. 但祂自己既赦免罪过，又怎能算有罪呢？因此，祂没有犯过罪；如果祂没有罪，祂就不是受造物。因为凡受造物都难免有犯罪的可能，唯有永恒的天主性才无罪、无玷。\par

137. 现在让我们看看圣神是否赦罪。但关于这一点，不能有疑惑，因为主亲自说过：\BibleQuote{你们领受圣神罢！你们赦免谁的罪，就给谁赦免。} \BibleRef{若 20:22} 可见，罪是藉着圣神而得赦免的。但人使用他们的职务来赦罪，并非行使任何己力的权利。因为他们不是因自己的名，而是因圣父、及圣子、及圣神的名赦罪。他们祈求，天主性赐予；服务属于人，恩赐却来自至高的大能。\par

138. 并且不容置疑的是，罪是藉着圣洗而得赦免的；而在圣洗中，那是圣父、及圣子、及圣神的工程。因此，如果圣神赦罪，因为经上记载：\BibleQuote{除了天主一个外，谁能赦罪呢？} \BibleRef{谷 2:7} 那不能与本性之名的合一性分离的，也定然不能从天主的大能中被分离出去。那么，如果祂不从天主的大能中被分离，又怎能从天主的名中被分离呢？\par

139. 现在让我们看看祂是受造物，还是造物主。但由于我们在前面已经极清楚地证明了祂是造物主，正如经上记载：\BibleQuote{天主的神造了我，} \BibleRef{约 33:4} 并且已申明，大地貌容因圣神而更新，万物若无圣神则萎靡，显然圣神是造物主。但谁还能怀疑这一点呢？因为正如我们前面所显示的，甚至连主由童贞女降生成人——这超越一切受造物——也非没有圣神的工程。\par

因此，圣神不是受造物，而是创造者，而祂是创造者，就决不是受造物。既然祂不是受造物，无疑祂就是与父和子一起创造万物的创造者。如果祂是创造者，那么宗徒在谴责外邦人时说：\BibleQuote{他们去敬拜事奉受造物，以代替造物主；祂就是永远受赞美的天主，}\BibleRef{罗 1:25} 并且如我在前面说过，告诫人们应当事奉圣神，这就既表明祂是创造者，又因为祂是创造者而证明祂应被称为天主。他在\BibleBook{致希伯来人}{书}中也总结说：\BibleQuote{因为创造万有的乃是天主。}\BibleRef{希 3:4} 因此，他们要么指出有什么是在父、子和圣神之外受造的，要么就应当承认圣神也与父和子具有同一的天主性。\par

该作者也教导说，那位被他称为主和天主的神应当受崇拜。因为宇宙的天主和主当然应受一切崇拜，因为经上这样写道：\BibleQuote{你要朝拜上主你的天主，惟独事奉祂。}\BibleRef{申 6:13}\par

或者，让他们说说，他们在哪里读到过圣神崇拜。因为关于天主圣子经上说：\BibleQuote{天主的天使都要崇拜祂；}\BibleRef{希 1:6} 但我们没有读到过\Quote{圣神要崇拜祂}。祂既不在仆役和使者之中，而是与父和子一起，使义人事奉祂，经上记着说：\BibleQuote{我们事奉天主圣神。}\BibleRef{斐 3:3} 所以，宗徒教导我们必须事奉的那一位，我们应当崇拜祂；而我们事奉祂，也钦崇祂，正如经上所写的，再说一次同样的话：\BibleQuote{你要朝拜上主你的天主，惟独事奉祂。}\BibleRef{申 6:13}\par

虽然宗徒并没有忽略这一点，没有不教导我们应当崇拜圣神。因为我们已经证明圣神在先知内，无人能怀疑预言是由圣神所赐，显然，当那位在先知内者受崇拜时，就是圣神受崇拜。因此你看到：\BibleQuote{全教会如果聚在一起，众人都说异语，即使有未受教导的人或不信者进来，他岂不不说你们疯了吗？但如果众人都说先知话，即使有未受教导的人或不信者进来，他必被众人说服，被众人审断；他心中的隐情显露出来，于是他必俯伏在地，朝拜天主，宣告说：天主真是在你们中间。}\BibleRef{格前 14:23}\par

% structure: p0183 source=h2 target=subsection reason=relative-heading-rank; base=h2

\subsection{第19章}

圣\Person{盎博罗削}在上面证明了圣神居住并发言于先知内，由此推断出，祂知晓一切属于天主的事，因此与父和子是一体。他又根据圣神拥有一切天主所拥有的，即天主性、洞悉人心、真理、超越一切名号的名号，以及复活死者的能力——正如从\Person{厄则克耳}所证明的——在这点上圣神与圣子同等。\par

所以，正如父与子是一体，因为子拥有父所有的一切，同样，圣神也与父和子是一体，因为祂也知晓一切属于天主的事。因为祂不是凭借暴力获得这一切，免得造成好像有人遭受损失的伤害；祂也不是抢夺来的，以免那似乎被劫掠者遭受损失。因为祂既不是出于需要而夺取，也不是凭借更大能力的优越而强取，而是通过能力的合一而拥有。因此，如果祂行作这一切事，因为同一个圣神，实在做成这一切，\BibleRef{格前 12:11} 那么，祂拥有天主所有的一切，怎么能不是天主呢？\par

或者让我们考虑一下，天主有什么是圣神所没有的。天主父有天主性，子——整个天主性的圆满居住在祂内的——也有，而圣神也有，因为经上记载：\BibleQuote{天主的神在我鼻孔内。}\BibleRef{约 27:3}\par

再者，天主洞察人心和肺腑，因为经上记载：\BibleQuote{天主洞察人心和肺腑。} 子也有这权能，祂说：\BibleQuote{为什么你们心里思念恶事呢？}\BibleRef{玛 9:4} 因为耶稣知道他们的心意。圣神也有同样的权能，祂也向先知们启示别人心中的秘密，正如我们上面说过：\Quote{因为他心内的隐情显露出来。} 既然圣神甚至探察天主的深奥事理，我们查考人心的隐秘事又有什么可惊奇的呢？\par

天主的一个属性就是祂是真实的，因为经上记载：\BibleQuote{天主是真实的，而人都是虚谎的。}\BibleRef{罗 3:4} 圣神既是真理之神，\BibleRef{若 16:13} 祂难道会说谎吗？我们已经指出祂被称为真理，因为\Person{若望}也称祂为真理，如同称子那样。\Person{达味}在圣咏中说：\BibleQuote{求你发出你的光明和你的真理，引导我，带我到你的圣山和你的居所。} 如果你从这段经文想到子是光明，那么圣神就是真理；或者如果你认为子是真理，那么圣神就是光明。\par

天主有一个名号，超越一切名号，并把这个名号赐给了子，正如我们读到，因耶稣的名，众膝都要跪拜。我们且想一想，圣神是否有这圣名。可是经上记载：\BibleQuote{你们去使万民成为门徒，因父及子及圣神之名给他们授洗。}\BibleRef{玛 28:19} 那么，祂也有一个超越一切名号的名。所以，父和子所拥有的，圣神也借着祂本性圣名的合一而拥有。\par

149. 使死者复活是天主的特权。\BibleQuote{因为就如父使死者复活并赐他们生命，子也照样随自己的意愿赐人生命。}\BibleRef{若 5:21} 但圣神（天主藉以复活者）也同样使他们复活，因为经上记载：\BibleQuote{祂也要藉着住在你们内的圣神，使你们必死的身体复活。}\BibleRef{罗 8:11} 但为使你不认为这恩宠微不足道，要知道圣神也使人复活，因为先知厄则克耳说：\BibleQuote{来吧，圣神！向这些死者吹气，他们便可以活了。我遵命说预言，生命之神便进入了他们内，他们便活了，站起来，成为一支极庞大的军队。}\BibleRef{则 37:9} 随后天主又说：\BibleQuote{当我打开你们的坟墓，将我的百姓从坟墓中领出来，将我的圣神注入你们内，你们便活了；那时，你们便要承认我是上主。}\BibleRef{则 37:13}\par

150. 当祂谈论祂的圣神时，祂是否提及了圣神之外的任何其他存在？因为祂既不会将祂的圣神描述为因吹气而产生，这圣神也不会从世界的四极而来，因为我们所经验的风的吹拂是局部而非普世的；我们赖以为生的这个气也是个别而非普世的。但圣神的特性正是既超越万有又寓居万有之内。因此，从先知的话中我们可以明白，当圣神赋予生命时，那些早已散落的肢体架构（虽然已分散）如何重新组合，复原成活的身体；那些骨灰如何重新附着于所属的肢体，在重新形成活人外貌之前，先由一种聚合的倾向所激活。\par

151. 在这事所成就的肖像中，我们岂不认出了天主能力的合一？圣神使人复活的方式，与主在祂受难之时使人复活的方式相同：祂受难时，突然之间，一眨眼间，死者的坟墓开了，重新活过来的身体从坟中起来，死亡的气味消散，生命的气息恢复，那些死者的骨灰又呈现出活人的样子。\par

152. 因此，圣神拥有基督所有的一切，因而也拥有天主所有的一切，因为凡父所有的一切，子也都拥有，所以祂说：\BibleQuote{凡父所有的一切，都是我的。}\BibleRef{若 16:15}\par

% structure: p0194 source=h2 target=subsection reason=relative-heading-rank; base=h2

\subsection{第二十章}

从天主宝座流出的河流是圣神的预像，但达味所说的水则意指天上的众\OriginalTerm{异能天使}{Powers}。天主的国是圣神的工作；如果祂在这国中与子一同为王，这并不奇怪，因为\Person{圣保禄}宗徒应许我们也将与子一同为王。\par

153. 我们所读到的\Quote{有一条河流从天主的宝座流出}也不是一件小事。因为你读到圣史\Person{若望}的话大意如此：\BibleQuote{天使又指示给我一条生命之水的河流，光亮有如水晶，从天主和羔羊的宝座那里流出。流在城的街道中央；沿河两岸，有生命树，一年结十二次果子，每月结果一次，树的叶子可医治万民。}\BibleRef{默 22:1}\par

154. 这无疑是从天主宝座流出的河流，即圣神；凡信基督的人都要畅饮祂，正如祂亲自说的：\BibleQuote{谁若渴，到我这里来喝罢！凡信从我的，就如经上所说：从他的心中要流出活水的江河。}\BibleRef{若 7:37} 这是指圣神说的。因此，这河流就是圣神。\par

155. 因此，这河就是在天主的宝座中，因为水并不洗涤天主的宝座。再者，不论你如何理解那水，达味并没有说它在天主的宝座之上，而是在诸天之上，因为经上写着：\BibleQuote{请天上的诸水赞美上主的名！}\BibleRef{咏 148:4} 他说：\Quote{请\Emph{它们}赞美}，而不是\Quote{请\Emph{它}赞美}。因为若他本意是要我们理解水元素，必定会说\Quote{请它赞美}，但他用了复数，意在指\OriginalTerm{异能天使}{Powers}。\par

156. 圣神在天主的宝座中又有什么奇怪呢？因为天主的国本身就是圣神的工作，如经上所载：\BibleQuote{因为天主的国并不在于吃喝，而在于义德、平安以及在圣神内的喜乐。}\BibleRef{罗 14:17} 当救主自己说：\BibleQuote{凡一国自相纷争，必成废墟，}\BibleRef{玛 12:25} 随后又加上：\BibleQuote{如果我仗赖天主的神驱魔，那么天主的国已来到你们中间了，}\BibleRef{玛 12:27} 祂便表明天主的国是由祂自己和圣神共同而不分裂地掌管。\par

157. 然而，有什么比否认圣神与基督一同为王更愚蠢的呢？因为宗徒说，连我们也要在基督的国里与基督一同为王：\BibleQuote{若我们与祂同死，也必与祂同生；若我们坚忍到底，也必与祂一同为王。}\BibleRef{弟后 2:11} 不过，我们乃是藉着\OriginalTerm{义子}{adoptio}的名分，祂却是凭大能；我们因着恩宠，祂却是因着本性。\par

158. 因此，圣神与父及子共享王国，祂与祂们同性同体，同一主位，同一能力。\par

% structure: p0202 source=h2 target=subsection reason=relative-heading-rank; base=h2

\subsection{第二十一章}

\Person{依撒意亚}被圣神所派遣，他也因此看见了同一位圣神。那些转动的轮子和各样翅膀有何寓意，以及既然圣神被\OriginalTerm{色辣芬}{Seraphim}宣称为\OriginalTerm{万军}{Sabaoth}的上主，无疑只有不虔敬的人才会否认祂这个尊号。\par

159. 因此，既然圣神有份于王国，有什么能阻止我们理解派遣\Person{依撒意亚}的正是圣神呢？因为凭\Person{保禄}的权威我们不能怀疑，圣史\Person{路加}在\BibleBook{宗徒}{大事录}中极其认同他的判断，以致他这样写下\Person{保禄}的话：\BibleQuote{圣神藉\Person{依撒意亚}先知向你们的祖先说得真对：他说：\InnerQuote{你去对这人民说：你们听是听，但不了解；看是看，却不明白。}}\BibleRef{宗 28:25}\par

160. 那么，是圣神派遣了\Person{依撒意亚}。如果圣神派遣了他，那么毫无疑问，在\Person{乌齐雅}去世后，\Person{依撒意亚}看见的那位就是圣神，他当时说：\BibleQuote{我看见\OriginalTerm{万军的主}{Lord of Sabaoth}坐在崇高的宝座上，他的衣边拖曳满殿。\OriginalTerm{色辣芬}{Seraphim}侍立在他左右，各有六个翅膀：两个盖着脸，两个盖着脚，两个用于飞翔。他们互相高呼说：\InnerQuote{圣！圣！圣！\OriginalTerm{万军的主}{Lord of Sabaoth}！他的光荣充满大地！}}\BibleRef{依 6:1}\par

161. 如果\OriginalTerm{色辣芬}{Seraphim}站着，他们怎么在飞翔？如果他们在飞翔，他们怎么站着？如果我们不能理解这点，我们又如何想要理解我们从未见过的天主呢？\par

162. 然而，正如先知看见一个轮子转动在另一个轮子中间\BibleRef{则 1:16}（这当然不是指肉眼可见的任何景象，而是指各\OriginalTerm{约}{Testament}的恩宠；因为圣者的生命是光洁的，而且如此首尾一贯，后来的部分与先前的一致。）所以，轮子里的轮子就是法律下的生命、恩宠下的生命；因为犹太人在教会内，法律包容于恩宠中。因为那暗中是犹太人的，也是在教会内；内心的\OriginalTerm{割损}{circumcision}是教会内的一件圣事。但经上论到那包含在教会内的\OriginalTerm{犹大}{Jewry}说：\Quote{在犹大，天主为人所认识；}因此，就如轮子转动在轮子中间一样，翅膀也既是静止，又是在飞翔。\par

163. 同样，\OriginalTerm{色辣芬}{Seraphim}也用两个翅膀遮脸，两个翅膀遮脚，两个翅膀飞翔。因为这里也蕴藏着灵性智慧的奥迹。时期有静止，时期有飞逝；过去的静止，未来的飞逝，正如\OriginalTerm{色辣芬}{Seraphim}的翅膀，同样遮掩着面容或脚的天主；因为那位无始无终的天主，整个时间的流程与季节，从其起始与终结的知识来看，都是安息的。所以，过去与未来的时期是静止的，现在的则飞逝。不要追问祂的起始或终结的奥秘，因为它们根本不存在。你拥有现在，但你应当赞美，而非质问。\par

164. \OriginalTerm{色辣芬}{Seraphim}以不倦的声音赞美，而你却在质问吗？的确，当他们如此行事时，他们指示我们，有时我们不当质问天主，而应常赞美祂。因此，圣神也就是\OriginalTerm{万军的主}{Lord of Sabaoth}。除非那位基督所拣选的导师不中悦那些不虔敬的人，或者他们能够否认圣神是\OriginalTerm{万军的主}{Lord of powers}，祂按自己的意愿赐予各种能力。\par

% structure: p0210 source=h2 target=subsection reason=relative-heading-rank; base=h2

\subsection{第二十二章}

为证明\OriginalTerm{天主圣三的合一}{Unity in Trinity}，我们考量上面引用的\Person{依撒意亚}的经文，并且指出，在解释此段时，无论是归于圣父、圣子或圣神，其意义并无分别。如果那位被钉在十字架上的乃是\OriginalTerm{光荣之主}{Lord of glory}，那么，圣神同样在一切事上与圣父和圣子平等，而\OriginalTerm{亚略派}{Arians}永远无法贬损祂的光荣。\par

165. 现在可以认清圣父、圣子及圣神在尊威与主权上的合一性。因为许多人说，那时\Person{依撒意亚}所看见的是天主圣父。\Person{保禄}说那是圣神，\Person{路加}也证实这一点。\Person{若望}却将它归于圣子。因为论及圣子时，他这样写道：\BibleQuote{耶稣说了这些话，就离开他们，隐藏起来了。虽然他在他们面前行了这么多的神迹，他们仍不信从他，这正应验了\Person{依撒意亚}先知所说的话：\InnerQuote{主，谁会相信我们的报道呢？主的手臂又向谁显示了出来呢？}\BibleRef{依 53:1}因此，他们不能信，因为\Person{依撒意亚}又说过：\InnerQuote{他使他们眼瞎，使他们心硬，免得他们眼睛看见，心里明白而悔改，我医治他们。}\BibleRef{依 6:10}\Person{依撒意亚}说这话时，是因为看到了他的光荣，而谈论他。}\BibleRef{若 12:36}\par

166. \Person{若望}说，\Person{依撒意亚}说了这些话，并且极清楚地揭示了，是圣子的光荣向他显现。然而，\Person{保禄}却叙述，是圣神说了这些话。那么，这种差异来自何处？\par

167. 的确，这是言词上的差异，而非意义的不同。因为他们虽说了不同的话，但两者都没有错，因为既可在圣子内看见圣父，如圣子所说：\BibleQuote{谁看见了我，就是看见了父，}\BibleRef{若 14:9}也能在圣神内看见圣子；因为正如\BibleQuote{除非受圣神感动，没有一个能说\InnerQuote{耶稣是主}的，}\BibleRef{格前 12:3}同样，看见基督也不是凭借肉身的眼目，而是凭借圣神的恩宠。因此经上也说：\BibleQuote{你这睡眠的，醒起来吧！从死者中起来吧！基督必要光照你！}\BibleRef{弗 5:14}而\Person{保禄}当时已看不见了，他如何看见基督，岂不是在圣神内吗？\BibleRef{宗 9:8}因此，主说：\BibleQuote{我显现给你，正是为了选派你作仆役和见证人，为你所见到的我，以及我将显现给你的事。}\BibleRef{宗 26:16}因为先知们也是领受了圣神而看见基督的。\par

168. 因此，神视是一个，命令之权是一个，光荣是一个。我们岂能否认圣神也是\OriginalTerm{光荣之主}{Lord of glory}，既然那由\Person{童贞玛利亚}因圣神而降生的\OriginalTerm{光荣之主}{Lord of glory}被钉在十字架上了呢？因为基督并非两个中的一位，而是一位，在万世之前，作为圣子由圣父所生；而在世界中，则因取肉身而作为人子诞生。\par

169. 我为何还要说，圣神一如圣父与圣子，无玷而全能呢？因为\Person{撒罗满}在\WorkTitle{智慧篇}中称祂为\OriginalTerm{全能且洞察一切的}{\GreekText{παντοδύναμον}, \GreekText{πᾶνέπίσχοπον}}，意即祂全能并监察万物，\BibleRef{智 7:22}正如我们在前面已经指明的。因此，圣神享有尊荣与荣耀。\par

170. 现在请你想一想，或许有什么不合乎祂的尊位；或者，\OriginalTerm{亚略}{Arian}啊，如果这使你不悦，你就把祂从与父及子的共融中拉下来吧！但如果你想要把祂拉下来，你会看见诸天在你头上翻转，因为诸天的一切力量都来自圣神。如果你想要把祂拉下来，你必须先对天主下手，因为圣神是天主。但你怎么能拉下那洞察天主的深处的那一位呢？\par

\SourceRecord{Translated by H. de Romestin, E. de Romestin and H.T.F. Duckworth.；Nicene and Post-Nicene Fathers, Second Series；Vol. 10.；Edited by Philip Schaff and Henry Wace.；Buffalo, NY: Christian Literature Publishing Co.,；1896.；<http://www.newadvent.org/fathers/34023.htm>.}
